% La correspondencia y los intercambios escritos (continuación) : cartas formales y administrativas — Espagnol 1ère A — Cours signé OnBuch+
% Programme officiel MINESEC. Compiler avec : tectonic EA20-correspondencia-intercambios-escritos-continuacion.tex
\newcommand{\DOCMATIERE}{Espagnol}
\newcommand{\DOCNIVEAU}{1ère A}
\newcommand{\DOCMODULE}{Módulo 5 : Les mass médias et la communication}
\newcommand{\DOCLECON}{EA20}
\newcommand{\DOCTITRE}{La correspondencia y los intercambios escritos (continuación) : cartas formales y administrativas}
% OnBuch+ — préambule « cours signés » ENRICHI (compilé avec tectonic / XeLaTeX)
% Système visuel « Pop » d'OnBuch+ : fond crème (jamais blanc), bordures encre
% épaisses, ombres « dures » décalées, titres Archivo Black en capitales.
% Version MATHÉMATIQUES : graphes (pgfplots/tikz), boîtes pédagogiques
% (définition, propriété, méthode, attention, le savais-tu, exercice résolu,
% exercice, TP, bilan), page de garde et signature OnBuch+ ancrée en bas.
%
% Macros attendues dans main.tex AVANT \input{preamble} :
%   \DOCMATIERE  \DOCNIVEAU  \DOCMODULE  \DOCLECON (numéro)  \DOCTITRE
\documentclass[11pt]{article}

\usepackage{fontspec}
\usepackage[spanish,french]{babel} % français = langue principale ; l'espagnol via \esp{..} / espagnol
\usepackage[a4paper,margin=2cm,top=3.4cm,bottom=2.8cm,headheight=1.4cm,footskip=1.3cm]{geometry}
\usepackage{amsmath,amssymb,amsfonts}
\usepackage{mathtools} % \xmapsto, \coloneqq... souvent utilisés par les modèles
\usepackage{cancel} % simplifications barrées (bilans de forces)
% \abs{z} (module, valeur absolue) et \norm{u} (norme), utilisés par les modèles
\providecommand{\abs}{}\let\abs\relax\DeclarePairedDelimiter{\abs}{\lvert}{\rvert}
\providecommand{\norm}{}\let\norm\relax\DeclarePairedDelimiter{\norm}{\lVert}{\rVert}
\usepackage[table]{xcolor} % [table] : fournit aussi \rowcolors (alternance de lignes)
\usepackage{pagecolor}
\usepackage{tikz}
\usetikzlibrary{arrows.meta,positioning,calc,shapes.geometric,shapes.misc,shapes.arrows,decorations.pathmorphing,decorations.pathreplacing,decorations.markings,patterns,shadows,mindmap,trees,fit,backgrounds,matrix,intersections,angles,quotes}
\usepackage{pgfplots}
\usepackage[version=4]{mhchem} % équations nucléaires \ce{^{235}_{92}U}
\usepackage[european,siunitx]{circuitikz} % schémas électriques
\pgfplotsset{compat=1.18}
\usepgfplotslibrary{fillbetween,groupplots} % aires entre courbes (calcul intégral)
\usepackage{enumitem}
\usepackage{fancyhdr}
% [most] charge toutes les bibliothèques tcolorbox (listings, minted, raster,
% documentation...) et gonfle inutilement la mémoire TeX sur un document long
% (~300 boîtes) : on ne charge que ce qui est réellement utilisé.
\usepackage[skins,breakable]{tcolorbox}
\usepackage{graphicx}
\usepackage{colortbl}
\usepackage{booktabs}
\usepackage{tabularx}
\usepackage{diagbox} % case d'en-tête diagonale des tableaux à double entrée
% Colonnes extensibles centrée / gauche / droite (C, L, R), souvent utilisées par les modèles
\newcolumntype{C}{>{\centering\arraybackslash}X}
\newcolumntype{L}{>{\raggedright\arraybackslash}X}
\newcolumntype{R}{>{\raggedleft\arraybackslash}X}
\usepackage{array}
\usepackage{multicol}
\usepackage{siunitx}
% parse-numbers=false : accepte aussi \SI{2,5 \times 10^{-3}}{...}
\sisetup{locale=FR,output-decimal-marker={,},group-minimum-digits=4,parse-numbers=false}
\DeclareSIUnit\ohm{\ensuremath{\symup{\Omega}}} % U+2126 absent de la police
\DeclareSIUnit\division{div} % réglages d'oscilloscope (ms/div, V/div)
\DeclareSIUnit\tr{tr} % tours (vitesse de rotation, ex. \SI{1500}{\tr\per\minute})
\DeclareSIUnit\molar{M} % concentration molaire (ex. \SI{3}{\molar}, \SI{1}{\micro\molar})
\DeclareSIUnit\an{an} % année (le modèle confond parfois avec \year, un compteur TeX) — ex. \SI{5}{\kilo\meter\per\an}
\DeclareSIUnit\FCFA{FCFA} % franc CFA (ex. \SI{500}{\FCFA\per\kilo\gram})
\DeclareSIUnit\degreeBrix{\textdegree Brix} % teneur en sucre d'un sirop/jus (réfractométrie)
\DeclareSIUnit\month{mois} % le modèle utilise parfois \month, un compteur TeX, comme unité de durée
\usepackage{qrcode}
% \degree : commande fréquemment utilisée par les modèles pour un angle ou une
% température en dehors de \SI{}{} (non fournie par les paquets chargés).
\providecommand{\degree}{^{\circ}}
% \qed (fin de démonstration), utilisé par les modèles sans amsthm
\providecommand{\qed}{\hfill$\square$}
% \barre : double barre des valeurs interdites dans un tableau de variations (tkz-tab)
\providecommand{\barre}{\ensuremath{\|}}
% environnement proof (amsthm non chargé) utilisé par les modèles
\ifdefined\proof\else\newenvironment{proof}[1][Démonstration]{\par\noindent\textit{#1.}\ }{\qed\par}\fi
\usepackage{lastpage}
\usepackage{hyperref}
\hypersetup{hidelinks}

% ============================================================
% Palette « Pop » OnBuch+ — mêmes hex que la classe P (plus_ui.dart)
% ============================================================
\definecolor{poporange}{HTML}{F59321}
\definecolor{popdark}{HTML}{E07A0C}
\definecolor{popcream}{HTML}{FFF6E8}
\definecolor{popink}{HTML}{1C1714}
\definecolor{poppurple}{HTML}{7A5AE0}
\definecolor{popgreen}{HTML}{2E9E6A}
\definecolor{poppink}{HTML}{E0457B}
\definecolor{popblue}{HTML}{2F7DE1}
\definecolor{popgold}{HTML}{C98A12}
\definecolor{popmuted}{HTML}{8A7F76}
\definecolor{popline}{HTML}{EADFD2}
\definecolor{popgray}{HTML}{8A7F76}
\colorlet{popgrey}{popgray}
% Alias tolérants (le modèle nomme parfois les couleurs différemment)
\colorlet{popwhite}{white}
\colorlet{popblack}{popink}
\colorlet{popred}{poppink}
\colorlet{popyellow}{popgold}
% Teintes claires (fonds de tableaux/encadrés) — le modèle nomme parfois
% les couleurs avec un suffixe L (ex. popblueL) sans les avoir définies.
\colorlet{poporangeL}{poporange!20}
\colorlet{popdarkL}{popdark!20}
\colorlet{poppurpleL}{poppurple!20}
\colorlet{popgreenL}{popgreen!20}
\colorlet{poppinkL}{poppink!20}
\colorlet{popblueL}{popblue!20}
\colorlet{popgoldL}{popgold!20}
\colorlet{popredL}{popred!20}
\colorlet{popgrayL}{popgray!20}
% Teintes claires pour fonds de tableaux / zones
\colorlet{popblueL}{popblue!10!white}
\colorlet{poporangeL}{poporange!14!white}
\colorlet{popgreenL}{popgreen!12!white}
\colorlet{poppurpleL}{poppurple!12!white}
\colorlet{poppinkL}{poppink!12!white}
\colorlet{popgoldL}{popgold!14!white}

\pagecolor{popcream}

% ============================================================
% Typographie
% ============================================================
\defaultfontfeatures{Ligatures=TeX}
\setmainfont{PlusJakartaSans-Medium.ttf}[Path=../../../fonts/,BoldFont=PlusJakartaSans-Bold.ttf]
% Maths en sans-serif (Fira Math) avec chiffres et lettres droites de Plus
% Jakarta, pour que les formules se fondent dans le texte.
\usepackage[math-style=ISO,bold-style=ISO]{unicode-math}
\setmathfont{FiraMath-Regular.otf}[Path=../../../fonts/]
\setmathfont{PlusJakartaSans-Medium.ttf}[Path=../../../fonts/,range={up/num,up/{Latin,latin},\mathup}]
\usepackage{icomma} % virgule décimale sans espace en mode math
% Caractères mathématiques Unicode absents des polices de texte (Plus Jakarta,
% Archivo Black) : rendus par leur équivalent mathématique, en texte comme en
% formule — sinon ils s'affichent en carré vide (ex. « Arithmétique dans ℤ »).
\usepackage{newunicodechar}
\newunicodechar{ℝ}{\ensuremath{\mathbb{R}}}
\newunicodechar{ℤ}{\ensuremath{\mathbb{Z}}}
\newunicodechar{ℂ}{\ensuremath{\mathbb{C}}}
\newunicodechar{ℕ}{\ensuremath{\mathbb{N}}}
\newunicodechar{ℚ}{\ensuremath{\mathbb{Q}}}
\newunicodechar{𝔻}{\ensuremath{\mathbb{D}}}
\newunicodechar{α}{\ensuremath{\alpha}}
\newunicodechar{β}{\ensuremath{\beta}}
\newunicodechar{γ}{\ensuremath{\gamma}}
\newunicodechar{δ}{\ensuremath{\delta}}
\newunicodechar{ε}{\ensuremath{\varepsilon}}
\newunicodechar{θ}{\ensuremath{\theta}}
\newunicodechar{λ}{\ensuremath{\lambda}}
\newunicodechar{μ}{\ensuremath{\mu}}
\newunicodechar{σ}{\ensuremath{\sigma}}
\newunicodechar{τ}{\ensuremath{\tau}}
\newunicodechar{φ}{\ensuremath{\varphi}}
\newunicodechar{ψ}{\ensuremath{\psi}}
\newunicodechar{ω}{\ensuremath{\omega}}
\newunicodechar{Σ}{\ensuremath{\Sigma}}
% majuscules produites par \MakeUppercase dans les titres (α-aminés -> Α)
\newunicodechar{Α}{\ensuremath{\alpha}}
\newunicodechar{Β}{\ensuremath{\beta}}
\newunicodechar{Γ}{\ensuremath{\Gamma}}
\newunicodechar{Δ}{\ensuremath{\Delta}}
\newunicodechar{Ε}{\ensuremath{\varepsilon}}
\newunicodechar{Θ}{\ensuremath{\Theta}}
\newunicodechar{Λ}{\ensuremath{\Lambda}}
\newunicodechar{Μ}{\ensuremath{\mu}}
\newunicodechar{Τ}{\ensuremath{\tau}}
\newunicodechar{Φ}{\ensuremath{\Phi}}
\newunicodechar{Ψ}{\ensuremath{\Psi}}
\newunicodechar{Ω}{\ensuremath{\Omega}}
\newunicodechar{↦}{\ensuremath{\mapsto}}
\newunicodechar{∈}{\ensuremath{\in}}
\newunicodechar{∉}{\ensuremath{\notin}}
\newunicodechar{∧}{\ensuremath{\wedge}}
\newunicodechar{⇔}{\ensuremath{\Leftrightarrow}}
\newunicodechar{⟺}{\ensuremath{\Longleftrightarrow}}
\newunicodechar{⇒}{\ensuremath{\Rightarrow}}
\newunicodechar{⟹}{\ensuremath{\Longrightarrow}}
\newunicodechar{∘}{\ensuremath{\circ}}
\newunicodechar{∪}{\ensuremath{\cup}}
\newunicodechar{∩}{\ensuremath{\cap}}
\newunicodechar{≡}{\ensuremath{\equiv}}
\newunicodechar{⊂}{\ensuremath{\subset}}
\newunicodechar{⊥}{\ensuremath{\perp}}
\newunicodechar{∃}{\ensuremath{\exists}}
\newunicodechar{∀}{\ensuremath{\forall}}
\newunicodechar{∛}{\ensuremath{\sqrt[3]{\,}}}
\newunicodechar{∜}{\ensuremath{\sqrt[4]{\,}}}
\newunicodechar{⃗}{\ensuremath{{}^{\to}}}
\newunicodechar{ⁿ}{\ifmmode^{n}\else\textsuperscript{\ensuremath{n}}\fi}
\newunicodechar{ˣ}{\ifmmode^{x}\else\textsuperscript{\ensuremath{x}}\fi}
\newunicodechar{ᵇ}{\ifmmode^{b}\else\textsuperscript{\ensuremath{b}}\fi}
\newunicodechar{ʳ}{\ifmmode^{r}\else\textsuperscript{\ensuremath{r}}\fi}
\newunicodechar{ᵘ}{\ifmmode^{u}\else\textsuperscript{\ensuremath{u}}\fi}
\newunicodechar{ʸ}{\ifmmode^{y}\else\textsuperscript{\ensuremath{y}}\fi}
\newunicodechar{ᵃ}{\ifmmode^{a}\else\textsuperscript{\ensuremath{a}}\fi}
\newunicodechar{ᵉ}{\ifmmode^{e}\else\textsuperscript{\ensuremath{e}}\fi}
\newunicodechar{ᵏ}{\ifmmode^{k}\else\textsuperscript{\ensuremath{k}}\fi}
\newunicodechar{ᵖ}{\ifmmode^{p}\else\textsuperscript{\ensuremath{p}}\fi}
\newunicodechar{ᴺ}{\ifmmode^{N}\else\textsuperscript{\ensuremath{N}}\fi}
\newunicodechar{ᶜ}{\ifmmode^{c}\else\textsuperscript{\ensuremath{c}}\fi}
\newunicodechar{ʰ}{\ifmmode^{h}\else\textsuperscript{\ensuremath{h}}\fi}
\newunicodechar{ᵠ}{\ifmmode^{q}\else\textsuperscript{\ensuremath{q}}\fi}
\newunicodechar{ₙ}{\ifmmode_{n}\else\textsubscript{\ensuremath{n}}\fi}
\newunicodechar{ₐ}{\ifmmode_{a}\else\textsubscript{\ensuremath{a}}\fi}
\newunicodechar{ᵢ}{\ifmmode_{i}\else\textsubscript{\ensuremath{i}}\fi}
\newunicodechar{ᵣ}{\ifmmode_{r}\else\textsubscript{\ensuremath{r}}\fi}
\newunicodechar{ₖ}{\ifmmode_{k}\else\textsubscript{\ensuremath{k}}\fi}
\newunicodechar{ᵦ}{\ifmmode_{\beta}\else\textsubscript{\ensuremath{\beta}}\fi}
\newunicodechar{ₓ}{\ifmmode_{x}\else\textsubscript{\ensuremath{x}}\fi}
\newfontfamily\popdisplay{ArchivoBlack-Regular.ttf}[Path=../../../fonts/]
\newfontfamily\poplogo{SpaceGrotesk-Bold.ttf}[Path=../../../fonts/]
\newfontfamily\popsemi{PlusJakartaSans-SemiBold.ttf}[Path=../../../fonts/]
\newfontfamily\popxbold{PlusJakartaSans-ExtraBold.ttf}[Path=../../../fonts/]
\newfontfamily\popxboldls{PlusJakartaSans-ExtraBold.ttf}[Path=../../../fonts/,LetterSpace=3.0]

\setlist[enumerate]{leftmargin=1.7em,itemsep=5pt,topsep=4pt}
\setlist[itemize]{leftmargin=1.7em,itemsep=4pt,topsep=4pt,label={\color{poporange}\textbullet}}
\setlength{\parindent}{0pt}
\setlength{\parskip}{5pt}
\renewcommand{\arraystretch}{1.25}

% pgfplots : style Pop par défaut
\pgfplotsset{
  every axis/.append style={axis line style={popink,line width=1pt},tick style={popink},
    grid style={popline,line width=0.5pt},label style={font=\small},tick label style={font=\footnotesize}},
  popcurve/.style={poporange,line width=1.8pt,no markers,smooth},
}
% Même style disponible dans un \draw TikZ simple (hors pgfplots)
\tikzset{popcurve/.style={poporange,line width=1.8pt}}
\tikzset{popgrid/.style={popline,very thin}}
\colorlet{popcurve}{poporange} % le nom du style est parfois utilisé comme couleur

% ============================================================
% Pill / squiggle
% ============================================================
\newcommand{\poppill}[3]{%
  \tikz[baseline=-0.55ex]\node[draw=popink,line width=1.2pt,rounded corners=2.6pt,%
    fill=#2,inner xsep=6.5pt,inner ysep=3pt]{%
    \popxboldls\fontsize{7}{7}\selectfont\color{#3}\MakeUppercase{#1}};%
}
\newcommand{\popsquiggle}[1]{%
  \tikz[baseline=-0.4ex]{
    \draw[#1,line width=2.6pt,line cap=round]
      plot [smooth] coordinates {(0,0.09) (0.34,0.01) (0.68,0.21) (1.02,0.01) (1.36,0.10)};
  }%
}
% Mot-clé surligné (marqueur orange sous le texte)
\newcommand{\cle}[1]{{\popxbold\color{popdark}#1}}

% ============================================================
% En-tête / pied
% ============================================================
\pagestyle{fancy}
\fancyhf{}
\renewcommand{\headrulewidth}{0pt}
\renewcommand{\footrulewidth}{0pt}
\lhead{\raisebox{-0.15ex}{\poplogo\fontsize{13}{13}\selectfont\color{popink}OnBuch\color{poporange}+}}
\rhead{\poppill{\DOCMATIERE\ \textbullet\ \DOCNIVEAU\ \textbullet\ Leçon \DOCLECON}{white}{popink}}
\lfoot{\popxbold\fontsize{7.6}{7.6}\selectfont\color{popmuted}\MakeUppercase{Cours signé OnBuch+}\ \textbullet\ {\popsemi\DOCTITRE}}
\rfoot{\tikz[baseline=-0.6ex]\node[rounded corners=8.5pt,fill=popink,minimum height=17pt,minimum width=17pt,inner xsep=5pt,inner sep=0]{\popxbold\fontsize{8}{8}\selectfont\color{popcream}\thepage\,/\,\pageref*{LastPage}};}

% ============================================================
% Titres
% ============================================================
\newcommand{\chaptitre}[2]{%
  \begin{tcolorbox}[enhanced,colback=poporange,colframe=popink,boxrule=2.6pt,arc=11pt,
      drop shadow=popink,left=15pt,right=15pt,top=12pt,bottom=11pt,before skip=3pt,after skip=18pt]
    \poppill{#2}{popink}{popcream}\par\vspace{9pt}
    {\raggedright\hyphenpenalty=10000\exhyphenpenalty=10000\popdisplay\fontsize{22}{24}\selectfont\color{popcream}\MakeUppercase{#1}\par}
  \end{tcolorbox}
}
% Section numérotée : pastille orange avec le numéro + titre Archivo
\newcounter{popsec}
\newcommand{\coursec}[1]{%
  \stepcounter{popsec}\setcounter{popsub}{0}%
  \par\vspace{16pt}\noindent
  \begin{minipage}{\linewidth}
  \tikz[baseline=-0.7ex]\node[draw=popink,line width=1.6pt,fill=poporange,rounded corners=4pt,minimum size=22pt,inner sep=0]{\popdisplay\fontsize{12}{12}\selectfont\color{popcream}\thepopsec};%
  \hspace{8pt}{\popdisplay\fontsize{15}{17}\selectfont\color{popink}\MakeUppercase{#1}}\par
  \vspace{4pt}\noindent{\color{poporange}\rule{36pt}{3.6pt}}
  \end{minipage}\par\nopagebreak\vspace{8pt}
}
\newcounter{popsub}
\newcommand{\courssub}[1]{%
  \stepcounter{popsub}%
  \par\vspace{9pt}\noindent{\popxbold\fontsize{11.5}{13}\selectfont\color{popink}{\color{poporange}\thepopsec.\thepopsub}\hspace{6pt}#1}\par\nopagebreak\vspace{4pt}
}

% ============================================================
% Boîtes pédagogiques — même recette : fond blanc, bordure épaisse colorée,
% ombre dure encre, badge plein. Titre optionnel [#1].
% ============================================================
\tcbset{popbase/.style={breakable,enhanced,colback=white,boxrule=2.3pt,arc=9pt,
  shadow={3pt}{-3pt}{0pt}{popink},left=14pt,right=14pt,top=11pt,bottom=11pt,before skip=11pt,after skip=15pt,
  fontupper=\popsemi\fontsize{10.6}{14.5}\selectfont\color{popink},
  fontlower=\popsemi\fontsize{10.6}{14.5}\selectfont\color{popink}}}
% Coche dessinée (le glyphe ✓ n'existe pas dans les polices du document)
\newcommand{\popcheck}[1]{\tikz[baseline=-0.5ex]\draw[#1,line width=1.6pt,line cap=round,line join=round](-0.13,0.0)--(-0.04,-0.09)--(0.13,0.1);}
\providecommand{\coche}{\popcheck{popgreen}}
\providecommand{\resultat}[1]{\par\smallskip\noindent\fcolorbox{popgreen}{popgreenL}{\parbox{\dimexpr\linewidth-2\fboxsep-2\fboxrule\relax}{\textbf{Résultat.} #1}}\par\smallskip}
\newcommand{\popboxhead}[3]{\poppill{#1}{#2}{white}\ifx\relax#3\relax\else\hspace{6pt}{\popxbold\fontsize{10.6}{12}\selectfont\color{popink}#3}\fi\par\vspace{7pt}}

\newtcolorbox{aretenir}[1][]{popbase,colframe=popgreen,before upper={\popboxhead{À retenir}{popgreen}{#1}}}
% Texte en espagnol (document de lecture, dialogue, poème, extrait) : cadre
% d'étude. Un titre optionnel (source, auteur) s'affiche dans l'en-tête.
\newtcolorbox{textebox}[1][]{popbase,colframe=popblue,before upper={\popboxhead{Texto}{popblue}{#1}\begin{otherlanguage}{spanish}},after upper={\end{otherlanguage}}}
% Marques ✓ / ✗ (forme correcte / fautive) souvent utilisées par les modèles
\providecommand{\cross}{\textcolor{poppink}{\ensuremath{\times}}}
\providecommand{\xmark}{\cross}\providecommand{\crossmark}{\cross}\providecommand{\cmark}{\ensuremath{\checkmark}}
% Ligne de réponse / justification à compléter (inventée par les modèles)
\providecommand{\justification}[1]{\par\noindent\textit{Justification~:} \dotfill\par}
% \textipa (package tipa non chargé) : la police ne couvre pas l'API -> simple passage
\providecommand{\textipa}[1]{#1}
% Soulignement (ulem non chargé) : \uline, \ul
\providecommand{\uline}[1]{\underline{#1}}\providecommand{\ul}[1]{\underline{#1}}
% Ordinaux français écrits en macro par les modèles : 1\ière, 3\ième
\newcommand{\ière}{\textsuperscript{re}}\newcommand{\ième}{\textsuperscript{e}}
% \exemple (inventée par les modèles) : étiquette « Exemple : »
\providecommand{\exemple}{\textbf{Exemple~:}\ }
% Mot ou phrase en espagnol à l'intérieur d'un paragraphe français
\newcommand{\esp}[1]{\ifmmode\text{\textit{#1}}\else\begingroup\selectlanguage{spanish}\itshape #1\endgroup\fi}
\newtcolorbox{exemplebox}[1][]{popbase,colframe=poporange,before upper={\popboxhead{Exemple}{poporange}{#1}}}
\newtcolorbox{definition}[1][]{popbase,colframe=popblue,before upper={\popboxhead{Définition}{popblue}{#1}}}
\newtcolorbox{propriete}[1][]{popbase,colframe=poppurple,before upper={\popboxhead{Propriété}{poppurple}{#1}}}
\newtcolorbox{methode}[1][]{popbase,colframe=popgold,before upper={\popboxhead{Méthode}{popgold}{#1}}}
\newtcolorbox{attention}[1][]{popbase,colframe=poppink,before upper={\popboxhead{Attention}{poppink}{#1}}}
\newtcolorbox{experience}[1][]{popbase,colframe=popblue,colback=popblueL,before upper={\popboxhead{Expérience}{popblue}{#1}}}
% « Le savais-tu ? » : carte violette pleine (contexte, culture, Cameroun)
\newtcolorbox{savaistu}[1][]{popbase,colback=poppurple,colframe=popink,
  fontupper=\popsemi\fontsize{10.4}{14.2}\selectfont\color{white},
  before upper={\poppill{Le savais-tu\,?}{white}{poppurple}\ifx\relax#1\relax\else\hspace{6pt}{\popxbold\color{white}#1}\fi\par\vspace{7pt}}}
% Exercice résolu : énoncé en haut, \tcblower puis la solution sur fond crème
\newcounter{popexo}\newcounter{popexr}
\newtcolorbox{exoresolu}[1][]{popbase,colframe=poporange,colbacklower=popcream,
  segmentation style={popink,line width=1.2pt,dashed},
  before upper={\stepcounter{popexr}\popboxhead{Exercice résolu \thepopexr}{poporange}{#1}},
  before lower={\poppill{Solution}{popink}{popcream}\par\vspace{6pt}}}
% Exercice (non résolu en place) — la correction est dans la partie Corrigés
\newtcolorbox{exercice}[1][]{popbase,colframe=popink,
  before upper={\stepcounter{popexo}\popboxhead{Exercice \thepopexo}{popink}{#1}}}
% Corrigé
\newtcolorbox{corrige}[1][]{popbase,colframe=popgreen,colback=popgreenL,
  before upper={\popboxhead{Corrigé}{popgreen}{#1}}}
% Objectifs / prérequis / bilan
\newtcolorbox{objectifs}{popbase,colframe=popink,colback=poporangeL,
  before upper={\popboxhead{Objectifs de la leçon}{popink}{}}}
\newtcolorbox{prerequis}{popbase,colframe=popink,colback=popblueL,
  before upper={\popboxhead{Prérequis}{popblue}{}}}
\newtcolorbox{bilan}{popbase,colframe=popink,colback=white,boxrule=2.8pt,
  before upper={\popboxhead{Fiche bilan}{poporange}{L'essentiel à connaître pour l'examen}}}
% Figure encadrée avec légende
\newtcolorbox{popfigure}[1][]{enhanced,breakable=false,colback=white,colframe=popline,boxrule=1.4pt,arc=8pt,
  left=10pt,right=10pt,top=10pt,bottom=8pt,before skip=10pt,after skip=12pt,halign=center,
  lower separated=false,halign lower=center,
  fontlower=\popsemi\fontsize{9}{11.5}\selectfont\color{popmuted},
  #1}
\newcommand{\legende}[1]{\par\vspace{6pt}{\popsemi\fontsize{9}{11.5}\selectfont\color{popmuted}#1}}

% ============================================================
% Signature OnBuch+
% ============================================================
\newcommand{\poplogoinline}[1]{{\poplogo\fontsize{#1}{#1}\selectfont\color{popink}OnBuch\color{poporange}+}}
\newcommand{\signatureonbuch}{%
  \begin{tcolorbox}[enhanced,colback=popink,colframe=popink,boxrule=2.2pt,arc=10pt,
      drop shadow=poporange,left=14pt,right=14pt,top=10pt,bottom=10pt,before skip=0pt,after skip=0pt]
    \tikz[baseline=-0.7ex]\node[circle,fill=popgreen,draw=popcream,line width=1.3pt,minimum size=22pt,inner sep=0]{\popcheck{white}};%
    \hspace{9pt}\begin{minipage}[c]{0.62\linewidth}
      {\popxbold\fontsize{12}{13}\selectfont\color{popcream}Signé\ \textbullet\ {\poplogo OnBuch\color{poporange}+}}\par\vspace{2pt}
      {\raggedright\popsemi\fontsize{8}{10.5}\selectfont\color{popline}\MakeUppercase{Équipe pédagogique OnBuch+}\\\MakeUppercase{Conforme au programme officiel MINESEC}\par}
    \end{minipage}\hfill
    {\poplogo\fontsize{22}{22}\selectfont\color{popcream}OnBuch\color{poporange}+}
  \end{tcolorbox}%
}

% ============================================================
% Page de garde
% #1 titre · #2 matière · #3 niveau · #4 module · #5 n° leçon · #6 durée ·
% #7 description / objectifs (texte ou itemize)
% ============================================================
\newcommand{\pagegarde}[7]{%
  \thispagestyle{empty}
  \newgeometry{a4paper,margin=1.7cm,top=1.6cm,bottom=1.4cm}%
  \begin{tikzpicture}[remember picture,overlay]
    \fill[poporange!18] ([xshift=-3.2cm,yshift=-3.6cm]current page.north east) circle (4.6cm);
    \fill[poppurple!14] ([xshift=2.4cm,yshift=7.6cm]current page.south west) circle (3.8cm);
  \end{tikzpicture}%
  {\poplogo\fontsize{30}{30}\selectfont\color{popink}OnBuch\color{poporange}+}\hfill
  \raisebox{0.6ex}{\poppill{Cours signé \textbullet\ Premium}{popink}{popcream}}\par
  \vspace{1pt}\popsquiggle{poppurple}\par
  \vspace{8pt}
  \begin{tcolorbox}[enhanced,colback=poporange,colframe=popink,boxrule=2.8pt,arc=13pt,
      drop shadow=popink,left=18pt,right=18pt,top=12pt,bottom=13pt,after skip=10pt]
    \poppill{#2 \textbullet\ #3}{popink}{popcream}\hspace{5pt}\poppill{#4}{white}{popink}\par\vspace{8pt}
    {\popdisplay\fontsize{14}{14}\selectfont\color{popink}LEÇON #5}\par\vspace{4pt}
    {\raggedright\hyphenpenalty=10000\exhyphenpenalty=10000\popdisplay\fontsize{25}{27}\selectfont\color{popcream}\MakeUppercase{#1}\par}
    \vspace{8pt}
    {\popxbold\fontsize{9.5}{9.5}\selectfont\color{popink}\MakeUppercase{Durée conseillée : #6}}
  \end{tcolorbox}
  \begin{tcolorbox}[enhanced,colback=white,colframe=popink,boxrule=2.2pt,arc=10pt,
      drop shadow=popink,left=16pt,right=16pt,top=10pt,bottom=10pt,after skip=8pt]
    \poppill{Dans cette leçon}{poppurple}{white}\par\vspace{5pt}
    \popsemi\fontsize{9.3}{12.6}\selectfont\color{popink}#7
  \end{tcolorbox}
  \noindent\begin{minipage}[t]{0.487\linewidth}
    \begin{tcolorbox}[enhanced,colback=popgreen,colframe=popink,boxrule=2.2pt,arc=10pt,
        drop shadow=popink,left=11pt,right=11pt,top=10pt,bottom=10pt,height=3.1cm,valign=center]
      \tcbox[colback=white,colframe=popink,boxrule=1.3pt,arc=5pt,boxsep=0pt,nobeforeafter,
        left=3pt,right=3pt,top=3pt,bottom=3pt,valign=center]{\qrcode[height=1.9cm]{https://play.google.com/store/apps/details?id=com.luvvix.onbuch&hl=fr}}%
      \hspace{8pt}\begin{minipage}[c]{0.5\linewidth}
        {\popdisplay\fontsize{11}{12.5}\selectfont\color{white}\MakeUppercase{Télécharge OnBuch}}\par\vspace{4pt}
        {\raggedright\popsemi\fontsize{8.4}{11}\selectfont\color{white}Gratuit \textbullet\ Google Play\\Tuteur IA, annales, concours, résultats\par}
      \end{minipage}
    \end{tcolorbox}
  \end{minipage}\hfill
  \begin{minipage}[t]{0.487\linewidth}
    \begin{tcolorbox}[enhanced,colback=poppurple,colframe=popink,boxrule=2.2pt,arc=10pt,
        drop shadow=popink,left=11pt,right=11pt,top=10pt,bottom=10pt,height=3.1cm,valign=center]
      \tikz[baseline=-0.7ex]\node[circle,fill=white,draw=popink,line width=1.3pt,minimum size=1.6cm,inner sep=0]{\popdisplay\fontsize{24}{24}\selectfont\color{poppurple}+};%
      \hspace{8pt}\begin{minipage}[c]{0.6\linewidth}
        {\popdisplay\fontsize{11}{12.5}\selectfont\color{white}\MakeUppercase{Passe à OnBuch+}}\par\vspace{4pt}
        {\raggedright\popsemi\fontsize{8.4}{11}\selectfont\color{white}Tous les cours signés, Tuteur IA illimité, fiches PDF — onglet OnBuch+ de l'app\par}
      \end{minipage}
    \end{tcolorbox}
  \end{minipage}
  \vspace{8pt}
  \popcontenu
  \vfill
  \signatureonbuch
  \restoregeometry
  \newpage
}

% Rangée « Ce cours contient » (page de garde)
\newcommand{\poptile}[3]{% #1 couleur #2 gros texte #3 légende
  \begin{tcolorbox}[enhanced,colback=#1,colframe=popink,boxrule=2pt,arc=9pt,shadow={2.5pt}{-2.5pt}{0pt}{popink},
      left=6pt,right=6pt,top=9pt,bottom=9pt,halign=center,valign=center,height=1.75cm,width=\linewidth,nobeforeafter]
    {\popdisplay\fontsize{12.5}{13}\selectfont\color{white}\MakeUppercase{#2}}\par\vspace{3pt}
    {\centering\popsemi\fontsize{7.8}{9.5}\selectfont\color{white}#3\par}
  \end{tcolorbox}}
\newcommand{\popcontenu}{%
  \par\noindent{\popxboldls\fontsize{7.5}{7.5}\selectfont\color{popmuted}CE COURS SIGNÉ CONTIENT}\par\vspace{7pt}
  \noindent\begin{minipage}[t]{0.235\linewidth}\poptile{popblue}{Cours}{complet, illustré, pas à pas}\end{minipage}\hfill
  \begin{minipage}[t]{0.235\linewidth}\poptile{poporange}{Méthodes}{et exercices résolus}\end{minipage}\hfill
  \begin{minipage}[t]{0.235\linewidth}\poptile{poppink}{Exercices}{type examen + corrigés}\end{minipage}\hfill
  \begin{minipage}[t]{0.235\linewidth}\poptile{popgreen}{Fiche bilan}{l'essentiel à retenir}\end{minipage}\par}

% Sommaire
\newtcolorbox{sommaire}{enhanced,colback=white,colframe=popink,boxrule=2.2pt,arc=10pt,
  drop shadow=popink,left=18pt,right=18pt,top=13pt,bottom=13pt,before skip=6pt,after skip=20pt,
  before upper={\poppill{Sommaire}{poppurple}{white}\par\vspace{9pt}\popsemi\fontsize{10.6}{14.5}\selectfont\color{popink}}}

% Signature de fin de cours — poussée tout en bas de la dernière page
\newcommand{\findecours}{%
  \par\vfill\nopagebreak
  \begin{center}\popsquiggle{poporange}\end{center}\vspace{4pt}
  \signatureonbuch
}
\AtBeginDocument{\def\llbracket{\mathopen{[\mkern-3.5mu[}}\def\rrbracket{\mathclose{]\mkern-3.5mu]}}} % intervalles d'entiers (glyphe ⟦ absent de la police)
% Glyphes absents de Fira Math : redéfinis à partir de glyphes disponibles
\AtBeginDocument{%
  \def\setminus{\mathbin{\backslash}}\let\smallsetminus\setminus
  \def\vdots{\mathord{\vbox{\baselineskip3.2pt\lineskiplimit0pt\kern4pt\hbox{.}\hbox{.}\hbox{.}}}}%
  \def\ddots{\mathinner{\mkern1mu\raise6.5pt\hbox{.}\mkern2mu\raise3.6pt\hbox{.}\mkern2mu\raise0.7pt\hbox{.}\mkern1mu}}%
  \def\triangle{\mathord{\scalebox{0.9}{$\symup{\Delta}$}}}%
  \def\nabla{\mathord{\raisebox{\depth}{\scalebox{1}[-1]{$\symup{\Delta}$}}}}%
  \def\checkmark{\popcheck{popdark}}%
  \def\star{\mathbin{\ast}}\def\bigstar{\mathord{\ast}}%
  \def\diamond{\mathbin{\scalebox{0.6}{$\Diamond$}}}%
  \def\bigcup{\mathop{\mathchoice{\raisebox{-0.3em}{\scalebox{1.7}{$\cup$}}}{\scalebox{1.25}{$\cup$}}{\cup}{\cup}}\displaylimits}%
  \def\bigcap{\mathop{\mathchoice{\raisebox{-0.3em}{\scalebox{1.7}{$\cap$}}}{\scalebox{1.25}{$\cap$}}{\cap}{\cap}}\displaylimits}%
  \def\lesseqqgtr{\mathrel{\lessgtr}}\def\bigoplus{\mathop{\oplus}\displaylimits}%
}
\providecommand{\ssi}{si et seulement si} % abréviation fréquente
\AtBeginDocument{\providecommand{\wideparen}{\overparen}} % arc de cercle

%% FIGFIX-BEGIN v5 — correctifs globaux des figures (docs/onbuch-generation/tools/figfix)
\usepackage{adjustbox}\usepackage{etoolbox}\tcbuselibrary{raster}
% toute figure TikZ/pgfplots plus large que la ligne est réduite à la largeur disponible
\BeforeBeginEnvironment{tikzpicture}{\begin{adjustbox}{max width=\linewidth}}
\AfterEndEnvironment{tikzpicture}{\end{adjustbox}}
% légendes pgfplots lisibles par-dessus les courbes
\pgfplotsset{every axis/.append style={legend style={fill=white,fill opacity=0.92,text opacity=1,draw=black!15,inner sep=3pt}}}
% étiquettes d'arêtes (syntaxe edge["..."]) avec fond blanc
\tikzset{every edge quotes/.append style={fill=white,inner sep=1.2pt}}
%% FIGFIX-END
\begin{document}

\pagegarde{La correspondencia y los intercambios escritos (continuación) : cartas formales y administrativas}{Espagnol}{1ère A}{Módulo 5 : Les mass médias et la communication}{EA20}{2 h}{Cette leçon mixte poursuit l'étude de la correspondance en abordant les lettres formelles et administratives : lettre de demande, lettre de réclamation et lettre de motivation. Elle articule lecture d'un texte support, lexique thématique, grammaire (conditionnel de courtoisie) et traduction de formules pour aboutir à la rédaction guidée d'une lettre de motivation de 100 mots.
\vspace{4pt}
\begin{multicols}{2}\small
\begin{itemize}
\item À la fin de cette leçon, je sais identifier la structure d'une lettre formelle espagnole (en-tête, destinatario, cuerpo, despedida)
\item À la fin de cette leçon, je sais distinguer lettre de demande, lettre de réclamation et lettre de motivation
\item À la fin de cette leçon, je sais utiliser le conditionnel de politesse (quisiera, me gustaría, le agradecería que + subjonctif)
\item À la fin de cette leçon, je sais employer les formules d'ouverture et de clôture de la correspondance formelle
\item À la fin de cette leçon, je sais traduire les formules de correspondance du français vers l'espagnol et inversement
\item À la fin de cette leçon, je sais rédiger une lettre de demande formelle complète et correcte
\end{itemize}\end{multicols}}

\begin{sommaire}
\begin{multicols}{2}
\begin{enumerate}
\item Texte support : una carta de solicitud de prácticas
\item Lexique thématique de la correspondance formelle
\item Grammaire : le conditionnel de courtoisie
\item Traduction des formules de correspondance
\item La lettre de réclamation
\item Production guidée : la lettre de motivation (100 mots)
\item Activité d'intégration
\item Méthodes et exercices résolus
\item Exercices
\item Corrigés des exercices
\item Fiche bilan
\end{enumerate}
\end{multicols}
\end{sommaire}

\begin{objectifs}
\begin{itemize}
\item À la fin de cette leçon, je sais identifier la structure d'une lettre formelle espagnole (en-tête, destinatario, cuerpo, despedida)
\item À la fin de cette leçon, je sais distinguer lettre de demande, lettre de réclamation et lettre de motivation
\item À la fin de cette leçon, je sais utiliser le conditionnel de politesse (quisiera, me gustaría, le agradecería que + subjonctif)
\item À la fin de cette leçon, je sais employer les formules d'ouverture et de clôture de la correspondance formelle
\item À la fin de cette leçon, je sais traduire les formules de correspondance du français vers l'espagnol et inversement
\item À la fin de cette leçon, je sais rédiger une lettre de demande formelle complète et correcte
\item À la fin de cette leçon, je sais rédiger une lettre de motivation simple d'environ 100 mots
\item À la fin de cette leçon, je sais formuler une réclamation polie mais ferme en espagnol
\end{itemize}
\end{objectifs}

\begin{prerequis}
\begin{itemize}
\item La correspondance informelle (carta informal) vue en Seconde : salutations, despedidas
\item Le présent de l'indicatif et le présent du subjonctif (1ère, leçons antérieures)
\item Le conditionnel simple (formation de base vue en début d'année)
\item Le vocabulaire des études et des métiers (Seconde)
\item Les pronoms compléments d'objet direct et indirect (le, la, les)
\end{itemize}
\end{prerequis}

\newpage

\coursec{Texte support : una carta de solicitud de prácticas}

Mariam, élève de Première à Yaoundé, veut faire un stage dans une entreprise de téléphonie mobile pendant les grandes vacances. Elle rédige une lettre de demande, mais son directeur la lui renvoie avec cette remarque : « Ton ton est trop familier, et tu as oublié les formules de politesse ! » En espagnol comme en français, une lettre formelle obéit à des codes précis : en-tête, formules de courtoisie, conditionnel de politesse. Que ce soit pour demander un stage, s'inscrire à l'université de Salamanque ou se plaindre d'un service, savoir écrire une \esp{carta formal} ouvre toutes les portes. \textbf{Problématique :} quels sont les éléments obligatoires d'une lettre formelle espagnole, et quelles formules permettent de formuler une demande avec politesse et efficacité ?

\courssub{Lecture du texte : una carta de solicitud de prácticas}

\begin{textebox}[Una carta de solicitud de prácticas]
Yaoundé, 15 de mayo de 2024\par
\medskip
Señor Director de Recursos Humanos\\
Empresa Telecam S.A.\\
Avenida de la Independencia, 45\\
Yaoundé\par
\medskip
\textbf{Asunto:} solicitud de prácticas de verano\par
\medskip
Estimado señor director:\par
\medskip
Me dirijo a usted con el fin de solicitar unas prácticas en su empresa durante los meses de julio y agosto. Actualmente curso el bachillerato en el Liceo de Yaoundé y quisiera adquirir experiencia profesional en el sector de las telecomunicaciones, un campo que me apasiona desde hace varios años.\par
\medskip
Durante mis estudios, he desarrollado competencias en informática y en comunicación. Además, hablo español, francés e inglés, lo que me permitiría colaborar con clientes de diferentes países. Estoy convencido de que estas prácticas serían una oportunidad excelente para mi formación y, al mismo tiempo, podría aportar mi energía y mi motivación a su equipo.\par
\medskip
Le agradecería que me enviara información sobre las condiciones de admisión y los documentos necesarios para completar mi expediente. Adjunto a la presente mi currículum vítae y una carta de recomendación de mi profesor de español. Quedo a su disposición para una entrevista en el horario que usted considere oportuno.\par
\medskip
En espera de su respuesta, le saluda atentamente,\par
\medskip
Andrés Mbarga
\end{textebox}

\begin{definition}[Glossaire du texte]
\begin{itemize}
\item \esp{la solicitud} : la demande (administrative). \esp{solicitar} : demander, solliciter.
\item \esp{las prácticas} : le stage (en entreprise). \esp{hacer unas prácticas} : faire un stage.
\item \esp{adjunto / adjunta} : ci-joint(e). \esp{Adjunto a la presente...} : Vous trouverez ci-joint...
\item \esp{el expediente} : le dossier (administratif, scolaire).
\item \esp{agradecer} : remercier. \esp{Le agradecería que...} : Je vous serais reconnaissant de...
\item \esp{atentamente} : cordialement, avec mes salutations distinguées.
\item \esp{el asunto} : l'objet (de la lettre).
\item \esp{el currículum vítae} : le curriculum vitae (CV).
\item \esp{quedar a su disposición} : se tenir à votre disposition.
\end{itemize}
\end{definition}

\begin{attention}[Faux amis et pièges de traduction]
\begin{itemize}
\item \esp{las prácticas} ne signifie pas « les pratiques » au sens abstrait : c'est le mot courant pour \cle{stage} en Espagne. En Amérique latine, on dit aussi \esp{la pasantía}.
\item \esp{asistir a} ne veut pas dire « assister » au sens d'aider : \esp{asistir a una reunión} = assister à (être présent à) une réunion. « Aider » se dit \esp{ayudar}.
\item \esp{demandar} en espagnol juridique signifie « poursuivre en justice » ; pour « demander », on utilise \esp{solicitar} ou \esp{pedir}.
\item Ne traduis jamais « actuellement » par \esp{actualmente}... si, justement, c'est correct ! \esp{actualmente} = actuellement. Mais attention : \esp{actual} = actuel, présent, pas « actuel » au sens de « vrai ».
\end{itemize}
\end{attention}

\courssub{La structure de la lettre formelle}

Observons d'abord comment la lettre d'Andrés est organisée, puis énonçons la règle générale.

\begin{propriete}[Les sept blocs obligatoires de la carta formal]
Toute lettre formelle espagnole comprend, dans cet ordre :
\begin{enumerate}
\item \esp{Lugar y fecha} : le lieu et la date (\esp{Yaoundé, 15 de mayo de 2024}). En espagnol, la date s'écrit : ville, virgule, jour + \esp{de} + mois + \esp{de} + année. Le mois ne prend \textbf{pas} de majuscule.
\item \esp{Destinatario} : le nom et l'adresse du destinataire, avec sa fonction (\esp{Señor Director de Recursos Humanos}).
\item \esp{Asunto} : l'objet de la lettre, en une phrase courte.
\item \esp{Saludo formal} : la formule d'appel (\esp{Estimado señor director:}).
\item \esp{Cuerpo} : le corps de la lettre, en général trois paragraphes (présentation de la demande ; arguments ; demande d'information et pièces jointes).
\item \esp{Despedida} : la formule de politesse finale (\esp{En espera de su respuesta, le saluda atentamente}).
\item \esp{Firma} : la signature (prénom et nom).
\end{enumerate}
\end{propriete}

\begin{popfigure}
\begin{center}
\begin{tikzpicture}[
  bloque/.style={draw=popblue, fill=popblueL, rounded corners=3pt, minimum width=8.6cm, minimum height=0.72cm, align=center, font=\small},
  etiqueta/.style={font=\small\itshape, text=popmuted, right=0.25cm}
]
\node[bloque] (b1) at (0,0) {\textbf{Lugar y fecha} : Yaoundé, 15 de mayo de 2024};
\node[bloque, fill=poporangeL, draw=poporange] (b2) at (0,-0.95) {\textbf{Destinatario} : Señor Director de Recursos Humanos};
\node[bloque, fill=popgreenL, draw=popgreen] (b3) at (0,-1.9) {\textbf{Asunto} : solicitud de prácticas de verano};
\node[bloque, fill=poppurpleL, draw=poppurple] (b4) at (0,-2.85) {\textbf{Saludo} : Estimado señor director:};
\node[bloque, fill=popgoldL, draw=popgold, minimum height=1.5cm] (b5) at (0,-4.15) {\textbf{Cuerpo} : 1.º la demande \quad 2.º los argumentos \quad 3.º información y adjuntos};
\node[bloque, fill=poppinkL, draw=poppink] (b6) at (0,-5.35) {\textbf{Despedida} : En espera de su respuesta, le saluda atentamente};
\node[bloque] (b7) at (0,-6.3) {\textbf{Firma} : Andrés Mbarga};
\foreach \a/\b in {b1/b2, b2/b3, b3/b4, b4/b5, b5/b6, b6/b7}
  \draw[-{Stealth}, popblue, thick] (\a.south) -- (\b.north);
\end{tikzpicture}
\end{center}
\legende{Plan d'une lettre formelle espagnole : les sept blocs dans l'ordre obligatoire.}
\end{popfigure}

\begin{savaistu}[Les deux points après le saludo]
En espagnol, le \esp{saludo} formel se termine obligatoirement par \textbf{deux points} : \esp{Estimado señor director:} et non par une virgule comme en français (« Cher Monsieur, »). Après les deux points, on va à la ligne et on commence le premier paragraphe par une majuscule. Autre différence : en espagnol, on écrit \esp{Estimado señor Mbarga} avec le \textbf{nom de famille}, alors qu'en français on écrit « Cher Monsieur » sans nom. Enfin, les mois de l'année ne prennent jamais de majuscule en espagnol : \esp{mayo}, \esp{julio}, \esp{agosto}.
\end{savaistu}

\courssub{Compréhension globale du texte}

\begin{exercice}[Questions de compréhension globale]
Réponds en français, puis formule la réponse en espagnol simple.
\begin{enumerate}
\item ¿Quién escribe la carta? (Qui écrit la lettre ?)
\item ¿A quién va dirigida? (À qui est-elle adressée ?)
\item ¿Por qué escribe Andrés? (Pourquoi écrit-il ?)
\item ¿Qué documentos adjunta a su carta? (Quelles pièces jointes accompagnent la lettre ?)
\end{enumerate}
\end{exercice}

\begin{corrige}[Exercice : compréhension globale]
\begin{enumerate}
\item Escribe Andrés Mbarga, un alumno de bachillerato del Liceo de Yaoundé. (Andrés Mbarga, un lycéen de terminale au Lycée de Yaoundé.)
\item La carta va dirigida al señor Director de Recursos Humanos de la empresa Telecam S.A. (Au directeur des ressources humaines de l'entreprise Telecam.)
\item Escribe para solicitar unas prácticas durante los meses de julio y agosto. (Pour demander un stage pendant juillet et août.)
\item Adjunta su currículum vítae y una carta de recomendación de su profesor de español. (Son CV et une lettre de recommandation de son professeur d'espagnol.)
\end{enumerate}
\end{corrige}

\courssub{Compréhension fine : les formules de politesse}

Relisons le texte à la loupe. Une lettre formelle se reconnaît à ses \cle{formules figées} et à l'emploi du \cle{conditionnel de courtoisie}. Relevons-les.

\begin{exemplebox}[Formules de politesse relevées dans le texte]
\begin{itemize}
\item \esp{Me dirijo a usted con el fin de...} = « Je m'adresse à vous afin de... » (ouverture classique d'une demande).
\item \esp{Quisiera adquirir experiencia profesional...} = « Je voudrais acquérir une expérience professionnelle... » (conditionnel de courtoisie).
\item \esp{Le agradecería que me enviara información...} = « Je vous serais reconnaissant de bien vouloir m'envoyer des informations... ».
\item \esp{Adjunto a la presente mi currículum vítae.} = « Vous trouverez ci-joint mon CV. »
\item \esp{Quedo a su disposición para una entrevista.} = « Je me tiens à votre disposition pour un entretien. »
\item \esp{En espera de su respuesta, le saluda atentamente.} = « Dans l'attente de votre réponse, je vous prie d'agréer mes salutations distinguées. »
\end{itemize}
\end{exemplebox}

\begin{aretenir}[Le conditionnel de courtoisie : aperçu]
Dans une lettre formelle, on adoucit la demande avec le \textbf{conditionnel} : \esp{quisiera} (je voudrais), \esp{me gustaría} (j'aimerais), \esp{le agradecería} (je vous serais reconnaissant). Comparer :
\begin{itemize}
\item \esp{Quiero información.} « Je veux des informations. » $\rightarrow$ direct, presque impoli.
\item \esp{Quisiera información.} « Je voudrais des informations. » $\rightarrow$ poli, adapté au registre formel.
\end{itemize}
Après \esp{le agradecería que}, l'espagnol exige l'\textbf{imparfait du subjonctif} : \esp{Le agradecería que me \textbf{enviara} información.} Nous étudierons cette grammaire en détail dans la section 3.
\end{aretenir}

\begin{center}
\begin{tabularx}{\linewidth}{|X|X|X|}
\hline
\rowcolor{popblueL} \textbf{Français} & \textbf{Español} & \textbf{Remarque} \\
\hline
Monsieur le Directeur, & Estimado señor director: & deux points obligatoires \\
\hline
Je me permets de vous écrire pour... & Me dirijo a usted con el fin de... & formule d'ouverture \\
\hline
Je voudrais / J'aimerais & Quisiera / Me gustaría & conditionnel de courtoisie \\
\hline
Je vous serais reconnaissant de... & Le agradecería que + subj. impf. & registre très soutenu \\
\hline
Vous trouverez ci-joint... & Adjunto a la presente... & \esp{adjunto} s'accorde : \esp{adjunta una carta} \\
\hline
Dans l'attente de votre réponse... & En espera de su respuesta... & formule de clôture \\
\hline
Je vous prie d'agréer... & Le saluda atentamente & 3\textsuperscript{e} personne en espagnol \\
\hline
\end{tabularx}
\end{center}

\courssub{Commentaire guidé : registre, ton, efficacité}

\begin{methode}[Commenter une lettre formelle en quatre étapes]
\begin{enumerate}
\item \textbf{Identifier l'émetteur et le destinataire} : qui écrit ? à qui ? quelle relation hiérarchique ? (Ici : un lycéen $\rightarrow$ un directeur ; relation de subordination, donc registre soutenu et vouvoiement \esp{usted}.)
\item \textbf{Analyser la structure} : vérifier la présence des sept blocs (lugar y fecha, destinatario, asunto, saludo, cuerpo, despedida, firma) et la logique des trois paragraphes du corps.
\item \textbf{Relever le registre et les formules de politesse} : souligner les conditionnels de courtoisie (\esp{quisiera}, \esp{agradecería}), les formules figées (\esp{Me dirijo a usted...}, \esp{Adjunto a la presente...}) et l'absence de tutoiement ou d'abréviations familières.
\item \textbf{Évaluer l'efficacité de la demande} : la demande est-elle claire dès le premier paragraphe ? Les arguments sont-ils pertinents (compétences, langues, motivation) ? La lettre facilite-t-elle la réponse (documents joints, disponibilité pour un entretien) ?
\end{enumerate}
\end{methode}

Appliquons la méthode. Le ton de la lettre d'Andrés est \textbf{respectueux sans être obséquieux} : le conditionnel adoucit chaque demande, et le vouvoiement (\esp{usted}, \esp{le}, \esp{su}) est constant. L'efficacité persuasive repose sur trois atouts : une demande précise dès la première phrase, des arguments concrets (informatique, trois langues), et une conclusion qui facilite la décision du directeur (CV joint, disponibilité totale). C'est exactement ce qui manquait à la lettre de Mariam dans notre situation d'accroche !

\begin{exoresolu}[Repérage des formules]
\textbf{Énoncé.} Dans la phrase suivante, souligne mentalement la formule de politesse et le verbe au conditionnel, puis traduis : \esp{Le agradecería que me enviara información sobre las condiciones de admisión.}
\tcblower
\textbf{Solution détaillée.} La formule de politesse est \esp{Le agradecería que...} ; le verbe au conditionnel est \esp{agradecería} (verbe \esp{agradecer}, 1\textsuperscript{re} personne du singulier du conditionnel). Le verbe \esp{enviara} est à l'imparfait du subjonctif, exigé après cette formule. Traduction : « Je vous serais reconnaissant de bien vouloir m'envoyer des informations sur les conditions d'admission. » Piège : ne pas traduire mot à mot par « Je vous remercierais que vous m'envoyiez », qui est lourd en français ; préférer la tournure idiomatique « Je vous serais reconnaissant de... ».
\end{exoresolu}

\begin{attention}[Erreurs fréquentes des francophones]
\begin{itemize}
\item \textbf{La virgule après le saludo} : on écrit \esp{Estimado señor director:} avec deux points, jamais avec une virgule.
\item \textbf{Le tutoiement parasite} : une fois \esp{usted} choisi, toute la lettre reste à la 3\textsuperscript{e} personne : \esp{le}, \esp{su}, \esp{usted}. Ne jamais mélanger \esp{tú} et \esp{usted}.
\item \textbf{Le calque « Je vous demande de... »} : en espagnol formel, on préfère \esp{Me dirijo a usted con el fin de...} ou \esp{Le escribo para solicitar...}.
\item \textbf{La clôture à la 1\textsuperscript{re} personne} : on écrit \esp{le saluda atentamente} (3\textsuperscript{e} personne, « celui qui vous salue »), formule figée, et non \esp{le saludo atentamente}, qui est tolérée mais moins classique.
\end{itemize}
\end{attention}

\begin{experience}[À toi de jouer : corrige la lettre de Mariam]
Par groupes de trois, réécrivez oralement le début de la lettre de Mariam en appliquant les sept blocs : imaginez le lieu et la date, le destinataire (une entreprise de téléphonie mobile de Douala), l'objet, le saludo correct avec deux points, et la première phrase de demande avec \esp{Me dirijo a usted con el fin de...}. Chaque groupe lit sa version à la classe ; les autres groupes vérifient la présence des sept blocs sur la grille du schéma.
\end{experience}

\coursec{Lexique thématique de la correspondance formelle}

Dans la section précédente, nous avons lu et commenté une \esp{carta de solicitud de prácticas}, la lettre de demande de stage d'un élève camerounais adressée à une entreprise. Nous y avons repéré des expressions comme \esp{Estimado señor director}, \esp{Le agradecería que me enviara información} ou encore \esp{En espera de su respuesta, le saluda atentamente}. Ces formules ne sont pas choisies au hasard : elles appartiennent à un \cle{lexique codifié} de la correspondance formelle. C'est ce lexique que nous allons maintenant étudier méthodiquement, champ par champ, afin de pouvoir rédiger nous-mêmes des lettres impeccables.

\begin{definition}[La carta formal]
Une \esp{carta formal} (lettre formelle ou administrative) est une lettre adressée à une \cle{administration}, une \cle{entreprise} ou une \cle{personne que l'on ne connaît pas} (directeur, responsable, fonctionnaire). Elle obéit à des règles strictes : registre \cle{soutenu}, emploi de \esp{usted} (jamais \esp{tú}), formules de politesse codifiées, structure fixe (en-tête, salutation, corps, clôture, signature). Elle s'oppose à la \esp{carta informal} (lettre à un ami, à la famille), étudiée dans la leçon précédente.
\end{definition}

\courssub{L'en-tête et les références}

Observons d'abord le haut de la lettre de demande de stage de la section 1 : en haut à gauche figurent le nom et l'adresse de l'expéditeur ; en haut à droite, le nom et l'adresse du destinataire ; puis une ligne \esp{Asunto: solicitud de prácticas}. Chaque élément porte un nom précis en espagnol.

\begin{definition}[Vocabulaire de l'en-tête]
\begin{itemize}
\item \esp{el remitente} : l'expéditeur, la personne qui écrit et envoie la lettre.
\item \esp{el destinatario} : le destinataire, la personne ou l'organisme qui reçoit la lettre.
\item \esp{la dirección} : l'adresse. \esp{el código postal} : le code postal.
\item \esp{el asunto} : l'objet de la lettre (ce dont elle parle).
\item \esp{la referencia} : la référence ; dans les échanges entre administrations on écrit \esp{su ref.} (votre référence) et \esp{nuestra ref.} (notre référence).
\end{itemize}
\end{definition}

\begin{center}
\begin{tabularx}{\linewidth}{|X|X|X|}
\hline
\rowcolor{popblueL} Español & Français & Exemple en contexte \\
\hline
el remitente & l'expéditeur & \esp{El remitente escribe su dirección arriba a la izquierda.} \\
\hline
el destinatario & le destinataire & \esp{El destinatario es el director de la empresa.} \\
\hline
la dirección & l'adresse & \esp{Mi dirección es: B.P. 1200, Yaoundé.} \\
\hline
el código postal & le code postal & \esp{¿Cuál es el código postal de Madrid?} \\
\hline
el asunto & l'objet & \esp{Asunto: solicitud de información.} \\
\hline
su ref. / nuestra ref. & votre réf. / notre réf. & \esp{Su ref.: carta del 12 de marzo.} \\
\hline
\end{tabularx}
\end{center}

\begin{attention}[« Asunto » n'est pas « sujet »]
Le mot \esp{asunto} est un faux ami partiel : il ne signifie pas « sujet » au sens scolaire (sujet de dissertation, sujet de grammaire), mais \cle{objet} de la lettre, \cle{affaire}, \cle{question}. On ne traduit donc jamais « Objet : demande de stage » par \esp{Objeto: demanda de prácticas} (double erreur : \esp{objeto} et \esp{demanda} au lieu de \esp{solicitud}) mais par \esp{Asunto: solicitud de prácticas}. De même, « le sujet de la phrase » se dit \esp{el sujeto de la oración}, jamais \esp{el asunto}.
\end{attention}

\courssub{Les salutations formelles}

Après l'en-tête vient la \esp{salutación} ou \esp{fórmula de encabezamiento}. En espagnol formel, elle se termine toujours par \cle{deux-points} (:) et la première ligne du corps commence par une majuscule, souvent à la ligne. \textbf{Les titres de civilité et fonctions s'écrivent en minuscule} (\esp{señor}, \esp{señora}, \esp{director}, \esp{directora}).

\begin{center}
\begin{tabularx}{\linewidth}{|X|X|X|}
\hline
\rowcolor{popblueL} Español & Français & Emploi \\
\hline
Estimado señor / Estimada señora & Monsieur / Madame & formule standard, la plus courante \\
\hline
Estimado señor director & Monsieur le Directeur & si l'on connaît la fonction (minuscule à \esp{director}) \\
\hline
Muy señor mío & Monsieur & très solennel, administrations \\
\hline
Distinguido señor & Monsieur (avec égards) & soutenu, relation protocolaire \\
\hline
A quien corresponda & À qui de droit & destinataire inconnu (nom et fonction ignorés) \\
\hline
\end{tabularx}
\end{center}

\begin{exemplebox}[Exemples traduits]
\esp{Estimada señora directora:} « Madame la Directrice, »

\esp{A quien corresponda:} « À qui de droit, » (ex. : demande d'attestation à une administration).

\esp{Muy señor mío:} « Monsieur, » (formule administrative maximale).
\end{exemplebox}

\begin{savaistu}[La formule la plus solennelle]
\esp{Muy señor mío} (littéralement « très mien monsieur ») est la formule la plus solennelle de la correspondance espagnole : elle est réservée aux administrations, aux ambassades, aux courriers officiels. Curieusement, elle se traduit en français administratif par le simple « Monsieur ». À l'inverse, le français « Cher Monsieur » est plus chaleureux que l'espagnol : il se rendra par \esp{Estimado señor}, et non par \esp{Querido señor}, qui sonnerait trop familier.
\end{savaistu}

\begin{attention}[Attention aux majuscules et à « querido »]
En espagnol, les titres de civilité et les fonctions prennent la \cle{minuscule} : \esp{señor}, \esp{señora}, \esp{director} (contrairement au français « Monsieur le Directeur »). Et surtout : \esp{Querido señor} est un calque du français « Cher Monsieur » à éviter dans une lettre formelle ; \esp{querido} est réservé aux lettres amicales (\esp{Querido Pablo}).
\end{attention}

\courssub{Les formules de clôture}

La \esp{despedida} (formule de clôture) termine la lettre avant la signature. Elle est aussi codifiée que la salutation. On distingue les formules \cle{autonomes} (qui ferment la lettre) et les formules \cle{préalables} (qui préparent la clôture).

\begin{center}
\begin{tabularx}{\linewidth}{|X|X|X|}
\hline
\rowcolor{popblueL} Español & Français & Type / Emploi \\
\hline
Le saluda atentamente & Veuillez agréer, Monsieur... (abrégé) & Autonome, standard, très fréquent \\
\hline
Atentamente & Cordialement (formel) & Autonome, simple et sûr \\
\hline
Reciba un cordial saludo & Recevez mes cordiales salutations & Autonome, un peu plus chaleureux \\
\hline
En espera de su respuesta & Dans l'attente de votre réponse & Préalable (se place avant la clôture) \\
\hline
Agradeciéndole de antemano & En vous remerciant d'avance & Préalable (lettre de demande) \\
\hline
Quedo a la espera de sus noticias & Je reste dans l'attente de vos nouvelles & Préalable (moins formel) \\
\hline
\end{tabularx}
\end{center}

\begin{exemplebox}[Une clôture complète]
\esp{En espera de su respuesta, le saluda atentamente.} « Dans l'attente de votre réponse, je vous prie d'agréer, Monsieur, mes salutations distinguées. »

Remarque : l'espagnol préfère une formule courte (\esp{Le saluda atentamente}) là où le français déploie une longue périphrase (« Je vous prie d'agréer, Monsieur, l'expression de mes salutations distinguées »). En version, on traduit donc par la formule française complète ; au thème, on condense.
\end{exemplebox}

\courssub{Le vocabulaire de la demande}

La lettre de demande (\esp{carta de solicitud}) possède son propre lexique. Observons ces phrases tirées de notre texte support : \esp{Solicito información sobre las prácticas} ; \esp{Adjunto mi currículum vítae} ; \esp{Quisiera pedir una beca}.

\begin{center}
\begin{tabularx}{\linewidth}{|X|X|X|}
\hline
\rowcolor{popblueL} Español & Français & Exemple en contexte \\
\hline
solicitar & demander, solliciter (formel) & \esp{Solicito un certificado de estudios.} \\
\hline
una solicitud & une demande (écrite), un formulaire & \esp{Debe enviar una solicitud por escrito.} \\
\hline
adjuntar & joindre (à une lettre) & \esp{Adjunto mi currículum vítae.} \\
\hline
el currículum vítae (CV) & le curriculum vitæ & \esp{Mi currículum vítae tiene dos páginas.} \\
\hline
el expediente & le dossier (administratif) & \esp{Mi expediente está completo.} \\
\hline
la inscripción & l'inscription & \esp{La inscripción termina el 30 de junio.} \\
\hline
la beca & la bourse (d'études) & \esp{Quisiera obtener una beca.} \\
\hline
\end{tabularx}
\end{center}

\begin{attention}[Faux amis et calques fréquents]
\begin{itemize}
\item « Demander » (formel) se traduit par \esp{solicitar} ou \esp{pedir}, \textbf{jamais} par \esp{demandar} : \esp{demandar} signifie « poursuivre en justice, intenter un procès » !
\item \esp{Una solicitud} est la demande écrite, le formulaire officiel ; « une demande » orale se rendra plutôt par \esp{una petición}.
\item « Joindre un document » = \esp{adjuntar un documento} ; on trouve aussi l'abréviation \esp{adj.} en bas de lettre, comme « P.-J. » (Pièce Jointe) en français.
\item \esp{Currículum vítae} s'écrit avec accents (sur le \emph{u} de \emph{currículum} et le \emph{i} de \emph{vítae}) ; l'abréviation \esp{CV} est invariable.
\end{itemize}
\end{attention}

\courssub{Le vocabulaire de la réclamation}

La \esp{carta de reclamación} (lettre de réclamation) sert à protester contre un produit défectueux, un retard, un service insatisfaisant. Son ton reste poli mais ferme.

\begin{center}
\begin{tabularx}{\linewidth}{|X|X|X|}
\hline
\rowcolor{popblueL} Español & Français & Exemple en contexte \\
\hline
una reclamación & une réclamation & \esp{Escribo para presentar una reclamación.} \\
\hline
quejarse de & se plaindre de & \esp{Me quejo del mal servicio.} \\
\hline
un defecto & un défaut & \esp{El producto tiene un defecto.} \\
\hline
un retraso & un retard & \esp{El paquete llegó con un retraso de diez días.} \\
\hline
exigir & exiger & \esp{Exijo una explicación inmediata.} \\
\hline
devolver & rendre (un produit), rembourser (de l'argent) & \esp{Quiero devolver el artículo.} \\
\hline
el reembolso & le remboursement & \esp{Solicito el reembolso total.} \\
\hline
insatisfecho/a & insatisfait(e) & \esp{Estoy insatisfecho con la compra.} \\
\hline
tramitar & traiter (un dossier) & \esp{Les ruego que tramiten mi reclamación.} \\
\hline
\end{tabularx}
\end{center}

\begin{exemplebox}[Exemples traduits]
\esp{Me dirijo a ustedes para quejarme del retraso de mi pedido.} « Je m'adresse à vous pour me plaindre du retard de ma commande. »

\esp{Exijo el reembolso del importe total.} « J'exige le remboursement du montant total. »

\esp{El aparato presenta un defecto de fabricación.} « L'appareil présente un défaut de fabrication. »
\end{exemplebox}

\begin{attention}[« Devolver », verbe à changement de voyelle]
\esp{Devolver} est un verbe à radical changeant (o $\to$ ue) : \esp{devuelvo}, \esp{devuelves}, \esp{devuelve}, \esp{devolvemos}, \esp{devolvéis}, \esp{devuelven} ; participe passé irrégulier : \esp{devuelto}. 
Distinction cruciale : \esp{devolver un producto} = rapporter/retourner un produit (au magasin) ; \esp{devolver el dinero} = rembourser l'argent. 
Enfin, « exiger » = \esp{exigir}, mais « demander poliment » = \esp{solicitar} : dans une réclamation, on commence par \esp{solicito} et l'on ne passe à \esp{exijo} que si l'on veut hausser le ton.
\end{attention}

\courssub{Les verbes de politesse : le cœur de la courtoisie}

Quatre verbes dominent la correspondance formelle. Ils permettent d'atténuer la demande (atténuation = politesse).

\begin{center}
\begin{tabularx}{\linewidth}{|X|X|X|}
\hline
\rowcolor{popblueL} Español & Français & Construction / Exemple \\
\hline
agradecer & remercier & \esp{Le agradezco su atención.} (indicatif présent) \\
\hline
agradecería & je vous serais reconnaissant de & \esp{Le agradecería que me enviara el folleto.} (conditionnel + \cle{subjonctif}) \\
\hline
rogar (ue) & prier de & \esp{Les ruego que me envíen el folleto.} (présent + \cle{subjonctif}) \\
\hline
quisiera & je voudrais, j'aimerais & \esp{Quisiera información sobre los cursos.} (conditionnel + \cle{infinitif}) \\
\hline
\end{tabularx}
\end{center}

\begin{exemplebox}[Exemples traduits]
\esp{Les ruego que me envíen el folleto.} « Je vous prie de bien vouloir m'envoyer la brochure. » (subjonctif \esp{envíen})

\esp{Quisiera saber las fechas de inscripción.} « Je voudrais connaître les dates d'inscription. » (infinitif \esp{saber})

\esp{Le agradezco de antemano su amabilidad.} « Je vous remercie d'avance de votre amabilité. »
\end{exemplebox}

\begin{attention}[Conjugaison et constructions : points de vigilance]
\begin{itemize}
\item \esp{Agradecer} : conjugaison comme \esp{conocer} $\to$ \esp{agradezco} (1ère pers. sg., \emph{z} devant \emph{c}). 
\item \esp{Rogar} : radical changeant (o $\to$ ue) $\to$ \esp{ruego}, \esp{ruegas}, \esp{ruega}... Formule figée : \esp{les ruego que} + \textbf{subjonctif présent} (jamais d'infinitif après \esp{que}).
\item \esp{Quisiera} : conditionnel de politesse de \esp{querer} (imparfait du subjonctif historiquement, figé à la 1ère pers. sg.). Toujours suivi de l'\textbf{infinitif}.
\item \esp{Le agradecería que} : conditionnel de \esp{agradar} (ou \esp{agradecer} selon analyses) + \textbf{subjonctif imparfait} (\esp{enviara} / \esp{enviase}).
\end{itemize}
\end{attention}

\courssub{Méthode : Structurer une lettre formelle}

\begin{methode}[Rédiger une lettre formelle en 5 étapes]
\begin{enumerate}
\item \textbf{En-tête (\esp{Encabezamiento}) :} À gauche : \esp{Remitente} (nom, adresse, téléphone, email). À droite : \esp{Destinatario} (nom, fonction, entreprise, adresse). Date (lieu, jour mois année) à droite sous le destinataire.
\item \textbf{Objet (\esp{Asunto}) :} Ligne claire : \esp{Asunto: solicitud de...} / \esp{reclamación por...}
\item \textbf{Salutation (\esp{Salutación}) :} \esp{Estimado señor director:} (minuscules, deux-points, passage à la ligne).
\item \textbf{Corps (\esp{Cuerpo}) :} 
\begin{itemize}
    \item Phrase d'intro : \esp{Me dirijo a usted para...} / \esp{Le escribo para...}
    \item Exposé des faits / demande précise (vocabulaire spécifique : \esp{solicito}, \esp{adjunto}, \esp{exijo}, \esp{reclamo}).
    \item Formules de politesse : \esp{Quisiera...}, \esp{Les ruego que...}, \esp{Le agradecería que...}
\end{itemize}
\item \textbf{Clôture et signature (\esp{Despedida y firma}) :} Formule préalable (\esp{En espera de su respuesta,}) + Formule autonome (\esp{le saluda atentamente.}) + Signature manuscrite + Nom imprimé.
\end{enumerate}
\end{methode}

\courssub{Carte mentale : la structure lexicale de la lettre formelle}

\begin{popfigure}
\begin{tikzpicture}[mindmap, grow cyclic, every node/.style={concept, align=center, font=\small},
root concept/.append style={concept color=popblue, font=\small\bfseries},
level 1/.append style={level distance=3.4cm, sibling angle=90},
level 2/.append style={level distance=2.4cm, sibling angle=45}]
\node[root concept]{La carta formal}
  child[concept color=poporange]{ node {Encabezamiento}
    child { node[concept color=poporangeL, align=center]{Remitente \\ Destinatario} }
    child { node[concept color=poporangeL, align=center]{Asunto \\ Referencia} }
  }
  child[concept color=popgreen]{ node {Salutación}
    child { node[concept color=popgreenL] {Estimado señor} }
    child { node[concept color=popgreenL] {A quien corresponda} }
  }
  child[concept color=popgold]{ node {Cuerpo}
    child { node[concept color=popgoldL, align=center]{Solicitar \\ Adjuntar} }
    child { node[concept color=popgoldL, align=center]{Quisiera \\ Les ruego que} }
  }
  child[concept color=poppurple]{ node {Despedida}
    child { node[concept color=poppurpleL] {Le saluda atentamente} }
    child { node[concept color=poppurpleL] {En espera de respuesta} }
  };
\end{tikzpicture}
\legende{Carte mentale : les quatre blocs lexicaux de la lettre formelle.}
\end{popfigure}

\courssub{Texte d'application : una carta de reclamación}

Lisons maintenant une lettre de réclamation qui réutilise presque tout le lexique de cette section.

\begin{textebox}[Una carta de reclamación]
Yaoundé, 15 de mayo

Estimado señor director de la empresa TransCamer:

Asunto: reclamación por un retraso

Me dirijo a usted para presentar una reclamación. El pasado 3 de mayo compré en su agencia de Douala un billete de autobús Douala–Yaoundé para el día 10 de mayo a las ocho de la mañana. Sin embargo, el autobús salió con un retraso de tres horas y llegó a Yaoundé a las cuatro de la tarde. A causa de este retraso, perdí una cita muy importante con el director de un instituto donde quiero hacer la inscripción.

Estoy muy insatisfecho con su servicio. Los pasajeros no recibieron ninguna explicación ni ninguna disculpa. Además, el autobús presentaba un defecto en el aire acondicionado y el viaje fue muy incómodo.

Por todo esto, exijo el reembolso total del billete. Adjunto una fotocopia del billete y del recibo de compra. Les ruego que tramiten mi reclamación lo antes posible.

Quisiera recibir una respuesta antes del 30 de mayo. En espera de su respuesta, le saluda atentamente.

André Mbarga
\end{textebox}

\textbf{Glossaire.} \esp{el billete} : le billet ; \esp{sin embargo} : cependant ; \esp{perdí una cita} : j'ai manqué un rendez-vous ; \esp{la disculpa} : l'excuse ; \esp{el aire acondicionado} : la climatisation ; \esp{tramitar} : traiter (un dossier) ; \esp{lo antes posible} : le plus tôt possible ; \esp{el recibo} : le reçu ; \esp{la fotocopia} : la photocopie.

\textbf{Questions de compréhension.}
\begin{enumerate}
\item ¿Quién es el remitente y quién es el destinatario de esta carta ?
\item ¿Cuál es el asunto de la carta ?
\item ¿Qué dos problemas menciona André (retard et confort) ?
\item ¿Qué exige el remitente y qué documentos adjunta ?
\item Identifica dos fórmulas de cortesía de la carta (une au conditionnel, une au subjonctif).
\end{enumerate}

\begin{exoresolu}[Compléter une lettre formelle avec le lexique étudié]
Énoncé. Complétez les espaces avec les mots suivants : \esp{asunto}, \esp{adjunto}, \esp{reclamación}, \esp{Estimado}, \esp{atentamente}, \esp{reembolso}.

\esp{\ldots\ldots\ldots señor director:}

\esp{\ldots\ldots\ldots: reclamación por un defecto del producto}

\esp{Le escribo para presentar una \ldots\ldots\ldots. El teléfono que compré la semana pasada tiene un defecto de pantalla. \ldots\ldots\ldots el recibo de compra. Solicito el \ldots\ldots\ldots total. Le saluda \ldots\ldots\ldots.}

\tcblower
Solution détaillée.

1. \esp{Estimado señor director:} — salutation formelle standard (minuscules aux titres), suivie de deux-points.

2. \esp{Asunto: reclamación por un defecto del producto} — l'objet de la lettre s'introduit par \esp{asunto} (jamais \esp{objeto}).

3. \esp{...para presentar una reclamación.} — expression figée : \esp{presentar una reclamación} = déposer une réclamation.

4. \esp{Adjunto el recibo de compra.} — \esp{adjuntar} = joindre ; première personne du présent : \esp{adjunto}.

5. \esp{...solicito el reembolso total.} — le remboursement = \esp{el reembolso} (verbe \esp{solicitar} au présent).

6. \esp{Le saluda atentamente.} — formule de clôture classique, équivalent condensé de « Je vous prie d'agréer, Monsieur, mes salutations distinguées ».
\end{exoresolu}

\begin{aretenir}[L'essentiel du lexique de la correspondance formelle]
\begin{itemize}
\item En-tête : \esp{remitente}, \esp{destinatario}, \esp{dirección}, \esp{asunto} (= objet, pas « sujet »).
\item Salutations : \esp{Estimado señor}, \esp{Muy señor mío}, \esp{A quien corresponda} (« À qui de droit »), toujours suivies de deux-points ; titres en minuscules.
\item Clôtures : \esp{Le saluda atentamente}, \esp{Atentamente}, \esp{En espera de su respuesta} (préalable), \esp{Agradeciéndole de antemano} (préalable).
\item Demande : \esp{solicitar}, \esp{una solicitud}, \esp{adjuntar}, \esp{el currículum vítae}, \esp{la beca}, \esp{la inscripción}.
\item Réclamation : \esp{quejarse de}, \esp{un defecto}, \esp{un retraso}, \esp{exigir}, \esp{el reembolso}, \esp{insatisfecho}, \esp{tramitar}.
\item Politesse (grammaire) : \esp{quisiera} + infinitif ; \esp{les ruego que} / \esp{le agradecería que} + subjonctif.
\end{itemize}
\end{aretenir}

Ce lexique désormais maîtrisé, nous pouvons passer à la section suivante, consacrée à la grammaire qui rend ces formules possibles : le conditionnel de courtoisie (\esp{quisiera}, \esp{le agradecería que}), dont nous étudierons la formation et les emplois en détail.

\coursec{Grammaire : le conditionnel de courtoisie}

Dans la section précédente, nous avons découvert le lexique de la correspondance formelle : \esp{estimado señor}, \esp{atentamente}, \esp{a la atención de}, \esp{en relación con su anuncio}... Mais un vocabulaire correct ne suffit pas pour écrire une bonne lettre administrative. Il faut aussi un \cle{ton} adapté : en espagnol comme en français, on ne s'adresse pas à un directeur d'entreprise comme à un camarade de classe. L'outil grammatical essentiel de la politesse écrite est le \cle{conditionnel de courtoisie}. C'est lui que nous allons étudier maintenant, à partir des phrases de la lettre de demande de stage lue au début de la leçon.

\courssub{Observation : repérons le conditionnel dans la lettre}

Relisons deux phrases extraites de la \esp{carta de solicitud de prácticas} étudiée en section 1 :

\begin{exemplebox}[Extraits de la lettre de demande de stage]
\esp{Quisiera adquirir experiencia profesional en su empresa.}

« Je voudrais acquérir une expérience professionnelle dans votre entreprise. »

\esp{Le agradecería que me enviara información sobre las fechas de las prácticas.}

« Je vous serais reconnaissant de bien vouloir m'envoyer des informations sur les dates du stage. »
\end{exemplebox}

Observons attentivement :

\begin{itemize}
\item \esp{quisiera} et \esp{agradecería} se terminent tous les deux en \textbf{-ía} : ce sont des formes du \cle{conditionnel présent}.
\item L'étudiant n'écrit pas \esp{quiero adquirir experiencia} (« je veux acquérir... »), qui sonnerait sec et exigeant, mais \esp{quisiera} : le conditionnel \textbf{adoucit} la demande.
\item Après \esp{le agradecería que}, le verbe de la subordonnée est \esp{enviara}, une forme particulière : l'\cle{imparfait du subjonctif}. Ce n'est pas un hasard, c'est une règle obligatoire.
\end{itemize}

\begin{definition}[Le conditionnel de courtoisie]
Le \cle{conditionnel de courtoisie} est l'emploi du conditionnel présent espagnol pour \textbf{atténuer} une demande, exprimer un souhait poli, poser une question délicate ou donner un conseil sans imposer. Il correspond exactement au conditionnel de politesse français : « je \emph{voudrais} », « pourriez-vous... ? », « vous \emph{devriez}... ».
\end{definition}

\courssub{Formation du conditionnel présent : rappel}

\begin{propriete}[Formation du conditionnel présent]
Le conditionnel présent se forme sur l'\textbf{infinitif entier} du verbe (sans enlever la terminaison), auquel on ajoute les terminaisons de l'imparfait de l'indicatif des verbes en \textbf{-ir / -er} :

\begin{center}
\textbf{-ía, -ías, -ía, -íamos, -íais, -ían}
\end{center}

Exemple avec \esp{hablar} : \esp{hablaría, hablarías, hablaría, hablaríamos, hablaríais, hablarían} (« je parlerais, tu parlerais... »).
\end{propriete}

\begin{attention}[Accent obligatoire]
Toutes les formes du conditionnel portent un \textbf{accent écrit sur le í} : \esp{hablar\textbf{í}a}, \esp{comer\textbf{í}amos}. Sans accent, la forme n'existe pas en espagnol. Erreur fréquente : écrire \esp{hablaria} sous l'influence du français « parlerais ».
\end{attention}

\textbf{Les verbes irréguliers}\par

Comme au futur simple, certains verbes modifient leur radical avant d'ajouter les terminaisons. Ce sont exactement les mêmes irrégularités qu'au futur : si vous connaissez le futur, vous connaissez le conditionnel.

\begin{center}
\begin{tabularx}{\linewidth}{|X|X|X|X|}
\hline
\rowcolor{popblueL} \textbf{Infinitif} & \textbf{Radical} & \textbf{1\iere{} personne} & \textbf{Français} \\
\hline
querer (vouloir) & querr- / quisier- & querría / quisiera & je voudrais \\
\hline
poder (pouvoir) & podr- & podría & je pourrais \\
\hline
hacer (faire) & har- & haría & je ferais \\
\hline
decir (dire) & dir- & diría & je dirais \\
\hline
tener (avoir) & tendr- & tendría & j'aurais \\
\hline
salir (sortir) & saldr- & saldría & je sortirais \\
\hline
poner (mettre) & pondr- & pondría & je mettrais \\
\hline
saber (savoir) & sabr- & sabría & je saurais \\
\hline
haber (avoir, auxiliaire) & habr- & habría & il y aurait \\
\hline
\end{tabularx}
\end{center}

\begin{attention}[Cas particulier : \esp{quisiera}]
Le verbe \esp{querer} possède \textbf{deux} conditionnels : le conditionnel régulier \esp{querría} et la forme héritée de l'imparfait du subjonctif \esp{quisiera}. Dans la correspondance formelle, on emploie presque toujours \esp{quisiera}, plus élégant et plus poli : \esp{Quisiera información} (« je voudrais des renseignements »). \esp{Querría} existe mais sonne plus littéraire.
\end{attention}

\courssub{Tableau de conjugaison des verbes clés de la correspondance}

Voici les quatre verbes indispensables de la lettre formelle, conjugués au conditionnel présent :

\begin{popfigure}
\begin{tikzpicture}[
  box/.style={draw=popblue, fill=popblueL, rounded corners=2pt, minimum width=2.6cm, minimum height=0.55cm, font=\small, align=center},
  head/.style={draw=popdark, fill=popdark, text=white, rounded corners=2pt, minimum width=2.6cm, minimum height=0.6cm, font=\small\bfseries, align=center},
  pers/.style={font=\small\itshape, text=popmuted, align=right}
]
% En-têtes
\node[head] (h0) at (0,0) {};
\node[head] (h1) at (3.0,0) {querer};
\node[head] (h2) at (6.0,0) {poder};
\node[head] (h3) at (9.0,0) {agradecer};
\node[head] (h4) at (12.0,0) {gustar};
% Lignes
\node[pers] at (-0.2,-0.8) {yo};
\node[box] at (3.0,-0.8) {quisiera / querría};
\node[box] at (6.0,-0.8) {podría};
\node[box] at (9.0,-0.8) {agradecería};
\node[box] at (12.0,-0.8) {me gustaría};
\node[pers] at (-0.2,-1.6) {tú};
\node[box] at (3.0,-1.6) {quisieras};
\node[box] at (6.0,-1.6) {podrías};
\node[box] at (9.0,-1.6) {agradecerías};
\node[box] at (12.0,-1.6) {te gustaría};
\node[pers] at (-0.2,-2.4) {él / ella / usted};
\node[box] at (3.0,-2.4) {quisiera};
\node[box] at (6.0,-2.4) {podría};
\node[box] at (9.0,-2.4) {agradecería};
\node[box] at (12.0,-2.4) {le gustaría};
\node[pers] at (-0.2,-3.2) {nosotros};
\node[box] at (3.0,-3.2) {quisiéramos};
\node[box] at (6.0,-3.2) {podríamos};
\node[box] at (9.0,-3.2) {agradeceríamos};
\node[box] at (12.0,-3.2) {nos gustaría};
\node[pers] at (-0.2,-4.0) {vosotros};
\node[box] at (3.0,-4.0) {quisierais};
\node[box] at (6.0,-4.0) {podríais};
\node[box] at (9.0,-4.0) {agradeceríais};
\node[box] at (12.0,-4.0) {os gustaría};
\node[pers] at (-0.2,-4.8) {ellos / ustedes};
\node[box] at (3.0,-4.8) {quisieran};
\node[box] at (6.0,-4.8) {podrían};
\node[box] at (9.0,-4.8) {agradecerían};
\node[box] at (12.0,-4.8) {les gustaría};
\end{tikzpicture}
\legende{Conditionnel présent de quatre verbes de la correspondance formelle. Attention : \esp{gustar} se construit avec un pronom complément (me, te, le...) et se conjugue presque toujours à la 3\ieme{} personne.}
\end{popfigure}

\begin{aretenir}[Les trois formes à mémoriser en priorité]
\begin{itemize}
\item \esp{Quisiera} + infinitif ou nom : « je voudrais... »
\item \esp{Me gustaría} + infinitif : « j'aimerais... »
\item \esp{Le agradecería que} + imparfait du subjonctif : « je vous serais reconnaissant de bien vouloir... »
\end{itemize}
Ces trois formules suffisent pour rédiger la quasi-totalité des demandes d'une lettre formelle de niveau Première.
\end{aretenir}

\courssub{Emploi 1 : atténuer une demande}

C'est l'emploi fondamental. Le présent de l'indicatif exprime un fait, donc une exigence ; le conditionnel exprime un souhait, donc une demande respectueuse.

\begin{exemplebox}[Du direct au poli]
\esp{Quiero información sobre el curso.} « Je veux des informations sur le cours. » — correct mais \textbf{trop direct}, à éviter dans une lettre.

\esp{Quisiera información sobre el curso.} « Je voudrais des informations sur le cours. » — poli, adapté.

\esp{Quisiera solicitar una plaza en su academia.} « Je voudrais solliciter une place dans votre académie. »

\esp{Quisiera pedir una entrevista con el director.} « Je voudrais demander un entretien avec le directeur. »
\end{exemplebox}

Pour les \textbf{questions polies}, on utilise \esp{¿Podría usted...?} :

\begin{exemplebox}[Questions polies]
\esp{¿Podría usted indicarme el horario?} « Pourriez-vous m'indiquer les horaires ? »

\esp{¿Podría usted enviarme el formulario de inscripción?} « Pourriez-vous m'envoyer le formulaire d'inscription ? »

\esp{¿Sería posible visitar sus instalaciones?} « Serait-il possible de visiter vos installations ? »
\end{exemplebox}

\courssub{Emploi 2 : exprimer un souhait poli}

Pour parler de ses désirs professionnels ou personnels sans paraître présomptueux, on emploie \esp{me gustaría} + infinitif ou le conditionnel d'un verbe de volonté.

\begin{exemplebox}[Souhaits polis]
\esp{Me gustaría trabajar en su empresa.} « J'aimerais travailler dans votre entreprise. »

\esp{Me gustaría inscribirme en el curso de verano.} « J'aimerais m'inscrire au cours d'été. »

\esp{Me gustaría mejorar mi nivel de español en Guinea Ecuatorial.} « J'aimerais améliorer mon niveau d'espagnol en Guinée équatoriale. »

\esp{Haría las prácticas con mucho entusiasmo.} « Je ferais le stage avec beaucoup d'enthousiasme. »
\end{exemplebox}

\courssub{Emploi 3 : la structure « le agradecería que » + imparfait du subjonctif}

C'est la formule la plus formelle, typique du style administratif. Elle exige une \cle{concordance des temps} stricte.

\begin{propriete}[Après « le agradecería que »]
Après \esp{le agradecería que}, le verbe de la subordonnée est obligatoirement à l'\textbf{imparfait du subjonctif} (terminaison \textbf{-ara / -iera}) :

\esp{Le agradecería que me \textbf{enviara} el expediente.}

« Je vous serais reconnaissant de bien vouloir m'envoyer le dossier. »

Rappel de formation de l'imparfait du subjonctif : 3\ieme{} personne du pluriel du prétérit \textbf{moins -ron + -ra} : \esp{enviaron} → \esp{enviara} ; \esp{pusieron} → \esp{pusiera} ; \esp{tuvieron} → \esp{tuviera}.
\end{propriete}

\begin{exemplebox}[La formule administrative en action]
\esp{Le agradecería que me enviara el formulario antes del viernes.} « Je vous serais reconnaissant de bien vouloir m'envoyer le formulaire avant vendredi. »

\esp{Le agradecería que tuviera en cuenta mi candidatura.} « Je vous serais reconnaissant de bien vouloir prendre en compte ma candidature. »

\esp{Les agradeceríamos que nos confirmaran la fecha de la entrevista.} « Nous vous serions reconnaissants de bien vouloir nous confirmer la date de l'entretien. »
\end{exemplebox}

\begin{attention}[L'erreur classique du francophone]
En français, on dit « je vous serais reconnaissant de bien vouloir m'\emph{envoyer} » : un infinitif. Le francophone traduit donc souvent par l'indicatif présent ou l'infinitif :

\esp{Le agradecería que me \textbf{envía} el expediente.} FAUX (indicatif)

\esp{Le agradecería que me \textbf{enviara} el expediente.} CORRECT (subjonctif imparfait)

Même vigilance avec \esp{quisiera que} : on écrit \esp{Quisiera que me enviara información}, jamais \esp{quisiera que me envía información}. Règle d'or : \textbf{verbe principal au conditionnel + que → subordonnée au subjonctif imparfait}.
\end{attention}

\courssub{Emploi 4 : conseils et suggestions polis}

Le conditionnel de \esp{deber} (\esp{debería}, « vous devriez ») permet de conseiller sans ordonner. C'est l'équivalent exact du français.

\begin{exemplebox}[Conseils polis]
\esp{Debería revisar el documento antes de enviarlo.} « Vous devriez relire le document avant de l'envoyer. »

\esp{Deberías escribir la carta en papel blanco.} « Tu devrais écrire la lettre sur du papier blanc. »

\esp{Yo que usted, adjuntaría una fotocopia del diploma.} « Si j'étais vous, je joindrais une photocopie du diplôme. »
\end{exemplebox}

\courssub{Comparaison français / espagnol}

\begin{center}
\begin{tabularx}{\linewidth}{|X|X|X|}
\hline
\rowcolor{popblueL} \textbf{Français} & \textbf{Español} & \textbf{Remarque} \\
\hline
Je veux un renseignement. & Quiero una información. & Direct, à éviter en lettre formelle dans les deux langues. \\
\hline
Je voudrais un renseignement. & Quisiera una información. & Conditionnel de politesse : identique dans les deux langues. \\
\hline
Pourriez-vous m'aider ? & ¿Podría usted ayudarme? & Le pronom objet se place après l'infinitif en espagnol. \\
\hline
J'aimerais m'inscrire. & Me gustaría inscribirme. & \esp{gustar} + pronom me : construction inversée. \\
\hline
Je vous serais reconnaissant de bien vouloir m'envoyer le dossier. & Le agradecería que me enviara el expediente. & Différence majeure : infinitif en français, subjonctif imparfait en espagnol. \\
\hline
Vous devriez relire la lettre. & Debería revisar la carta. & Parallélisme parfait. \\
\hline
\end{tabularx}
\end{center}

\begin{methode}[Choisir le bon degré de politesse]
Pour graduer une demande dans une lettre formelle, suivez l'échelle de politesse :

\begin{enumerate}
\item \textbf{Niveau direct} (à éviter avec un supérieur) : \esp{Quiero información.} « Je veux des informations. »
\item \textbf{Niveau poli} (courant, sûr) : \esp{Me gustaría recibir información.} « J'aimerais recevoir des informations. »
\item \textbf{Niveau très poli} (lettre formelle) : \esp{Quisiera recibir información.} « Je voudrais recevoir des informations. »
\item \textbf{Niveau administratif} (demande importante, réclamation, autorité) : \esp{Le agradecería que me enviara la información.} « Je vous serais reconnaissant de bien vouloir m'envoyer les informations. »
\end{enumerate}

Conseil pratique : dans une lettre de demande ou de motivation à l'examen, utilisez au moins une formule de niveau 3 et une de niveau 4 — c'est ce que le correcteur attend.
\end{methode}

\begin{exoresolu}[Mettez les demandes au bon degré de politesse]
Énoncé — Un élève de Première A de Yaoundé écrit au consulat d'Espagne. Transformez ces demandes trop directes en demandes polies, en utilisant le verbe entre parenthèses au conditionnel.

1. \esp{Quiero un formulario de solicitud de visado.} (querer, très poli)

2. \esp{¿Puede usted decirme el precio del visado?} (poder, question polie)

3. \esp{Envíenme la lista de documentos necesarios.} (agradecer, niveau administratif)

4. \esp{Quiero trabajar como traductor en su embajada.} (gustar, souhait poli)
\tcblower
Solution détaillée :

1. \esp{Quisiera un formulario de solicitud de visado.} — \esp{querer} au conditionnel de courtoisie : \esp{quisiera} (forme préférée à \esp{querría} dans la correspondance). « Je voudrais un formulaire de demande de visa. »

2. \esp{¿Podría usted decirme el precio del visado?} — \esp{poder} → radical irrégulier \esp{podr-} + \esp{-ía} ; le pronom \esp{me} se soude à l'infinitif \esp{decirme}. « Pourriez-vous me dire le prix du visa ? »

3. \esp{Les agradecería que me enviaran la lista de documentos necesarios.} — après \esp{agradecería que}, subjonctif imparfait : \esp{enviaron} → \esp{enviaran} (3\ieme{} personne du pluriel, car on s'adresse à \esp{ustedes}). « Je vous serais reconnaissant de bien vouloir m'envoyer la liste des documents nécessaires. »

4. \esp{Me gustaría trabajar como traductor en su embajada.} — \esp{gustar} se construit avec le pronom \esp{me} et reste à la 3\ieme{} personne du singulier. « J'aimerais travailler comme traducteur dans votre ambassade. »
\end{exoresolu}

\begin{attention}[Récapitulatif des pièges]
\begin{itemize}
\item Ne pas oublier l'accent : \esp{gustaría}, \esp{podría}, \esp{debería} — jamais \esp{gustaria}.
\item Ne pas confondre \esp{habría} (« il y aurait », verbe \esp{haber}) et \esp{haría} (« je ferais », verbe \esp{hacer}).
\item Après \esp{quisiera que} et \esp{le agradecería que} : toujours le subjonctif imparfait (\esp{enviara}, \esp{tuviera}, \esp{fuera}), jamais l'indicatif présent.
\item Ne pas calquer « je vous serais reconnaissant de + infinitif » : l'espagnol exige \esp{que} + subjonctif.
\item \esp{¿Podría usted...?} garde le \esp{usted} explicite ou implicite, mais le verbe reste à la 3\ieme{} personne : \esp{¿Podría usted ayudarme?}, non \esp{¿Podrías usted...?}.
\end{itemize}
\end{attention}

\begin{aretenir}[L'essentiel du conditionnel de courtoisie]
Le conditionnel présent (infinitif + \textbf{-ía, -ías, -ía, -íamos, -íais, -ían}) transforme une exigence en demande respectueuse. En lettre formelle, retenez la progression : \esp{quiero} (trop direct) → \esp{me gustaría} (poli) → \esp{quisiera} (très poli) → \esp{le agradecería que} + subjonctif imparfait (administratif). C'est cette maîtrise du ton qui distinguera votre lettre de motivation dans la section 6.
\end{aretenir}

\coursec{Traduction des formules de correspondance}

Dans la section précédente, nous avons appris à former et à employer le conditionnel de courtoisie (\esp{quisiera}, \esp{le agradecería}) pour adoucir une demande. Mais une lettre formelle ne se résume pas à quelques verbes bien conjugués : elle repose sur des \cle{formules codifiées}, c'est-à-dire des expressions figées que l'on retrouve dans toutes les lettres administratives. Or ces formules ne se traduisent presque jamais mot à mot entre le français et l'espagnol. C'est précisément l'objet de cette section : apprendre à traduire correctement les formules d'ouverture, de demande et de clôture, compétence directement évaluée au baccalauréat dans les épreuves de thème et de version.

\courssub{Observation : deux lettres, deux langues, une même fonction}

Observons d'abord ces deux extraits de lettres formelles, l'une en français, l'autre en espagnol. Ils disent exactement la même chose, mais pas avec les mêmes mots.

\begin{exemplebox}[Deux ouvertures de lettre comparées]
\textbf{Français :} « Monsieur le Directeur, je me permets de vous écrire afin de solliciter un stage au sein de votre entreprise. Suite à votre annonce parue dans le journal, je vous serais reconnaissant de bien vouloir examiner ma candidature. Dans l'attente de votre réponse, veuillez agréer, Monsieur le Directeur, l'expression de mes salutations distinguées. »

\textbf{Español :} \esp{¿Muy señor mío? Me permito dirigirme a usted para solicitar unas prácticas en su empresa. En relación con su anuncio publicado en el periódico, le agradecería que examinara mi candidatura. En espera de su respuesta, le saluda atentamente.}
\end{exemplebox}

Que remarque-t-on ? La lettre française est longue, cérémonieuse, presque solennelle. La lettre espagnole est plus \cle{directe} et plus \cle{concise}. Chaque formule française possède un \cle{équivalent codifié} en espagnol, mais cet équivalent n'est pas une traduction littérale : « je me permets de vous écrire » ne devient pas \esp{me permito escribirle} (calque syntaxique maladroit), mais \esp{me permito dirigirme a usted}.

\begin{definition}[Formule codifiée]
Une \cle{formule codifiée} (ou formule figée) est une expression fixe, conventionnelle, utilisée dans un type de texte précis. En correspondance administrative, on ne la traduit pas mot à mot : on la remplace par la formule qui joue le \emph{même rôle} dans l'autre langue. On parle d'\cle{équivalence fonctionnelle}.
\end{definition}

\courssub{Les grandes familles de formules et leurs équivalents}

Classons les formules selon leur \cle{fonction} dans la lettre : ouverture, introduction du sujet, demande, clôture.

\textbf{1. Formules d'ouverture et d'introduction du sujet}

\begin{center}
\begin{tabularx}{\linewidth}{|X|X|X|}
\hline
\rowcolor{popblueL} Français & Español & Remarque \\
\hline
Je me permets de vous écrire & Me permito dirigirme a usted & Jamais \esp{escribirle} seul ici \\
\hline
J'ai l'honneur de... & Tengo el honor de... & Très solennel, registre officiel \\
\hline
Suite à votre annonce... & En relación con su anuncio... & Aussi : \esp{Con referencia a...} \\
\hline
En réponse à votre lettre du 5 mai & En respuesta a su carta del 5 de mayo & Date : \esp{del + jour + de + mois} \\
\hline
J'accuse réception de votre lettre & Acuso recibo de su carta & Formule administrative type \\
\hline
\end{tabularx}
\end{center}

\textbf{2. Formules de demande}

\begin{center}
\begin{tabularx}{\linewidth}{|X|X|X|}
\hline
\rowcolor{popgreenL} Français & Español & Remarque \\
\hline
Je vous serais reconnaissant de... & Le agradecería que + subj. imparfait & Conditionnel + subjonctif imparfait \\
\hline
Je vous prie de bien vouloir... & Les ruego que + subj. présent & \esp{Les ruego} (présent) + \esp{que} + subj. présent \\
\hline
Je souhaiterais obtenir... & Quisiera obtener información sobre... & Conditionnel de courtoisie (\esp{quisiera}) \\
\hline
Veuillez trouver ci-joint... & Adjunto encontrará... / Les adjunto... & \esp{Adjunto}, jamais \esp{junto} \\
\hline
\end{tabularx}
\end{center}

\textbf{3. Formules de clôture}

\begin{center}
\begin{tabularx}{\linewidth}{|X|X|X|}
\hline
\rowcolor{poporangeL} Français & Español & Remarque \\
\hline
Dans l'attente de votre réponse & En espera de su respuesta & Souvent suivi de \esp{le saluda atentamente} \\
\hline
Veuillez agréer, Monsieur, l'expression de mes salutations distinguées & Le saluda atentamente / Reciba un cordial saludo & L'espagnol préfère la concision \\
\hline
Cordialement & Atentamente / Un saludo & Éviter \esp{cordialmente} au registre formel \\
\hline
Je reste à votre disposition & Quedo a su disposición & Formule de disponibilité \\
\hline
\end{tabularx}
\end{center}

La figure suivante résume visuellement ces correspondances essentielles.

\begin{popfigure}
\begin{tikzpicture}[
  fr/.style={rectangle, draw=popblue, fill=popblueL, rounded corners=2pt, text width=5.6cm, align=left, font=\small, inner sep=3pt},
  es/.style={rectangle, draw=popgreen, fill=popgreenL, rounded corners=2pt, text width=5.6cm, align=left, font=\small, inner sep=3pt},
  arr/.style={-{Stealth[length=2mm]}, thick, poporange}]
\node[font=\bfseries\small, text=popdark] at (2.9,0.4) {Français};
\node[font=\bfseries\small, text=popdark] at (9.6,0.4) {Español};
\foreach \y/\f/\e in {
  0/{Je me permets de vous écrire}/{Me permito dirigirme a usted},
  -0.95/{J'ai l'honneur de...}/{Tengo el honor de...},
  -1.9/{Suite à votre annonce...}/{En relación con su anuncio...},
  -2.85/{J'accuse réception de votre lettre}/{Acuso recibo de su carta},
  -3.8/{Je vous serais reconnaissant de...}/{Le agradecería que + subj. imparfait},
  -4.75/{Je vous prie de bien vouloir...}/{Les ruego que + subj. présent},
  -5.7/{Veuillez trouver ci-joint...}/{Adjunto encontrará...},
  -6.65/{Dans l'attente de votre réponse}/{En espera de su respuesta},
  -7.6/{Je reste à votre disposition}/{Quedo a su disposición},
  -8.55/{Veuillez agréer... mes salutations distinguées}/{Le saluda atentamente}}
{
  \node[fr] (f) at (2.9,\y) {\f};
  \node[es] (e) at (9.6,\y) {\e};
  \draw[arr] (f.east) -- (e.west);
}
\end{tikzpicture}
\legende{Dix formules françaises et leurs équivalents espagnols codifiés}
\end{popfigure}

\courssub{Exemples traduits et commentés}

\begin{exemplebox}[Formules en contexte]
\begin{itemize}
\item \esp{¿En relación con su anuncio publicado en El País, me permito dirigirme a usted?} « Suite à votre annonce parue dans \emph{El País}, je me permets de vous écrire. »
\item \esp{¿Acuso recibo de su carta del 12 de marzo?} « J'accuse réception de votre lettre du 12 mars. »
\item \esp{¿Le agradecería que me enviara la información solicitada?} « Je vous serais reconnaissant de bien vouloir m'envoyer les informations demandées. »
\item \esp{¿Les ruego que acepten mi solicitud de inscripción?} « Je vous prie de bien vouloir accepter ma demande d'inscription. »
\item \esp{¿Adjunto encontrará mi currículum vítae y una fotocopia de mi diploma?} « Veuillez trouver ci-joint mon curriculum vitae et une photocopie de mon diplôme. »
\item \esp{¿En espera de su respuesta, le saluda atentamente?} « Dans l'attente de votre réponse, je vous prie d'agréer mes salutations distinguées. »
\item \esp{¿Quedo a su disposición para cualquier información adicional?} « Je reste à votre disposition pour tout renseignement complémentaire. »
\end{itemize}
\end{exemplebox}

Remarquez la construction de \esp{adjunto encontrará} : littéralement « ci-joint vous trouverez », avec inversion du verbe. C'est la formule espagnole standard ; on peut aussi dire \esp{les adjunto mi currículum} (« je vous joins mon CV »).

\courssub{Méthode : traduire une formule de correspondance en trois étapes}

\begin{methode}[Traduire une formule sans calquer]
\begin{enumerate}
\item \textbf{Identifier la fonction} de la formule dans la lettre : est-ce une ouverture, une introduction du sujet, une demande, une pièce jointe, une clôture ?
\item \textbf{Choisir l'équivalent codifié} dans la langue d'arrivée, en puisant dans le répertoire des formules apprises (tableaux ci-dessus), et non en traduisant mot à mot.
\item \textbf{Vérifier le registre et la grammaire} : lettre à un ministre, à un directeur d'entreprise ou à un ami ? En espagnol formel, on privilégie \esp{atentamente}, \esp{le saluda atentamente}, \esp{reciba un cordial saludo} ; on évite les formules trop longues ou trop familières. Vérifier la concordance des temps : \esp{le agradecería que} (conditionnel) exige le \cle{subjonctif imparfait} (\esp{enviara}, \esp{aceptara}) ; \esp{les ruego que} (présent de l'indicatif) exige le \cle{subjonctif présent} (\esp{envíen}, \esp{acepten}).
\end{enumerate}
\end{methode}

\begin{attention}[Le piège du mot à mot : la clôture à la française]
La formule française « Veuillez agréer, Monsieur, l'expression de mes sentiments distingués » ne se traduit \textbf{jamais} mot à mot. Une tentative comme \esp{le ruego que acepte la expresión de mis sentimientos distinguidos} est un calque ridicule et incompréhensible pour un hispanophone. L'espagnol préfère des formules \textbf{courtes} : \esp{Atentamente}, \esp{Le saluda atentamente}, \esp{Reciba un cordial saludo}. Autres pièges : « ci-joint » se dit \esp{adjunto} (et non \esp{junto}, qui signifie « ensemble ») ; « accusé de réception » (nom) se dit \esp{acuse de recibo} ; « cordialement » ne se traduit pas par \esp{cordialmente}, jugé trop familier ou ambigu dans plusieurs pays hispanophones au registre administratif : préfère \esp{atentamente}.
\end{attention}

\begin{exoresolu}[Thème : traduire une formule de demande]
Énoncé : traduis en espagnol la phrase suivante : « Je vous serais reconnaissant de bien vouloir m'envoyer un formulaire d'inscription. »
\tcblower
Solution détaillée. Étape 1 : la fonction est une \textbf{demande polie}. Étape 2 : l'équivalent codifié de « je vous serais reconnaissant de » est \esp{le agradecería que}, qui exige le subjonctif imparfait. « M'envoyer » devient donc \esp{me enviara} (ou \esp{me enviase}). « Un formulaire d'inscription » se dit \esp{un formulario de inscripción}. Étape 3 : vérification du registre — formel, avec \esp{le} (usted). Traduction finale : \esp{Le agradecería que me enviara un formulario de inscripción.} Erreurs à éviter : \esp{agradecería de enviarme} (calque de l'infinitif français) et \esp{estaría agradecido} (grammaticalement possible mais non codifié dans ce contexte).
\end{exoresolu}

\begin{savaistu}[La formule de politesse au baccalauréat]
Au baccalauréat camerounais, l'épreuve d'espagnol évalue à la fois la \cle{compétence linguistique} (grammaire, vocabulaire) et la \cle{compétence socioculturelle} (savoir se comporter dans la culture hispanophone). Une formule de politesse mal traduite — par exemple une clôture calquée du français — fait perdre des points sur les \emph{deux} compétences à la fois. À l'inverse, une formule comme \esp{En espera de su respuesta, le saluda atentamente}, correctement placée, montre au correcteur que tu maîtrises les usages réels de la correspondance hispanophone. En Guinée équatoriale, voisin hispanophone du Cameroun, les administrations utilisent quotidiennement ces formules : les connaître, c'est aussi se préparer à communiquer avec le monde hispanophone le plus proche.
\end{savaistu}

\courssub{Exercices d'entraînement}

\begin{exercice}[Thème : du français vers l'espagnol]
Traduis en espagnol les six formules suivantes.
\begin{enumerate}
\item Je me permets de vous écrire au sujet de votre annonce.
\item J'ai l'honneur de solliciter un stage dans votre entreprise.
\item Je vous prie de bien vouloir examiner ma candidature.
\item Veuillez trouver ci-joint mon curriculum vitae.
\item Dans l'attente de votre réponse, je reste à votre disposition.
\item Veuillez agréer, Madame, l'expression de mes salutations distinguées.
\end{enumerate}
\end{exercice}

\begin{corrige}[Thème : du français vers l'espagnol]
\begin{enumerate}
\item \esp{Me permito dirigirme a usted en relación con su anuncio.}
\item \esp{Tengo el honor de solicitar unas prácticas en su empresa.}
\item \esp{Les ruego que examinen mi candidatura.} (subjonctif présent car \esp{les ruego} est au présent de l'indicatif)
\item \esp{Adjunto encontrará mi currículum vítae.}
\item \esp{En espera de su respuesta, quedo a su disposición.}
\item \esp{Reciba, señora, un cordial saludo.} ou \esp{Le saluda atentamente.}
\end{enumerate}
\end{corrige}

\begin{exercice}[Version : de l'espagnol vers le français administratif]
Traduis en français administratif correct les six formules suivantes.
\begin{enumerate}
\item \esp{Acuso recibo de su carta del 3 de junio.}
\item \esp{En relación con su anuncio publicado en el periódico...}
\item \esp{Le agradecería que me enviara más información.}
\item \esp{Les ruego que acepten mi solicitud.}
\item \esp{Adjunto encontrará los documentos necesarios.}
\item \esp{En espera de su respuesta, le saluda atentamente.}
\end{enumerate}
\end{exercice}

\begin{corrige}[Version : de l'espagnol vers le français administratif]
\begin{enumerate}
\item « J'accuse réception de votre lettre du 3~juin. »
\item « Suite à votre annonce parue dans le journal... »
\item « Je vous serais reconnaissant de bien vouloir m'envoyer de plus amples informations. »
\item « Je vous prie de bien vouloir accepter ma demande. »
\item « Veuillez trouver ci-joint les documents nécessaires. »
\item « Dans l'attente de votre réponse, je vous prie d'agréer mes salutations distinguées. »
\end{enumerate}
\end{corrige}

\begin{aretenir}[L'essentiel de la section]
\begin{itemize}
\item Les formules de correspondance sont \cle{codifiées} : on traduit la \emph{fonction}, pas les mots.
\item Ouverture : \esp{me permito dirigirme a usted} ; introduction : \esp{en relación con su anuncio} ; réception : \esp{acuso recibo de su carta}.
\item Demande : \esp{le agradecería que} + \cle{subjonctif imparfait} ; \esp{les ruego que} (présent) + \cle{subjonctif présent}.
\item Clôture : l'espagnol préfère la concision — \esp{en espera de su respuesta, le saluda atentamente}.
\item Pièges : \esp{adjunto} (ci-joint), \esp{acuse de recibo} (accusé de réception), \esp{atentamente} plutôt que \esp{cordialmente} au registre formel.
\end{itemize}
\end{aretenir}

\coursec{La lettre de réclamation}

Après avoir appris à traduire les formules de correspondance, nous allons maintenant les mettre en pratique dans un type de lettre très fréquent dans la vie quotidienne et professionnelle : la lettre de réclamation. Un produit défectueux, un colis perdu, un service mal rendu, une erreur de facturation : dans tous ces cas, il faut savoir écrire une lettre \emph{ferme mais courtoise}, en espagnol correct et structuré. C'est une compétence attendue à l'examen du Baccalauréat (expression écrite) et indispensable pour la vie active.

\courssub{Observation : une réclamation authentique}

Lisons d'abord ce modèle de lettre. Observez son organisation : chaque phrase a un rôle précis, le ton est maîtrisé, les faits sont datés et référencés.

\begin{textebox}[Una carta de reclamación]
Estimados señores:

El pasado 5 de marzo compré en su tienda en línea un teléfono móvil (pedido n.º 4582). Al recibirlo, comprobé que la pantalla estaba rota.

Les ruego que procedan al cambio del producto o al reembolso del importe en un plazo de quince días.

Atentamente,

Lucie Etoundi
\end{textebox}

\textbf{Traduction.} « Messieurs, Le 5 mars dernier, j'ai acheté dans votre boutique en ligne un téléphone portable (commande n.º 4582). En le recevant, j'ai constaté que l'écran était cassé. Je vous prie de procéder à l'échange du produit ou au remboursement du montant sous quinze jours. Cordialement, Lucie Etoundi. »

\begin{aretenir}[Glossaire : mots et expressions clés]
\begin{itemize}
\item \esp{tienda en línea} : boutique en ligne / site marchand
\item \esp{pedido (n.º)} : commande (numéro de commande)
\item \esp{al recibirlo} : en le recevant / à sa réception (gérondif \esp{al} + infinitif)
\item \esp{comprobar} : constater, vérifier, s'assurer que
\item \esp{la pantalla} : l'écran
\item \esp{roto, -a} : cassé, brisé (participe passé de \esp{romper}, accordé ici avec \esp{pantalla})
\item \esp{el cambio} : l'échange (remplacement du produit)
\item \esp{el reembolso} : le remboursement
\item \esp{el importe} : le montant, la somme
\item \esp{el plazo} : le délai (masculin : \esp{un plazo}, \esp{el plazo})
\item \esp{a la mayor brevedad} : dans les meilleurs délais / au plus vite (formule figée)
\end{itemize}
\end{aretenir}

\textbf{Questions de compréhension.}
\begin{enumerate}
\item ¿Qué compró Lucie y cuándo? \emph{(Qu'a acheté Lucie et quand ?)}
\item ¿Cuál es el problema del teléfono? \emph{(Quel est le problème du téléphone ?)}
\item ¿Qué soluciones propone Lucie? \emph{(Quelles solutions Lucie propose-t-elle ?)}
\item ¿Qué plazo fija a la empresa? \emph{(Quel délai fixe-t-elle à l'entreprise ?)}
\end{enumerate}

\courssub{Les spécificités de la lettre de réclamation}

\begin{definition}[¿Qué es una carta de reclamación ?]
Une lettre de réclamation (\esp{carta de reclamación}) est une lettre formelle dans laquelle on signale un problème (produit défectueux, retard de livraison, erreur de facturation, service insatisfaisant) et on \cle{demande une solution concrète} : remboursement, échange, réparation ou dédommagement.
\end{definition}

Trois qualités font une bonne réclamation, recevable par un service client et notée à l'examen :

\begin{itemize}
\item \textbf{Un ton ferme mais courtois.} On exprime son mécontentement sans insulte ni agressivité. La politesse (\esp{les ruego}, \esp{quisiera}, \esp{le agradecería}) renforce la crédibilité de la demande et montre qu'on connaît ses droits.
\item \textbf{Des faits précis et vérifiables.} Dates exactes, numéros de commande ou de facture, description objective du problème : \esp{El pasado 5 de marzo compré... (pedido n.º 4582)}. Sans ces détails, la lettre n'a aucune valeur juridique ni administrative.
\item \textbf{Une demande explicite et chiffrée.} Il faut dire clairement ce que l'on attend : \esp{Solicito el reembolso} (je demande le remboursement), \esp{el cambio del producto} (l'échange du produit), \esp{la reparación gratuita} (la réparation gratuite).
\end{itemize}

\courssub{La structure type d'une réclamation (5 pasos)}

Une lettre de réclamation suit toujours le même plan en cinq étapes. Mémorisez-le : il vous servira à l'examen comme dans la vie.

\begin{popfigure}
\begin{center}
\begin{tikzpicture}[
  node distance=0.3cm,
  etapa/.style={rectangle, rounded corners=3pt, draw=popblue, fill=popblueL, text width=12cm, align=left, font=\small, inner sep=5pt},
  flecha/.style={-{Stealth[length=2.5mm]}, thick, poporange}
]
\node[etapa] (e1) {\textbf{1. Hechos} : rappel des faits (date, lieu, numéro de commande/facture) — \esp{El pasado 5 de marzo compré en su tienda en línea... (pedido n.º 4582)}};
\node[etapa, below=of e1] (e2) {\textbf{2. Problema} : description objective du problème constaté — \esp{Al recibirlo, comprobé que la pantalla estaba rota.}};
\node[etapa, below=of e2] (e3) {\textbf{3. Petición} : demande explicite de solution — \esp{Solicito el cambio del producto o el reembolso del importe.}};
\node[etapa, below=of e3] (e4) {\textbf{4. Plazo} : délai raisonnable pour la réponse — \esp{En un plazo de quince días} / \esp{a la mayor brevedad}};
\node[etapa, below=of e4] (e5) {\textbf{5. Despedida} : formule de politesse de clôture — \esp{Atentamente,} / \esp{Quedo a la espera de su respuesta.}};
\draw[flecha] (e1) -- (e2);
\draw[flecha] (e2) -- (e3);
\draw[flecha] (e3) -- (e4);
\draw[flecha] (e4) -- (e5);
\end{tikzpicture}
\end{center}
\legende{Organigramme : les cinq étapes d'une réclamation efficace}
\end{popfigure}

\begin{methode}[Rédiger une lettre de réclamation en cinco pasos]
\begin{enumerate}
\item \textbf{Rappeler les faits (Hechos)} : Quoi, quand, où, quel numéro de référence. Utiliser le \textbf{passé simple} (\esp{compré}, \esp{recibí}, \esp{llegó}) pour des actions ponctuelles et terminées.
\item \textbf{Décrire le problème (Problema)} objectivement, sans exagération ni émotion excessive : \esp{El producto llegó defectuoso.} / \esp{El paquete no ha llegado.}
\item \textbf{Formuler la demande (Petición)} avec un verbe précis et le registre de politesse adapté : \esp{solicito} (indicatif, direct), \esp{les ruego que} + \textbf{subjonctif} (courtois), \esp{quisiera que} + \textbf{subjonctif} (très poli).
\item \textbf{Fixer un délai (Plazo)} si nécessaire : \esp{en un plazo de quince días}, \esp{a la mayor brevedad}, \esp{antes del día 20}.
\item \textbf{Terminer poliment (Despedida)} : \esp{Atentamente}, \esp{Quedo a la espera de su respuesta.}, \esp{Saludos cordiales}.
\end{enumerate}
\end{methode}

\courssub{Les formules utiles de la réclamation}

\begin{aretenir}[Frases clave para reclamar]
\begin{itemize}
\item \esp{Me dirijo a ustedes para presentar una reclamación.} « Je m'adresse à vous pour présenter une réclamation. » (Ouverture standard)
\item \esp{El producto llegó defectuoso / en mal estado.} « Le produit est arrivé défectueux / en mauvais état. »
\item \esp{El paquete no ha llegado todavía.} « Le colis n'est pas encore arrivé. »
\item \esp{La factura presenta un error en el importe.} « La facture présente une erreur sur le montant. »
\item \esp{Solicito el reembolso del importe.} « Je demande le remboursement du montant. » (Direct, indicatif)
\item \esp{Solicito el cambio del producto por uno en perfecto estado.} « Je demande l'échange du produit contre un en parfait état. »
\item \esp{Les ruego que solucionen este problema a la mayor brevedad.} « Je vous prie de résoudre ce problème dans les meilleurs délais. » (\textbf{Subjonctif présent} après \esp{rogar que})
\item \esp{Quisiera que me enviaran un técnico para la reparación.} « Je voudrais qu'on m'envoie un technicien pour la réparation. » (\textbf{Subjonctif imparfait} après \esp{quisiera que})
\item \esp{Quedo a la espera de su respuesta.} « Dans l'attente de votre réponse. » (Formule de clôture standard)
\end{itemize}
\end{aretenir}

\begin{propriete}[Les expressions de délai : précision et registre]
\begin{itemize}
\item \esp{a la mayor brevedad} = « dans les meilleurs délais », « au plus vite ». \textbf{Attention :} ne jamais traduire mot à mot par « avec la plus grande brièveté ».
\item \esp{en un plazo de quince días} = « sous quinze jours », « dans un délai de quinze jours ». Le mot \esp{plazo} signifie « délai » ; il est \textbf{masculin} : \esp{un plazo}, \esp{el plazo}.
\item \esp{lo antes posible} = « le plus tôt possible », « dès que possible ».
\item \esp{antes del día [fecha]} = « avant le [date] ».
\end{itemize}
\end{propriete}

\begin{savaistu}[Le consommateur au Cameroun et en zone hispanophone]
Au Cameroun, la \esp{Comisión Nacional de la Competencia y la Protección del Consumidor} (CNCC) et les associations de consommateurs (comme l'ACOC) encadrent les réclamations. En Espagne, le système est très codifié : on utilise les \esp{hojas de reclamaciones} (formulaires officiels de réclamation) que tout commerce doit mettre à disposition. En Guinée équatoriale, le cadre légal s'inspire du droit espagnol. Conserver la facture (\esp{factura} ou \esp{ticket}) et le bon de commande (\esp{pedido}) est la règle n°1 partout.
\end{savaistu}

\courssub{Nuancer la fermeté : conditionnel et subjonctif}

La difficulté de la réclamation est de rester \emph{ferme} sans devenir \emph{impoli}. L'espagnol dispose de deux outils grammaticaux que vous connaissez déjà : le \cle{conditionnel de courtoisie} et le \cle{subjonctif} après les verbes de volonté, demande ou émotion.

\begin{propriete}[Fermeté polie : les trois degrés de registre]
\begin{enumerate}
\item \textbf{Degré direct / ferme} (Indicatif présent) : \esp{Exijo el reembolso inmediato.} « J'exige le remboursement immédiat. » — Très fort, à réserver aux cas graves (non-réponse à une première lettre, dangerosité du produit).
\item \textbf{Degré courtois standard} (Conditionnel présent) : \esp{Esperaría una solución rápida.} « J'attendrais une solution rapide. » ; \esp{Le agradecería que me enviaran otro producto.} « Je vous serais reconnaissant de m'envoyer un autre produit. » (Conditionnel \esp{agradecería} + \textbf{Subjonctif imparfait} \esp{enviaran}).
\item \textbf{Degré très poli / administratif} (Subjonctif après verbes de demande) : \esp{Les ruego que solucionen este problema.} « Je vous prie de résoudre ce problème. » — Après \esp{rogar que}, \esp{agradecer que}, \esp{esperar que}, \esp{querer que}, \esp{necesitar que}, le verbe est \textbf{toujours au subjonctif} (présent pour l'avenir/proche, imparfait pour le conditionnel/hypothétique).
\end{enumerate}
\end{propriete}

\begin{attention}[Piège francophone : le subjonctif obligatoire]
En français, on dit « Je vous prie de résoudre » (infinitif) ou « Je souhaite qu'il résolve » (subjonctif). En espagnol, \textbf{après \esp{rogar que}, \esp{pedir que}, \esp{exigir que}, \esp{agradecer que}}, le subjonctif est \textbf{obligatoire}, même si le sujet est le même.
\begin{itemize}
\item Correct : \esp{Les ruego que me \textbf{devuelvan} el dinero.} (Subjonctif présent, 3e pers. pl.)
\item Incorrect : *\esp{Les ruego me devolver el dinero.} (Infinitif interdit ici)
\item Correct (conditionnel) : \esp{Le agradecería que me \textbf{devolvieran} el dinero.} (Subjonctif imparfait)
\end{itemize}
\end{attention}

\begin{exemplebox}[Ejemplos comentados : concordance des temps]
\esp{Quisiera que me \textbf{cambiaran} el teléfono.} « Je voudrais qu'on m'échange le téléphone. » (\esp{quisiera} conditionnel $\rightarrow$ \esp{cambiaran} subj. imparfait)

\esp{Esperaría una respuesta antes del viernes.} « J'attendrais une réponse avant vendredi. » (Conditionnel seul, pas de \esp{que} $\rightarrow$ pas de subjonctif)

\esp{Les ruego que \textbf{revisen} mi factura.} « Je vous prie de vérifier ma facture. » (\esp{ruego} présent $\rightarrow$ \esp{revisen} subj. présent)

\esp{Solicito que me \textbf{informen} del estado del envío.} « Je demande qu'on m'informe de l'état de l'envoi. » (\esp{solicito} présent indicatif + \esp{que} $\rightarrow$ subj. présent)
\end{exemplebox}

\courssub{Tableau récapitulatif : le lexique de la réclamation}

\begin{center}
\begin{tabularx}{\linewidth}{|X|X|X|}
\hline
\rowcolor{popblueL} Español & Français & Remarque / Collocation \\
\hline
la reclamación & la réclamation & \esp{presentar / poner una reclamación} \\
\hline
quejarse de & se plaindre de & \textbf{Réfléchi} + \esp{de} : \esp{Me quejo del retraso.} \\
\hline
reclamar & réclamer, exiger & \textbf{Non réfléchi} + COD : \esp{Reclamo el reembolso.} \\
\hline
el pedido / la comanda & la commande & \esp{número de pedido} (Am. Lat. souvent \esp{comanda}) \\
\hline
defectuoso / roto / dañado & défectueux / cassé / abîmé & \esp{llegó defectuoso}, \esp{está dañado} \\
\hline
el reembolso & le remboursement & Verbe : \esp{reembolsar} (ex. \esp{me reembolsaron}) \\
\hline
el cambio & l'échange & \esp{el cambio del producto}, \esp{solicitar el cambio} \\
\hline
la reparación & la réparation & \esp{reparación gratuita}, \esp{en garantía} \\
\hline
el plazo & le délai & \textbf{Masculin} : \esp{un plazo de 10 días} \\
\hline
a la mayor brevedad & dans les meilleurs délais & Formule figée, registre formel \\
\hline
la factura & la facture & \textbf{Faux ami} : pas « fait », mais document comptable \\
\hline
la garantía & la garantie & \esp{en período de garantía} \\
\hline
\end{tabularx}
\end{center}

\begin{attention}[Quejarse ou Reclamar ? L'erreur classique]
Ne confondez surtout pas ces deux verbes :
\begin{itemize}
\item \esp{quejarse} est \textbf{réfléchi} (\esp{me quejo, te quejas, se queja...}) et signifie « se plaindre, porter plainte ». Il s'emploie avec la préposition \esp{de} : \esp{Me quejo \textbf{de} la mala atención.}
\item \esp{reclamar} n'est \textbf{pas réfléchi} et signifie « réclamer, exiger quelque chose qui est dû ». Il prend un COD direct : \esp{Reclamo \textbf{el reembolso}.}
\end{itemize}
\emph{Erreur typique du francophone :} *\esp{Me reclamo el reembolso} $\rightarrow$ \textbf{Incorrect !} \\
\emph{Correct :} \esp{Reclamo el reembolso.} (J'exige le remboursement) \quad OU \quad \esp{Me quejo del retraso.} (Je me plains du retard).
\end{attention}

\courssub{Comparaison : Lettre de demande vs Lettre de réclamation}

La lettre de réclamation et la lettre de demande (stage, inscription, information) partagent la \emph{même architecture formelle} : en-tête (expéditeur/destinataire/date), objet, salutations, corps organisé, formule de politesse finale, signature. Ce qui change radicalement, c'est le \cle{registre} et le \cle{contenu du corps}.

\begin{center}
\begin{tabularx}{\linewidth}{|X|X|X|}
\hline
\rowcolor{popblueL} Élément & Lettre de demande (\esp{carta de solicitud}) & Lettre de réclamation (\esp{carta de reclamación}) \\
\hline
\textbf{Objectif} & Obtenir une faveur : stage, inscription, info & Obtenir réparation : échange, remboursement, réparation \\
\hline
\textbf{Ton} & Poli, valorisant, motivé & Ferme, factuel, courtois mais exigeant \\
\hline
\textbf{Corps (Contenu)} & Présentation du candidat + motivation + demande & Faits (date/réf) + Problème + Demande + Délai \\
\hline
\textbf{Verbes clés} & \esp{quisiera, me gustaría, solicito} & \esp{solicito, les ruego que, exijo, quisiera que} \\
\hline
\textbf{Temps dominants} & Présent, Conditionnel & \textbf{Passé simple} (faits) + Conditionnel/Subjonctif (demande) \\
\hline
\textbf{Formule finale} & \esp{Le agradezco de antemano su atención.} & \esp{Quedo a la espera de su respuesta.} \\
\hline
\end{tabularx}
\end{center}

Remarquez le rôle crucial du \cle{passé simple (pretérito indefinido)} dans la réclamation : il ancre les faits dans le passé révolu (\esp{compré}, \esp{recibí}, \esp{llegó}, \esp{comprobé}), alors que la lettre de demande vit surtout au présent (qui je suis) et au conditionnel (ce que je voudrais).

\begin{experience}[Jeu de rôle : « Au service clientèle »]
\textbf{Consigne.} Par binômes. L'un est le client (au téléphone ou au guichet), l'autre le responsable du service client.
\textbf{Scénario A :} Achat d'un ventilateur à Douala (magasin \emph{ElectroDuala}), facture n.º 215. Panne au bout de 2 jours. Le client veut l'échange.
\textbf{Scénario B :} Commande en ligne (site \emph{TecnoYa}), pedido n.º 889. Colis reçu ouvert, écran de tablette fissuré. Le client veut le remboursement.
\textbf{Contraintes linguistiques :} Le client doit utiliser \esp{les ruego que} + subjonctif / \esp{quisiera que} + subj. imparfait. Le responsable doit répondre avec \esp{le ofrecemos...}, \esp{procederemos a...}, \esp{en un plazo de...}.
\textbf{Débriefing :} Noter les formules utilisées, corriger les accords et les modes (indicatif/subjonctif).
\end{experience}

\courssub{Exercice résolu et entraînement}

\begin{exoresolu}[Ordenar una carta de reclamación]
Énoncé. Remettez dans l'ordre les phrases suivantes pour reconstituer une lettre de réclamation cohérente :
\begin{itemize}
\item[a)] \esp{Atentamente, Mballa Serge.}
\item[b)] \esp{Les ruego que me envíen otro producto en un plazo de diez días.}
\item[c)] \esp{El pasado lunes recibí el pedido n.º 731 de su tienda en línea.}
\item[d)] \esp{Al abrir el paquete, comprobé que los zapatos eran de talla diferente.}
\item[e)] \esp{Estimados señores:}
\end{itemize}
\tcblower
Solution détaillée. On applique le plan en cinq étapes :
\begin{enumerate}
\item Salutations (étape 5 de la structure) $\rightarrow$ \textbf{e}
\item Rappel des faits : date, numéro de commande (étape 1) $\rightarrow$ \textbf{c}
\item Description du problème (étape 2) $\rightarrow$ \textbf{d}
\item Demande explicite avec délai (étapes 3 \& 4) $\rightarrow$ \textbf{b}
\item Formule de politesse finale (étape 5) $\rightarrow$ \textbf{a}
\end{enumerate}
Ordre correct : \textbf{e – c – d – b – a}.

Lettre complète reconstituée :
\esp{Estimados señores: El pasado lunes recibí el pedido n.º 731 de su tienda en línea. Al abrir el paquete, comprobé que los zapatos eran de talla diferente. Les ruego que me envíen otro producto en un plazo de diez días. Atentamente, Mballa Serge.}
\end{exoresolu}

\begin{exercice}[A tu turno : Entraînement individuel]
\begin{enumerate}
\item \textbf{Traduction (Version) :} Traduisez en espagnol correct et formel :
« Le colis n'est pas encore arrivé. Je demande le remboursement du montant dans les meilleurs délais. »
\item \textbf{Grammaire (Quejarse / Reclamar) :} Complétez les phrases avec la forme conjuguée correcte du verbe indiqué (présent de l'indicatif) :
\begin{itemize}
\item Yo \_\_\_\_\_\_\_\_\_\_ (quejarse) del retraso del tren.
\item Los clientes \_\_\_\_\_\_\_\_\_\_ (reclamar) el reembolso de sus entradas.
\item Nosotros \_\_\_\_\_\_\_\_\_\_ (quejarse) de la mala calidad del servicio.
\item La empresa \_\_\_\_\_\_\_\_\_\_ (reclamar) el pago de la factura.
\end{itemize}
\item \textbf{Production écrite (Redacción) :} Rédigez une lettre de réclamation complète (80--100 mots).
\textbf{Situation :} Vous avez acheté un smartphone dans une boutique de Yaoundé (\emph{MovilCenter}) le 10 octobre. La facture est la n.º 442. Au bout de trois jours, la batterie ne tient plus la charge (se décharge en 2 heures). Vous voulez l'échange ou le remboursement sous 10 jours.
\textbf{Contraintes :} Respectez les 5 étapes (Hechos, Problema, Petición, Plazo, Despedida). Utilisez au moins une forme de conditionnel de courtoisie (\esp{quisiera que} / \esp{le agradecería que} + subjonctif) et une formule \esp{les ruego que} + subjonctif.
\end{enumerate}
\end{exercice}

\begin{corrige}[Exercice « A tu turno »]
\begin{enumerate}
\item \esp{El paquete no ha llegado todavía. Solicito el reembolso del importe a la mayor brevedad.} \\
(Alternatif : \esp{Les ruego que me reembolsen el importe a la mayor brevedad.})
\item
\begin{itemize}
\item Yo \textbf{me quejo} del retraso del tren. (Réfléchi + \esp{de})
\item Los clientes \textbf{reclaman} el reembolso de sus entradas. (Non réfléchi + COD)
\item Nosotros \textbf{nos quejamos} de la mala calidad del servicio. (Réfléchi 1ère pers. pl. + \esp{de})
\item La empresa \textbf{reclama} el pago de la factura. (Non réfléchi 3e pers. sg. + COD)
\end{itemize}
\item \textbf{Modèle de rédaction (environ 90 mots) :}
\begin{textebox}[Modelo de carta]
Estimados señores:

El pasado 10 de octubre compré un smartphone en su tienda de Yaundé (factura n.º 442). Al cabo de tres días, he comprobado que la batería no mantiene la carga: se descarga completamente en dos horas.

Les ruego que procedan al cambio del aparato por uno nuevo en un plazo de diez días. Le agradecería que me confirmaran la recepción de esta carta.

Quedo a la espera de su respuesta.

Atentamente,

[Votre Nom]
\end{textebox}
\textbf{Analyse du modèle :} \esp{compré} (passé simple, fait) ; \esp{he comprobé} (passé composé, conséquence présente) ; \esp{mantiene} (présent indicatif, description état actuel) ; \esp{les ruego que procedan} (courtoisie + subj. présent) ; \esp{le agradecería que me confirmaran} (conditionnel + subj. imparfait) ; \esp{Quedo a la espera} (formule finale).
\end{enumerate}
\end{corrige}

\begin{aretenir}[Lo esencial : La réclamation parfaite]
Une réclamation efficace = des \cle{faits précis} (passé simple : \esp{compré, recibí}) + un \cle{problème décrit objectivement} (\esp{llegó defectuoso, no funciona}) + une \cle{demande explicite} (\esp{solicito} / \esp{les ruego que} + \textbf{subjonctif} / \esp{quisiera que} + \textbf{subj. imparfait}) + un \cle{délai raisonnable} (\esp{en un plazo de...}, \esp{a la mayor brevedad}) + une \cle{formule de politesse de clôture} (\esp{Atentamente, Quedo a la espera...}).
\\
\textbf{Trio gagnant :} \emph{Ferme, Précis, Courtois.}
\end{aretenir}

\coursec{Production guidée : la lettre de motivation (100 mots)}

Après avoir appris à rédiger une lettre de réclamation, nous allons maintenant utiliser les mêmes outils (formules de politesse, conditionnel de courtoisie, structure formelle) pour un exercice légèrement différent et très fréquent à l'examen : la \cle{carta de motivación}, la lettre de motivation. C'est la lettre que tu écriras un jour pour demander un stage, une inscription dans une université ou une bourse d'études. Autant s'entraîner dès maintenant !

\courssub{Objectif et observation : un modèle de lettre}

L'objectif de cette production est simple : rédiger une lettre de motivation \textbf{courte} (environ 100 mots), \textbf{claire} (3 paragraphes) et \textbf{polie} (registre formel, conditionnel de courtoisie). Observons d'abord un modèle rédigé par le professeur, avant d'en dégager la méthode.

\begin{textebox}[Una carta de motivación : modelo del profesor]
Yaoundé, 15 de mayo de 2024\par
\medskip
A la atención del Director de la Fundación Nguema,\par
\medskip
Muy señor mío:\par
\medskip
Me llamo Amina Mbarga y curso la clase de Primera en el Liceo Bilingüe de Yaoundé. Quisiera solicitar una beca de estudios para el próximo curso académico.\par
\medskip
Estoy muy interesada en las lenguas extranjeras, especialmente el español, y obtengo buenas notas en todas las asignaturas. Creo que tengo las cualidades necesarias para aprovechar esta oportunidad, soy trabajadora, puntual y me gusta ayudar a los demás. Además, participo en el club de teatro de mi liceo.\par
\medskip
Estoy disponible inmediatamente para una entrevista. Le agradecería que me enviara información sobre los documentos necesarios.\par
\medskip
En espera de su respuesta, le saluda atentamente,\par
\medskip
Amina Mbarga
\end{textebox}

\textbf{Glossaire :} \esp{curso la clase de Primera} = je suis en classe de Première ; \esp{una beca de estudios} = une bourse d'études ; \esp{aprovechar} = profiter de, tirer parti de ; \esp{puntual} = ponctuel ; \esp{una entrevista} = un entretien ; \esp{en espera de su respuesta} = dans l'attente de votre réponse.

Compte les mots : la lettre contient environ 110 mots, c'est-à-dire exactement ce que l'on attend de toi à l'examen.

\courssub{Le plan en trois paragraphes}

Une lettre de motivation de 100 mots ne laisse aucune place au bavardage. Elle suit toujours la même architecture en trois paragraphes, chacun avec une fonction précise.

\begin{popfigure}
\begin{center}
\begin{tikzpicture}[
  box/.style={rectangle, rounded corners=4pt, draw=popblue, fill=popblueL, text width=3.4cm, align=center, minimum height=1.6cm, font=\small},
  arr/.style={-{Stealth[length=3mm]}, thick, poporange}
]
\node[box, align=center] (p1) at (0,0){\textbf{Párrafo 1}\\Presentación y objeto\\ \emph{¿Quién soy? ¿Qué pido?}};
\node[box, draw=popgreen, fill=popgreenL, align=center] (p2) at (4.6,0){\textbf{Párrafo 2}\\Motivación y cualidades\\ \emph{¿Por qué yo?}};
\node[box, draw=poppurple, fill=poppurpleL, align=center] (p3) at (9.2,0){\textbf{Párrafo 3}\\Disponibilidad y gracias\\ \emph{¿Qué espero?}};
\draw[arr] (p1) -- (p2);
\draw[arr] (p2) -- (p3);
\end{tikzpicture}
\end{center}
\legende{Les trois paragraphes de la lettre de motivation et leurs fonctions}
\end{popfigure}

\begin{aretenir}[La structure en 3 paragraphes]
\begin{enumerate}
\item \textbf{Présentation et objet} : qui je suis (nom, classe, établissement) et ce que je demande (stage, inscription, bourse).
\item \textbf{Motivations et qualités} : pourquoi cette demande m'intéresse et pourquoi je suis le bon candidat (2 ou 3 arguments, pas une liste).
\item \textbf{Disponibilité et remerciements} : je me dis disponible, je remercie, je salue poliment.
\end{enumerate}
\end{aretenir}

\courssub{La banque de phrases}

Pour rédiger vite et bien, apprends par cœur cette banque de phrases. Chaque phrase correspond à une étape du plan.

\begin{center}
\begin{tabularx}{\linewidth}{|X|X|}
\hline
\rowcolor{popblueL} \textbf{Español} & \textbf{Français} \\
\hline
Me llamo... y curso la clase de Primera en... & Je m'appelle... et je suis en Première à... \\
\hline
Quisiera solicitar... & Je voudrais demander / solliciter... \\
\hline
Estoy muy interesado/a en... & Je suis très intéressé(e) par... \\
\hline
Creo que tengo las cualidades necesarias. & Je pense avoir les qualités nécessaires. \\
\hline
Soy trabajador/a, puntual y serio/a. & Je suis travailleur(-euse), ponctuel(le) et sérieux(-euse). \\
\hline
Estoy disponible desde el lunes. & Je suis disponible à partir de lundi. \\
\hline
Estoy disponible inmediatamente. & Je suis disponible immédiatement. \\
\hline
Le agradecería que me enviara información. & Je vous serais reconnaissant(e) de m'envoyer des informations. \\
\hline
En espera de su respuesta, le saluda atentamente. & Dans l'attente de votre réponse, je vous salue cordialement. \\
\hline
\end{tabularx}
\end{center}

\begin{exemplebox}[Exemples traduits]
\esp{Creo que tengo las cualidades necesarias para este puesto.} « Je pense avoir les qualités nécessaires pour ce poste. »\par
\esp{Estoy disponible inmediatamente.} « Je suis disponible immédiatement. »\par
\esp{Quisiera solicitar un período de prácticas en su empresa.} « Je voudrais demander un stage dans votre entreprise. »\par
\esp{Le agradecería que tuviera en cuenta mi candidatura.} « Je vous serais reconnaissant(e) de prendre en compte ma candidature. »
\end{exemplebox}

\begin{attention}[Pièges des francophones]
\begin{itemize}
\item \esp{Quisiera solicitar...} et non \esp{*quiero solicitar} : le présent \esp{quiero} (« je veux ») est trop direct, presque impoli dans une lettre formelle. Utilise toujours le conditionnel de courtoisie.
\item Attention au genre : un garçon écrit \esp{interesado}, \esp{trabajador} ; une fille écrit \esp{interesada}, \esp{trabajadora}.
\item \esp{Curso la clase de Primera} : le verbe \esp{cursar} signifie « suivre des études », pas « courir » !
\item N'oublie pas le subjonctif après \esp{le agradecería que} : \esp{que me enviara}, \esp{que tuviera en cuenta}.
\end{itemize}
\end{attention}

\courssub{Les contraintes formelles à respecter}

\begin{propriete}[Les cinq contraintes de la lettre de motivation]
\begin{enumerate}
\item \textbf{Le vouvoiement} : toujours \esp{usted} (et les possessifs \esp{su}, \esp{sus}), jamais \esp{tú}.
\item \textbf{Le conditionnel de politesse} : \esp{quisiera}, \esp{me gustaría}, \esp{le agradecería}.
\item \textbf{Les formules d'ouverture} : \esp{Muy señor mío:} / \esp{Estimado señor Director:} (avec deux-points, jamais de virgule).
\item \textbf{Les formules de clôture} : \esp{En espera de su respuesta, le saluda atentamente} / \esp{Atentamente}.
\item \textbf{La longueur} : environ 100 mots, ni 50 ni 200.
\end{enumerate}
\end{propriete}

\begin{attention}[Une lettre de motivation n'est pas un CV]
Ne transforme pas ta lettre en liste : \esp{*Soy serio. Soy puntual. Soy trabajador. Tengo dieciséis años.} est une énumération scolaire qui ne convainc personne. Une lettre de motivation \textbf{argumente} : choisis 2 ou 3 qualités et relie-les à ta demande en phrases complètes. Compare :\par
\esp{*Soy trabajador. Me gusta el español.} (liste sèche)\par
\esp{Estoy muy interesado en el español y, como soy trabajador, creo que puedo aprovechar esta beca.} (argumentation : qualité + lien avec la demande)
\end{attention}

\courssub{Méthode : rédiger en cinq étapes}

\begin{methode}[Comment rédiger sa lettre de motivation en 5 étapes]
\begin{enumerate}
\item \textbf{Lire la consigne et identifier le destinataire} : qui reçoit la lettre ? Un directeur d'entreprise, un recteur d'université, une fondation ? Cela détermine la formule d'ouverture.
\item \textbf{Rédiger le plan au brouillon} : note 2 ou 3 idées par paragraphe (présentation / motivations / disponibilité). N'oublie pas la mise en page : lieu et date en haut à droite, destinataire à gauche, formules d'ouverture et de clôture.
\item \textbf{Écrire avec la banque de phrases} : assemble les phrases apprises en les adaptant (genre, nom du poste, dates).
\item \textbf{Vérifier le conditionnel et le subjonctif} : chaque demande doit être au conditionnel (\esp{quisiera}, \esp{agradecería}) ; après \esp{agradecería que}, subjonctif imparfait (\esp{enviara}, \esp{tuviera}).
\item \textbf{Compter les mots et relire} : vise 100 mots (90--110 acceptés) ; vérifie les accents (\esp{está}, \esp{información}, \esp{atención}) et la concordance du genre.
\end{enumerate}
\end{methode}

\begin{exoresolu}[Sujet type examen]
\textbf{Énoncé :} \esp{Usted es alumno/a de Primera en el Liceo de Douala. Escriba una carta de motivación (unas 100 palabras) al director del Instituto Cervantes para solicitar una inscripción en un curso de verano de español.}
\tcblower
\textbf{Solution détaillée (lettre corrigée) :}\par
\medskip
\esp{Douala, 3 de junio de 2024}\par
\esp{Al Director del Instituto Cervantes,}\par
\medskip
\esp{Muy señor mío:}\par
\medskip
\esp{Me llamo Junior Etoundi y curso Primera en el Liceo de Douala. Quisiera solicitar una inscripción en su curso de verano de español.}\par
\medskip
\esp{Estoy muy interesado en la lengua y la cultura españolas. Creo que tengo las cualidades necesarias para seguir este curso, soy trabajador, serio y me encanta hablar con los demás.}\par
\medskip
\esp{Estoy disponible desde el 1 de julio. Le agradecería que me enviara información sobre las fechas y los precios.}\par
\medskip
\esp{En espera de su respuesta, le saluda atentamente,}\par
\esp{Junior Etoundi}\par
\medskip
\textbf{Commentaire :} la lettre compte environ 105 mots. Paragraphe 1 = présentation + objet (\esp{quisiera solicitar}) ; paragraphe 2 = motivations et qualités argumentées ; paragraphe 3 = disponibilité + demande polie avec subjonctif (\esp{enviara}) ; formules d'ouverture et de clôture correctes.
\end{exoresolu}

\courssub{Grille d'auto-évaluation}

Avant de rendre ta copie, évalue-toi toi-même avec cette grille, identique à celle de l'examen.

\begin{center}
\begin{tabularx}{\linewidth}{|X|c|}
\hline
\rowcolor{popgreenL} \textbf{Critère} & \textbf{Points} \\
\hline
Structure respectée (3 paragraphes, ouverture, clôture) & 2 pts \\
\hline
Formules de politesse (\esp{usted}, ouverture, clôture) & 2 pts \\
\hline
Conditionnel correct (\esp{quisiera}, \esp{agradecería} + subjonctif) & 2 pts \\
\hline
Lexique approprié (vocabulaire de la demande et des qualités) & 2 pts \\
\hline
Orthographe et accents & 2 pts \\
\hline
\textbf{Total} & \textbf{10 pts} \\
\hline
\end{tabularx}
\end{center}

\begin{savaistu}[La lettre de motivation au Bac]
Au baccalauréat, la lettre de motivation est un sujet très fréquent d'expression écrite en espagnol. Retiens cette conversion utile : 100 mots correspondent à environ 10 lignes de copie, soit 8 à 10 phrases. Un bon candidat rédige d'abord son plan au brouillon en 2 minutes, puis écrit sans ratures. Les correcteurs récompensent surtout trois choses : le respect de la structure, les formules de politesse et l'emploi correct du conditionnel. C'est exactement la grille d'auto-évaluation ci-dessus !
\end{savaistu}

\courssub{Correction collective de la lettre modèle}

Reprenons ensemble la lettre d'Amina Mbarga (modèle du début de la section) et vérifions chaque critère de la grille.

\begin{itemize}
\item \textbf{Structure (2/2)} : trois paragraphes nets, ouverture \esp{Muy señor mío:}, clôture \esp{le saluda atentamente}.
\item \textbf{Politesse (2/2)} : \esp{usted} partout (\esp{le agradecería}, \esp{su respuesta}).
\item \textbf{Conditionnel (2/2)} : \esp{Quisiera solicitar}, \esp{Le agradecería que me enviara} (conditionnel + subjonctif imparfait).
\item \textbf{Lexique (2/2)} : \esp{beca}, \esp{curso académico}, \esp{entrevista}, qualités argumentées.
\item \textbf{Orthographe (2/2)} : accents corrects (\esp{atención}, \esp{está}, \esp{información}).
\end{itemize}

Total : 10/10. C'est le niveau visé. À ton tour de produire la même qualité !

\courssub{Complément Bac (hors programme strict) : le courriel formel}

\begin{attention}[Hors programme strict — pour aller plus loin]
De plus en plus, l'examen propose de rédiger un \esp{correo electrónico formal} plutôt qu'une lettre sur papier. La bonne nouvelle : le contenu reste identique (mêmes 3 paragraphes, mêmes formules de politesse). Seule la présentation change.
\end{attention}

\begin{center}
\begin{tabularx}{\linewidth}{|X|X|}
\hline
\rowcolor{poppurpleL} \textbf{Carta formal (papier)} & \textbf{Correo electrónico formal} \\
\hline
Lieu et date en haut à droite & \esp{Asunto: Solicitud de beca} (objet du message) \\
\hline
\esp{Muy señor mío:} & \esp{Estimado señor Director:} \\
\hline
Adresse du destinataire en haut & Pas d'adresse postale \\
\hline
Signature manuscrite & Nom complet + établissement sous le message \\
\hline
\end{tabularx}
\end{center}

\begin{exemplebox}[Modèle de courriel formel]
\esp{Asunto: Solicitud de inscripción al curso de verano}\par
\medskip
\esp{Estimado señor Director:}\par
\medskip
\esp{Me llamo Amina Mbarga y curso la clase de Primera en el Liceo Bilingüe de Yaoundé. Quisiera solicitar una inscripción en su curso de verano de español. [...]}\par
\medskip
\esp{En espera de su respuesta, reciba un cordial saludo.}\par
\esp{Amina Mbarga — Liceo Bilingüe de Yaoundé}
\end{exemplebox}

\begin{exercice}[À toi de jouer]
Rédige une lettre de motivation d'environ 100 mots : tu demandes un stage (\esp{un período de prácticas}) dans une entreprise de cacao à Douala pendant les vacances. Respecte le plan en 3 paragraphes, utilise au moins deux conditionnels de courtoisie et une formule avec subjonctif. Auto-évalue-toi ensuite avec la grille sur 10 points.
\end{exercice}

\begin{corrige}[Exercice — éléments de correction]
Une bonne copie contient : ouverture \esp{Muy señor mío:} ; paragraphe 1 avec \esp{Me llamo... y curso Primera en...} + \esp{Quisiera solicitar un período de prácticas en su empresa} ; paragraphe 2 avec \esp{Estoy muy interesado/a en...} + 2 ou 3 qualités argumentées ; paragraphe 3 avec \esp{Estoy disponible desde...} + \esp{Le agradecería que me enviara...} ; clôture \esp{En espera de su respuesta, le saluda atentamente}. Longueur : 90 à 110 mots. Accents obligatoires : \esp{período}, \esp{información}, \esp{está}.
\end{corrige}

\newpage

\coursec{Activité d'intégration}

\courssub{Objectifs de l'activité}

À l'issue de cette activité d'intégration, l'élève sera capable de :

\begin{itemize}
\item rédiger une lettre de motivation formelle complète (environ 100 mots) en respectant la structure en trois paragraphes (présentation, motivation, demande et remerciements) ;
\item employer correctement au moins deux \cle{conditionnels de courtoisie} (\esp{quisiera}, \esp{le agradecería que}, \esp{me gustaría}) ;
\item utiliser les \cle{formules de correspondance} formelles : en-tête, formule d'appel (\esp{Estimado señor consejero cultural:}), formule de clôture (\esp{Le saluda atentamente}) ;
\item mobiliser le lexique de la demande (\esp{solicitar}, \esp{prácticas}, \esp{candidatura}, \esp{adjuntar}) ;
\item jouer le rôle de destinataire : évaluer la lettre d'un camarade avec une grille et rédiger une courte \cle{lettre de réponse formelle} (acceptation ou demande de compléments) ;
\item traduire mentalement, de l'espagnol vers le français et inversement, les formules figées de la correspondance administrative.
\end{itemize}

\courssub{Situation}

Tu es élève hispanisant en classe de Première A au lycée de Yaoundé. Ce matin, le proviseur a affiché au tableau d'information du CDI l'annonce suivante, transmise par le service culturel de l'ambassade d'Espagne au Cameroun.

\begin{textebox}[Anuncio — Embajada de España en Yaundé]
\textbf{CONVOCATORIA DE PRÁCTICAS DE DESCUBRIMIENTO}

La Embajada de España en Yaundé ofrece a los alumnos de español de Primera una \textbf{estancia de prácticas de dos semanas} en su servicio cultural, durante las vacaciones escolares.

Los estudiantes seleccionados descubrirán el trabajo diplomático, ayudarán en la organización de la Semana del Libro y participarán en actividades culturales (cine, música, exposiciones).

Para presentar su candidatura, los interesados deben enviar una \textbf{carta de motivación en español} (unas 100 palabras) al señor consejero cultural, antes del 15 de mayo.

\textbf{Dirección:} Embajada de España, Servicio Cultural, Yaundé, Camerún.

\textit{¡Anímense a participar!}
\end{textebox}

\textbf{Glossaire :} \esp{convocatoria} = appel à candidatures ; \esp{estancia de prácticas} = stage ; \esp{seleccionados} = sélectionnés ; \esp{presentar su candidatura} = poser sa candidature ; \esp{antes del} = avant le ; \esp{anímense} = lancez-vous (impératif de \esp{animarse}).

Tu décides de tenter ta chance : tu rédiges ta lettre de motivation. Puis, avec ton binôme, tu échanges les rôles : chacun devient le conseiller culturel qui lit, évalue et répond à la lettre reçue.

\courssub{Consignes}

\textbf{Tâche 1 — Compréhension de l'annonce (10 min).} Lis attentivement l'annonce ci-dessus et réponds en français : a) Qui propose le stage ? b) Quelle est sa durée ? c) Quelles activités sont proposées ? d) Quel document faut-il envoyer et avant quelle date ? e) À qui doit-on adresser la lettre ?

\textbf{Tâche 2 — Préparation du lexique (10 min).} Complète ton brouillon avec au moins six mots ou expressions de la leçon : \esp{solicitar}, \esp{estancia de prácticas}, \esp{presentar mi candidatura}, \esp{adjuntar}, \esp{quedar a su disposición}, \esp{en espera de su respuesta}. Vérifie le sens de chacun dans ton cours.

\textbf{Tâche 3 — Rédaction de la lettre de motivation (30 min).} Rédige ta lettre (environ 100 mots) en respectant le plan en trois paragraphes :
\begin{enumerate}
\item \textbf{Paragraphe 1 :} présentation (nom, classe, établissement) et objet de la lettre ;
\item \textbf{Paragraphe 2 :} ta motivation (pourquoi l'espagnol, pourquoi ce stage, tes qualités) ;
\item \textbf{Paragraphe 3 :} la demande formulée avec politesse, ta disponibilité et les remerciements.
\end{enumerate}
Obligations : au moins \textbf{deux conditionnels de courtoisie} (\esp{quisiera}, \esp{me gustaría}, \esp{le agradecería que} + subjonctif), une \textbf{formule d'appel} formelle, une \textbf{formule de clôture} formelle, et l'en-tête (lieu, date, destinataire).

\textbf{Tâche 4 — Échange et évaluation par binôme (20 min).} Échange ta lettre avec celle de ton camarade. Tu es maintenant le conseiller culturel. Évalue la lettre reçue avec la grille ci-dessous (2 points par critère, sur 10) :

\begin{center}
\begin{tabularx}{\linewidth}{|X|c|}
\hline
\rowcolor{popblueL} \textbf{Critère} & \textbf{Points} \\
\hline
Structure formelle complète (en-tête, appel, clôture) & /2 \\
\hline
Plan en trois paragraphes respecté & /2 \\
\hline
Deux conditionnels de courtoisie corrects & /2 \\
\hline
Lexique de la demande (au moins quatre mots) & /2 \\
\hline
Orthographe et accents (ñ, á, é, í, ó, ú) & /2 \\
\hline
\end{tabularx}
\end{center}

\textbf{Tâche 5 — Lettre de réponse formelle (20 min).} Toujours dans le rôle du conseiller culturel, rédige une courte réponse (5 à 8 lignes) : soit une \textbf{acceptation} (\esp{Nos complace informarle que su candidatura ha sido aceptada...}), soit une \textbf{demande de compléments} (\esp{Le agradeceríamos que nos enviara...}). Utilise au moins un conditionnel de politesse et une formule de clôture.

\textbf{Tâche 6 — Mise en commun (15 min).} Les trois meilleures lettres de la classe sont lues à voix haute. La classe commente collectivement : les points forts, les erreurs corrigées, la prononciation des formules.

\courssub{Ressources à mobiliser}

\begin{aretenir}[Boîte à outils de la lettre formelle]
\begin{itemize}
\item \textbf{En-tête :} \esp{Yaundé, a 10 de mayo de 2025} ; destinataire : \esp{Señor consejero cultural de la Embajada de España}.
\item \textbf{Appel :} \esp{Estimado señor consejero:} (deux-points, jamais de virgule).
\item \textbf{Conditionnels de courtoisie :} \esp{Quisiera solicitar...} « je voudrais demander... » ; \esp{Me gustaría participar en...} « j'aimerais participer à... » ; \esp{Le agradecería que me enviara...} « je vous serais reconnaissant de m'envoyer... » (attention : \esp{agradecería que} + \textbf{subjonctif imparfait}).
\item \textbf{Lexique :} \esp{solicitar} (demander), \esp{la estancia de prácticas} (le stage), \esp{presentar mi candidatura} (poser ma candidature), \esp{adjuntar} (joindre), \esp{quedar a su disposición} (rester à votre disposition).
\item \textbf{Clôture :} \esp{En espera de su respuesta, le saluda atentamente.} « Dans l'attente de votre réponse, je vous salue respectueusement. »
\end{itemize}
\end{aretenir}

\begin{attention}[Pièges à éviter]
Ne traduis pas mot à mot « je vous serais reconnaissant » par \esp{yo sería reconocido} : dis \esp{le agradecería}. N'oublie pas le subjonctif après \esp{le agradecería que} : \esp{que me enviara}, et non \esp{que me envía}. Enfin, \esp{asistir a} signifie « assister à » et non « aider » ; pour « aider », utilise \esp{ayudar}.
\end{attention}

\courssub{Proposition de production et corrigé commenté}

\begin{exoresolu}[Modèle de lettre de motivation (environ 100 mots)]
Rédige ta lettre de motivation pour le stage de l'ambassade d'Espagne, puis compare-la au modèle ci-dessous.
\tcblower
\textbf{Production modèle :}

\esp{Yaundé, a 10 de mayo de 2025}

\esp{Señor consejero cultural de la Embajada de España}

\esp{Estimado señor consejero:}

\esp{Me llamo Amina Mbarga y soy alumna de Primera A del Liceo de Yaundé. Quisiera presentar mi candidatura para la estancia de prácticas de dos semanas en su servicio cultural.}

\esp{Estudio español desde hace cuatro años y me apasiona la cultura hispánica. Me gustaría descubrir el trabajo diplomático y ayudar en la organización de la Semana del Libro. Soy seria, puntual y hablo también francés e inglés.}

\esp{Le agradecería que me enviara más información sobre las fechas. Quedo a su disposición para una entrevista.}

\esp{En espera de su respuesta, le saluda atentamente.}

\esp{Amina Mbarga}

\textbf{Commentaire des choix de langue :}
\begin{itemize}
\item \textbf{En-tête et appel} conformes : lieu et date, destinataire, \esp{Estimado señor consejero:} avec deux-points.
\item \textbf{Deux conditionnels de courtoisie} : \esp{Quisiera presentar} (paragraphe 1) et \esp{Me gustaría descubrir} (paragraphe 2) ; plus un troisième, \esp{Le agradecería que me enviara}, suivi correctement du subjonctif imparfait \esp{enviara}.
\item \textbf{Plan en trois paragraphes} respecté : présentation et objet ; motivation et qualités ; demande, disponibilité et remerciements.
\item \textbf{Lexique de la demande} : \esp{presentar mi candidatura}, \esp{estancia de prácticas}, \esp{quedo a su disposición}, \esp{en espera de su respuesta}.
\item \textbf{Clôture formelle} : \esp{le saluda atentamente}, suivie de la signature.
\item \textbf{Accents} soignés : \esp{información}, \esp{apasiona}, \esp{hispánica}, \esp{también}.
\end{itemize}

\textbf{Modèle de réponse du conseiller culturel (acceptation) :}

\esp{Estimada señorita Mbarga:}

\esp{Nos complace informarle que su candidatura ha sido aceptada. Le agradeceríamos que confirmara su disponibilidad antes del 20 de mayo.}

\esp{Le saluda cordialmente,}

\esp{El consejero cultural}

(« Nous avons le plaisir de vous informer que votre candidature a été acceptée. Nous vous serions reconnaissants de confirmer votre disponibilité avant le 20 mai. »)
\end{exoresolu}

\courssub{Bilan de l'activité}

Cette activité t'a permis de mobiliser dans une situation authentique l'ensemble des savoir-faire de la leçon : la \cle{structure} de la lettre formelle (en-tête, appel, trois paragraphes, clôture), le \cle{conditionnel de courtoisie} pour formuler une demande polie, le \cle{lexique} de la candidature et la \cle{traduction} des formules figées entre le français et l'espagnol. En jouant le rôle du conseiller culturel, tu as aussi développé un regard critique : évaluer une lettre, repérer les erreurs d'accents ou de registre, et rédiger une réponse administrative sont des compétences directement réutilisables à l'examen (expression écrite et thème) comme dans la vie réelle — stage, inscription universitaire, demande d'information. Retiens l'essentiel : une lettre formelle réussie est une lettre \textbf{claire, structurée et courtoise}.

\newpage

\coursec{Méthodes et exercices résolus}

\courssub{Savoir-faire 1 : Rédiger une lettre de demande formelle}

Après avoir étudié la lettre amicale, nous abordons ici la correspondance formelle : demande de stage, d'inscription, d'information. C'est un exercice imposé fréquent à l'examen de classe de Première et au Baccalauréat (expression écrite). Sa réussite repose sur deux piliers : une architecture fixe et un registre de langue soutenu.

\begin{definition}[Lexique thématique de la correspondance formelle]
\begin{itemize}
\item \textbf{L'en-tête et l'appel} : \esp{el remitente} (l'expéditeur), \esp{el destinatario} (le destinataire), \esp{el asunto} (l'objet), \esp{Estimado señor / Estimada señora:} (Monsieur / Madame).
\item \textbf{La demande} : \esp{solicitar} (demander officiellement), \esp{una solicitud} (une demande écrite), \esp{pedir información} (demander des renseignements), \esp{un formulario de inscripción} (un formulaire d'inscription), \esp{las prácticas} (le stage), \esp{un puesto de trabajo} (un poste de travail).
\item \textbf{La réclamation} : \esp{quejarse de} (se plaindre de), \esp{reclamar} (réclamer), \esp{el producto defectuoso} (le produit défectueux), \esp{la factura} (la facture), \esp{la garantía} (la garantie), \esp{el reembolso} (le remboursement), \esp{devolver el dinero} (rembourser l'argent).
\item \textbf{Les formules finales} : \esp{En espera de su respuesta} (dans l'attente de votre réponse), \esp{Atentamente} (cordialement, formel), \esp{Reciba un cordial saludo} (recevez mes salutations).
\end{itemize}
Prononciation : \esp{solicitud} se prononce « so-li-si-toud » ; \esp{reembolso} « ré-ém-bol-so ».
\end{definition}

\begin{methode}[Rédiger une lettre de demande (stage, inscription, information)]
Une lettre de demande formelle suit une architecture rigide. Respecte toujours ces étapes :
\begin{enumerate}
\item \textbf{L'en-tête} : lieu et date à droite (\esp{Yaoundé, a 15 de mayo de 2024}), puis le destinataire à gauche (\esp{Al Director de la Empresa CacaoCam, Douala}).
\item \textbf{La formule d'appel} : \esp{Estimado señor Director:} ou \esp{Muy señor mío:} — toujours suivie de deux points, jamais de virgule.
\item \textbf{L'objet} : \esp{Asunto: solicitud de prácticas} (facultatif mais apprécié).
\item \textbf{Le premier paragraphe} : se présenter brièvement et annoncer la demande (\esp{Me dirijo a usted para solicitar...}).
\item \textbf{Le développement} : justifier la demande (motivation, qualités, disponibilités) en deux ou trois phrases.
\item \textbf{La formule de politesse finale} : \esp{Le agradecería que me enviara una respuesta lo antes posible.} ou \esp{En espera de su respuesta, le saluda atentamente.}
\item \textbf{La signature} : prénom et nom, éventuellement la fonction.
\end{enumerate}
Réflexes : tutoiement interdit ; employer \esp{usted} et les possessifs \esp{su, sus} ; phrases courtes ; aucune abréviation familière.
\end{methode}

\begin{exoresolu}[Rédiger une demande de stage]
Énoncé : Tu es élève de Première A au Lycée de Nkol-Eton (Yaoundé). Rédige en espagnol une lettre de demande de stage (\esp{prácticas}) de deux semaines adressée au directeur de la radio \esp{Radio Litoral} de Douala (60 à 80 mots, sans compter l'en-tête et les formules).
\tcblower
\textbf{Raisonnement.} On applique la méthode : en-tête, appel avec deux points, présentation, demande avec le conditionnel de courtoisie \esp{quisiera}, justification, formule finale, signature. On emploie \esp{usted} partout.

\textbf{Réponse rédigée :}

\esp{Yaoundé, a 12 de junio de 2024}

\esp{Al Director de Radio Litoral, Douala}

\esp{Estimado señor Director:}

\esp{Soy alumno de Primera A del Liceo de Nkol-Eton, en Yaoundé. Me dirijo a usted para solicitar unas prácticas de dos semanas en su emisora durante las vacaciones de julio.}

\esp{Me apasiona el periodismo y quisiera descubrir el trabajo de una redacción de radio. Soy serio, puntual y hablo español y francés con fluidez.}

\esp{Le agradecería que me informara sobre las condiciones de admisión. En espera de su respuesta, le saluda atentamente.}

\esp{Jean-Paul Mbarga}

\textbf{Justifications grammaticales :} \esp{Me dirijo a usted} (verbe pronominal \esp{dirigirse a}, registre formel) ; \esp{quisiera} (forme de politesse, plus élégante que \esp{quiero}) ; \esp{me informara} (imparfait du subjonctif, obligatoire après \esp{le agradecería que}) ; \esp{hablo} (présent pour un fait actuel) ; \esp{sobre} = « au sujet de » (et non \esp{de} seul après \esp{informar}).

\textbf{Conclusion :} la lettre compte environ 70 mots, respecte la structure formelle et cumule les marqueurs de courtoisie attendus au Bac.
\end{exoresolu}

\courssub{Savoir-faire 2 : Utiliser le conditionnel de politesse}

Passons maintenant à l'outil grammatical central de la lettre formelle : le conditionnel de courtoisie. Sans lui, une demande sonne comme un ordre.

\begin{propriete}[Le conditionnel de courtoisie : forme et emploi]
\textbf{Formation du conditionnel} : infinitif + terminaisons \esp{-ía, -ías, -ía, -íamos, -íais, -ían} (identiques pour les trois groupes). Exemples : \esp{sería, podría, gustaría, agradecería}. \textbf{Toutes les formes portent un accent sur le í.}

\textbf{Radicaux irréguliers} (les mêmes qu'au futur simple) : \esp{querer $\rightarrow$ querr-} (\esp{querría}), \esp{poder $\rightarrow$ podr-}, \esp{saber $\rightarrow$ sabr-}, \esp{hacer $\rightarrow$ har-}, \esp{decir $\rightarrow$ dir-}, \esp{tener $\rightarrow$ tendr-}, \esp{salir $\rightarrow$ saldr-}, \esp{poner $\rightarrow$ pondr-}.

\textbf{Emploi} : le conditionnel adoucit la demande. \esp{Quisiera} (forme héritée de l'imparfait du subjonctif de \esp{querer}, figée comme conditionnel de politesse) vaut \esp{querría} dans ce registre : \esp{Quisiera un formulario} = « Je voudrais un formulaire ».

\textbf{Accord des temps} : après \esp{le agradecería que} ou \esp{quisiera que}, la subordonnée est à l'\textbf{imparfait du subjonctif} : \esp{Le agradecería que me enviara el dossier.} — jamais de futur, jamais de conditionnel dans la subordonnée.
\end{propriete}

\begin{methode}[Employer le conditionnel de courtoisie]
\begin{enumerate}
\item \textbf{Identifier le verbe de volonté} : \esp{querer, poder, gustar, agradecer, importar}.
\item \textbf{Le mettre au conditionnel} : \esp{quisiera} (je voudrais), \esp{podría usted...} (pourriez-vous), \esp{me gustaría} (j'aimerais), \esp{le agradecería} (je vous serais reconnaissant).
\item \textbf{Accorder le temps qui suit} : après \esp{le agradecería que} ou \esp{quisiera que}, le subordonné est à l'imparfait du subjonctif : \esp{Le agradecería que me enviara el dossier.}
\item \textbf{Éviter le présent direct} : \esp{Quiero información} est trop sec ; préférer \esp{Quisiera recibir información}.
\item \textbf{Conclure la lettre} par une formule figée : \esp{En espera de su respuesta...}
\end{enumerate}
Structure à retenir : \esp{Le agradecería que + imparfait du subjonctif} = « Je vous serais reconnaissant de bien vouloir... ».
\end{methode}

\begin{exoresolu}[Transformer des demandes trop directes]
Énoncé : Réécris ces phrases au registre formel avec le conditionnel de courtoisie.
\begin{enumerate}
\item \esp{Quiero un formulario de inscripción.}
\item \esp{Envíeme el horario de los cursos.}
\item \esp{Deme una respuesta rápida.}
\end{enumerate}
\tcblower
\textbf{Raisonnement.} Chaque phrase contient un verbe trop direct (\esp{quiero}, impératifs \esp{envíeme}, \esp{deme}). On remplace par un conditionnel de politesse ; quand on introduit \esp{que}, on déclenche l'imparfait du subjonctif.

\textbf{Réponses rédigées :}
\begin{enumerate}
\item \esp{Quisiera recibir un formulario de inscripción.} — \esp{quisiera} remplace \esp{quiero} ; l'infinitif \esp{recibir} précise la demande.
\item \esp{Le agradecería que me enviara el horario de los cursos.} — structure \esp{agradecería que} + imparfait du subjonctif \esp{enviara} (verbe \esp{enviar}, radical régulier au subjonctif imparfait : \esp{enviara, enviaras, enviara, enviáramos, enviarais, enviaran}).
\item \esp{¿Podría usted darme una respuesta lo antes posible?} — \esp{podría} (conditionnel de \esp{poder}, radical irrégulier \esp{podr-}) + infinitif \esp{dar} ; \esp{lo antes posible} = « dès que possible », formule idiomatique. Noter les signes \esp{¿...?} encadrant la question.
\end{enumerate}
\textbf{Conclusion :} le conditionnel de courtoisie est un marqueur de registre formel que le correcteur du Bac attend explicitement dans toute lettre administrative.
\end{exoresolu}

\courssub{Savoir-faire 3 : Traduire les formules de correspondance}

La traduction (thème et version) des formules de lettre est un exercice d'équivalence, non de mot à mot. Voici la démarche et le répertoire à mémoriser.

\begin{methode}[Traduire les formules de correspondance (version et thème)]
Les formules de lettre sont des blocs figés : on ne les traduit pas mot à mot, on les remplace par leur équivalent.
\begin{enumerate}
\item \textbf{Repérer la fonction} de la formule : appel, demande, remerciement, salutation finale.
\item \textbf{Choisir l'équivalent figé} dans l'autre langue (voir tableau ci-dessous).
\item \textbf{Vérifier le temps} : « je vous serais reconnaissant de » exige \esp{le agradecería que + subj. imparfait}, jamais un futur.
\item \textbf{Respecter la ponctuation espagnole} : deux points après l'appel, \esp{¿...?} pour une question.
\end{enumerate}
\begin{center}
\begin{tabularx}{\linewidth}{|X|X|X|}
\hline
\rowcolor{popblueL} Français & Español & Remarque \\
\hline
Monsieur le Directeur, & Estimado señor Director: & deux points obligatoires \\
\hline
Je me permets de vous écrire pour... & Me dirijo a usted para... & verbe pronominal \\
\hline
Je vous serais reconnaissant de bien vouloir... & Le agradecería que + subj. imparfait & pas de futur \\
\hline
Dans l'attente de votre réponse... & En espera de su respuesta... & formule figée \\
\hline
Veuillez agréer mes salutations distinguées. & Reciba un cordial saludo. & équivalent, pas calque \\
\hline
\end{tabularx}
\end{center}
\end{methode}

\begin{exoresolu}[Thème : formules de correspondance]
Énoncé : Traduis en espagnol.
\begin{enumerate}
\item Monsieur, je me permets de vous écrire pour demander des informations sur vos cours d'été.
\item Je vous serais reconnaissant de bien vouloir m'envoyer le dossier d'inscription.
\item Dans l'attente de votre réponse, je vous prie d'agréer mes salutations distinguées.
\end{enumerate}
\tcblower
\textbf{Raisonnement.} Ce sont des formules figées : on mobilise le tableau d'équivalents et on surveille les temps (subjonctif imparfait après \esp{agradecería que}).

\textbf{Traductions :}
\begin{enumerate}
\item \esp{Estimado señor: me dirijo a usted para solicitar información sobre sus cursos de verano.} — \esp{solicitar} = « demander » (faux ami : \esp{demandar} = « poursuivre en justice ») ; \esp{sobre} = « au sujet de » ; \esp{información} est singulier en espagnol là où le français dit « des informations ».
\item \esp{Le agradecería que me enviara el dossier de inscripción.} — \esp{enviara} : imparfait du subjonctif ; \esp{inscripción} avec accent sur la \emph{o}.
\item \esp{En espera de su respuesta, le ruego que acepte mis saludos más cordiales.} — \esp{le ruego que} + subjonctif présent \esp{acepte} (verbe \esp{aceptar}, régulier) ; on évite le calque \esp{agradar} pour « agréer ».
\end{enumerate}
\textbf{Conclusion :} en thème de correspondance, le correcteur évalue la fidélité au registre : équivalents figés, subjonctif correct, accents exacts.
\end{exoresolu}

\courssub{Savoir-faire 4 : Rédiger une lettre de réclamation}

La réclamation est une variante de la lettre de demande : même architecture, mais on y expose un grief. La difficulté est de rester ferme sans devenir agressif.

\begin{methode}[Rédiger une lettre de réclamation]
La réclamation reste formelle : on proteste avec fermeté mais sans agressivité.
\begin{enumerate}
\item \textbf{En-tête et appel} identiques à la lettre de demande.
\item \textbf{Rappeler les faits} avec précision (date, lieu, référence) : \esp{El pasado 3 de abril compré en su tienda un teléfono móvil.}
\item \textbf{Exposer le problème} objectivement : \esp{Desde el primer día, el aparato no funciona correctamente.}
\item \textbf{Formuler la demande} au conditionnel : \esp{Le agradecería que me cambiara el producto o me devolviera el dinero.}
\item \textbf{Annoncer une suite éventuelle} avec mesure : \esp{De lo contrario, me veré obligado a dirigirme a la asociación de consumidores.}
\item \textbf{Salutation finale} sobre : \esp{Atentamente,} + signature.
\end{enumerate}
Lexique utile : \esp{el producto defectuoso, la factura, la garantía, el reembolso, quejarse de, reclamar}.
\end{methode}

\begin{exoresolu}[Version : une lettre de réclamation]
Énoncé : Traduis en français le passage suivant, extrait d'une lettre de réclamation.

\esp{Estimado señor: El pasado lunes compré en su tienda del barrio de Briqueterie una bicicleta para mi hijo. Desgraciadamente, desde el primer día, los frenos no funcionan bien y la cadena se sale constantemente. Le agradecería que me cambiara la bicicleta o me devolviera el dinero lo antes posible. En espera de su respuesta, reciba un cordial saludo.}
\tcblower
\textbf{Raisonnement.} On repère les pièges : \esp{el pasado lunes} = « lundi dernier » (et non « le lundi passé » en français soigné) ; \esp{desgraciadamente} = « malheureusement » ; \esp{los frenos} = « les freins » ; \esp{se sale} = « déraille » (littéralement « sort ») ; \esp{devolviera} se rapporte ici à l'argent : « remboursât ».

\textbf{Traduction :}

« Monsieur, lundi dernier, j'ai acheté dans votre magasin du quartier de la Briqueterie un vélo pour mon fils. Malheureusement, dès le premier jour, les freins ne fonctionnent pas bien et la chaîne déraille constamment. Je vous serais reconnaissant de bien vouloir m'échanger le vélo ou me rembourser l'argent dès que possible. Dans l'attente de votre réponse, je vous adresse mes cordiales salutations. »

\textbf{Justifications :} \esp{compré} (prétérit simple) se rend par un passé composé ; \esp{lo antes posible} = « dès que possible » ; la formule finale est rendue par un équivalent français naturel, non par un calque.

\textbf{Conclusion :} en version, la précision du vocabulaire commercial (\esp{frenos, cadena, devolver el dinero}) et le naturel du français sont les deux critères de réussite.
\end{exoresolu}

\courssub{Savoir-faire 5 : Rédiger une lettre de motivation simple (100 mots)}

Dernière application, la plus fréquente à l'examen : la lettre de motivation d'environ 100 mots. La contrainte de longueur impose un plan serré.

\begin{methode}[Rédiger une lettre de motivation de 100 mots]
En 100 mots, chaque phrase doit compter. Plan en trois paragraphes :
\begin{enumerate}
\item \textbf{Paragraphe 1 (20 mots)} : qui je suis et ce que je demande. \esp{Actualmente soy alumno de Primera A y quisiera presentar mi candidatura para el puesto de...}
\item \textbf{Paragraphe 2 (50 mots)} : mes qualités et mon expérience, avec des connecteurs (\esp{además, también, por eso}). \esp{Hablo español y francés con fluidez; además, soy puntual y me gusta trabajar en equipo.}
\item \textbf{Paragraphe 3 (30 mots)} : disponibilité et politesse finale. \esp{Estoy disponible desde el mes de julio. Le agradecería que me concediera una entrevista.}
\end{enumerate}
Réflexes : un seul registre (formel), des adjectifs de qualité (\esp{serio, responsable, motivado, puntual}), aucune faute d'accord (\esp{estoy disponible}, \esp{soy alumno}).
\end{methode}

\begin{exoresolu}[Rédaction : lettre de motivation]
Énoncé (type Bac) : \esp{Has visto un anuncio: la librería « El Lector » de Yaoundé busca un vendedor para las vacaciones. Escribe una carta de motivación de unas 100 palabras.}
\tcblower
\textbf{Raisonnement.} On suit le plan en trois paragraphes ; on intègre au moins un conditionnel de courtoisie et deux connecteurs ; on compte les mots du corps de la lettre.

\textbf{Réponse rédigée :}

\esp{Estimado señor:}

\esp{Soy alumna de Primera A del Liceo Bilingüe de Yaoundé y quisiera presentar mi candidatura para el puesto de vendedora durante las vacaciones.}

\esp{Me encantan los libros y leo mucho, sobre todo novelas hispanas. Además, hablo español y francés con fluidez, lo que me permite atender a todos los clientes. Soy seria, puntual y me gusta trabajar en equipo. El verano pasado ya trabajé en la tienda de mi tío, en el mercado de Mokolo.}

\esp{Estoy disponible desde el 1 de julio. Le agradecería que me concediera una entrevista. En espera de su respuesta, reciba un cordial saludo.}

\esp{Amina Ngo Bell}

\textbf{Justifications :} \esp{quisiera} et \esp{concediera} forment la paire conditionnel + subjonctif imparfait ; \esp{lo que me permite} : relatif neutre correct ; \esp{trabajé} : prétérit simple pour une expérience terminée ; accords féminins respectés (\esp{alumna, vendedora, seria, disponible}) ; \esp{sobre todo} = « surtout » (deux mots).

\textbf{Conclusion :} le corps compte environ 100 mots, le plan est visible, les marqueurs de courtoisie et les connecteurs assurent la note maximale en expression écrite.
\end{exoresolu}

\begin{attention}[Les 5 erreurs qui coûtent des points au Bac]
\begin{enumerate}
\item \textbf{Oublier l'accent sur les formes verbales} : \esp{enviaria} au lieu de \esp{enviaría}, \esp{inscripcion} au lieu de \esp{inscripción}. Tout conditionnel espagnol porte un accent sur le \emph{í} : \esp{sería, podría, gustaría}.
\item \textbf{Le faux ami \esp{demandar}} : il signifie « poursuivre en justice ». Pour « demander », utilise \esp{solicitar, pedir} ou \esp{me dirijo a usted para...}.
\item \textbf{Le calque du futur français} : « je vous serais reconnaissant si vous m'enverrez » ne se traduit jamais par un futur. Après \esp{le agradecería que}, c'est l'imparfait du subjonctif : \esp{me enviara}, jamais \esp{me enviará}.
\item \textbf{La virgule après l'appel} : en espagnol, on écrit \esp{Estimado señor:} avec deux points, puis une majuscule à la ligne suivante. La virgule est un calque du français.
\item \textbf{Le mélange des registres} : passer de \esp{usted} à \esp{tú} dans la même lettre (\esp{Le agradecería que me enviaras tu folleto}) est éliminatoire. Choisis \esp{usted} et ses formes verbales (\esp{enviara, su, le}) du début à la fin.
\end{enumerate}
\end{attention}

\begin{savaistu}[La lettre formelle dans le monde hispanophone]
Dans les pays hispanophones, la lettre administrative garde un style très protocolaire : on y trouve encore souvent \esp{Muy señor mío:} ou \esp{A quien corresponda:} (« À qui de droit »), et la formule finale \esp{Atentamente} abrégée en \esp{Atte.} dans les courriers professionnels. En Guinée équatoriale, seul pays hispanophone d'Afrique subsaharienne, ce registre est celui de toute la correspondance officielle : un élève camerounais hispanophone y dispose donc d'un atout précieux pour les échanges avec Malabo ou Bata. À noter aussi : la date s'écrit toujours \esp{a 15 de mayo de 2024}, avec la préposition \esp{a} et le mois en minuscules.
\end{savaistu}

\begin{exercice}[Pour s'entraîner]
\begin{enumerate}
\item Réécris au registre formel : \esp{Quiero trabajar en su empresa. Dame el formulario.}
\item Traduis en espagnol : « Monsieur, je vous serais reconnaissant de bien vouloir m'envoyer la liste des documents nécessaires. Dans l'attente de votre réponse, recevez mes salutations distinguées. »
\item Rédige une lettre de réclamation (80 mots) : tu as acheté un téléphone défectueux dans un magasin de Douala ; tu demandes l'échange ou le remboursement.
\end{enumerate}
\end{exercice}

\newpage

\coursec{Exercices}

\coursec{Leçon EA20 : La correspondencia y los intercambios escritos (continuación) : cartas formales y administrativas}

\courssub{Texte support : une lettre de demande de stage}

\begin{textebox}[Carta de solicitud de prácticas]
Yaoundé, 15 de mayo de 2025

Estimado señor director :

Me dirijo a usted para solicitar información sobre la posibilidad de realizar unas prácticas profesionales en su empresa durante el próximo verano. Soy estudiante de Primera A en el Liceo de Nkolbisson y deseo adquirir experiencia en el sector de la comunicación digital.

Quisiera conocer las fechas disponibles, la duración de las prácticas y los requisitos necesarios. Le agradecería que me enviara el dossier de inscripción y cualquier documento útil.

En espera de su respuesta, le saluda atentamente,

Manga Bébé
\end{textebox}

\textbf{Glossaire :} \esp{solicitar} = demander ; \esp{prácticas profesionales} = stage professionnel ; \esp{sector de la comunicación digital} = secteur de la communication numérique ; \esp{dossier de inscripción} = dossier d'inscription ; \esp{En espera de su respuesta} = Dans l'attente de votre réponse.

\courssub{Lexique thématique de la correspondance formelle}

\begin{definition}[Vocabulaire clé de la lettre formelle]
\begin{itemize}
\item \esp{remitente} : expéditeur (celui qui écrit)
\item \esp{destinatario} : destinataire
\item \esp{sobre} : enveloppe
\item \esp{sello} : timbre
\item \esp{solicitud} : demande (lettre de demande)
\item \esp{reclamación} : réclamation
\item \esp{matrícula} : inscription / frais d'inscription
\item \esp{Estimado señor / Estimada señora} : Monsieur / Madame (formule d'appel)
\item \esp{Me dirijo a usted para...} : Je m'adresse à vous pour...
\item \esp{Quisiera / Le agradecería que...} : Je voudrais / Je vous serais reconnaissant de...
\item \esp{En espera de su respuesta} : Dans l'attente de votre réponse
\item \esp{Le saluda atentamente / Atentamente} : Je vous salue respectueusement / Cordialement
\end{itemize}
\end{definition}

\courssub{Grammaire : le conditionnel de courtoisie}

\begin{propriete}[Le conditionnel de courtoisie (cortesía)]
En espagnol, la politesse formelle s'exprime souvent par des formes verbales atténuées :
\begin{itemize}
\item \textbf{Quisiera} (imparfait du subjonctif de \esp{querer}) : « Je voudrais » (plus poli que \esp{querría}).
\item \textbf{Le agradecería} (conditionnel de \esp{agradecer}) : « Je vous serais reconnaissant de... » ; suivi du subjonctif (imparfait ou présent) : \esp{Le agradecería que me enviara / envíe el dossier}.
\item \textbf{Podría} (conditionnel de \esp{poder}) : « Je pourrais / Pourrais-je... » ; \esp{¿Podría decirme...?}
\item \textbf{Gustaría} (conditionnel de \esp{gustar}) : « J'aimerais... » ; \esp{Me gustaría hacer unas prácticas}.
\end{itemize}
\emph{Règle :} Après \esp{quisiera que}, \esp{le agradecería que}, \esp{les ruego que}, le verbe qui suit est au \textbf{subjonctif} (généralement imparfait : \esp{-ra/-se} ou présent).
\end{propriete}

\begin{attention}[Pièges fréquents pour francophones]
\begin{itemize}
\item Ne pas confondre \esp{quisiera} (subjonctif imparfait) et \esp{querría} (conditionnel) : \esp{quisiera} est la forme de politesse standard.
\item Après \esp{le agradecería que}, n'utilisez \textbf{jamais} l'indicatif (\esp{*responde}) mais le subjonctif (\esp{respondiera / responda}).
\item \esp{Estimado señor director :} (avec deux points) et non \esp{Estimado señor director,} (virgule).
\item La formule de clôture \esp{Le saluda atentamente} s'emploie même si l'expéditeur est « je » ; c'une formule figée.
\end{itemize}
\end{attention}

\courssub{Je vérifie mes connaissances}

\begin{exercice}[QCM : la correspondance formelle]
Entoure la bonne réponse.
\begin{enumerate}
\item Pour commencer une lettre formelle adressée à un directeur, on écrit :
\begin{enumerate}
\item \esp{Hola, director}
\item \esp{Estimado señor director :}
\item \esp{Querido amigo :}
\end{enumerate}
\item \esp{Quisiera} est :
\begin{enumerate}
\item un présent de l'indicatif
\item une forme de courtoisie (imparfait du subjonctif utilisé comme conditionnel de politesse)
\item un impératif
\end{enumerate}
\item Après \esp{Le agradecería que...}, on utilise :
\begin{enumerate}
\item l'imparfait du subjonctif (ou le présent du subjonctif)
\item le présent de l'indicatif
\item l'infinitif
\end{enumerate}
\item La formule de clôture correcte est :
\begin{enumerate}
\item \esp{En espera de su respuesta, le saluda atentamente}
\item \esp{En attendant votre réponse, salutations}
\item \esp{Adiós, hasta pronto}
\end{enumerate}
\end{enumerate}
\end{exercice}

\begin{exercice}[Vrai ou faux ? Justifie ta réponse]
Dis si les affirmations suivantes sont vraies ou fausses, puis justifie en une phrase.
\begin{enumerate}
\item Dans une lettre formelle, on tutoie le destinataire.
\item \esp{Me dirijo a usted} signifie « Je m'adresse à vous ».
\item On peut écrire \esp{Le agradecería que me responde pronto}.
\item Une lettre de réclamation doit rester polie même si l'on est mécontent.
\end{enumerate}
\end{exercice}

\begin{exercice}[Vocabulaire de la correspondance]
Complète chaque phrase avec le mot espagnol qui convient, choisi dans la liste : \esp{sello, sobre, remitente, destinatario, solicitud, matrícula}.
\begin{enumerate}
\item El \_\_\_\_\_\_\_\_\_\_\_\_ es la persona que escribe y envía la carta.
\item El \_\_\_\_\_\_\_\_\_\_\_\_ es la persona que recibe la carta.
\item Pongo la carta en un \_\_\_\_\_\_\_\_\_\_\_\_ y pego un \_\_\_\_\_\_\_\_\_\_\_\_.
\item Escribo una carta de \_\_\_\_\_\_\_\_\_\_\_\_ para pedir unas prácticas en una empresa.
\item Para estudiar en la universidad, hay que pagar la \_\_\_\_\_\_\_\_\_\_\_\_.
\end{enumerate}
\end{exercice}

\begin{exercice}[Conjugaison : formes de courtoisie (conditionnel et subjonctif imparfait)]
Complète avec la forme verbale appropriée (conditionnel de courtoisie ou subjonctif imparfait selon la structure).
\begin{enumerate}
\item \esp{Quisiera que usted me (enviar) \_\_\_\_\_\_\_\_\_\_\_\_ información sobre los cursos.}
\item \esp{Le agradecería que (responder) \_\_\_\_\_\_\_\_\_\_\_\_ pronto.}
\item \esp{(Poder, yo) \_\_\_\_\_\_\_\_\_\_\_\_ usted decirme el precio de la matrícula ?}
\item \esp{(Gustar, me) \_\_\_\_\_\_\_\_\_\_\_\_ hacer unas prácticas en su empresa.}
\end{enumerate}
\end{exercice}

\courssub{Je m'entraîne}

\begin{methode}[Traduire une formule de politesse administrative]
\begin{enumerate}
\item Identifier la fonction : appel, exposition du motif, demande, remerciement anticipé, clôture.
\item Choisir l'équivalent espagnol figé (ne pas traduire mot à mot).
\item Vérifier l'accord (genre/nombre) et la ponctuation (deux-points après l'appel, pas de virgule).
\item Utiliser \esp{usted} / \esp{ustedes} et les formes de courtoisie (\esp{quisiera}, \esp{le agradecería}, \esp{podría}, \esp{gustaría}).
\end{enumerate}
\end{methode}

\begin{exercice}[Traduction de formules de correspondance]
Traduis en espagnol les formules suivantes.
\begin{enumerate}
\item Je vous serais reconnaissant de bien vouloir m'envoyer le dossier d'inscription avant le 30 juin.
\item Dans l'attente de votre réponse, veuillez agréer, Monsieur, mes salutations distinguées.
\item J'ai l'honneur de vous demander des informations sur vos cours d'été.
\item Je vous prie de bien vouloir accepter mes excuses pour ce retard.
\end{enumerate}
\end{exercice}

\begin{exercice}[Transformation : du familier au formel]
Réécris les phrases suivantes dans un registre formel, avec \esp{usted} et les formes de courtoisie.
\begin{enumerate}
\item \esp{Hola, quiero información sobre el curso.}
\item \esp{Envíame el dossier rápido, por favor.}
\item \esp{Gracias, escríbeme pronto.}
\item \esp{¿Me puedes decir el precio ?}
\end{enumerate}
\end{exercice}

\begin{exercice}[Texte à trous : une lettre de demande]
Complète la lettre suivante avec les mots de la liste : \esp{Estimado, dirijo, quisiera, agradecería, espera, atentamente}.

\begin{textebox}[Carta de solicitud]
Yaoundé, 15 de mayo

\_\_\_\_\_\_\_\_\_\_\_\_ señor director :

Me \_\_\_\_\_\_\_\_\_\_\_\_ a usted para pedir información sobre los cursos de español de su instituto. \_\_\_\_\_\_\_\_\_\_\_\_ saber las fechas y los precios.

Le \_\_\_\_\_\_\_\_\_\_\_\_ que me enviara el dossier de inscripción.

En \_\_\_\_\_\_\_\_\_\_\_\_ de su respuesta, le saluda \_\_\_\_\_\_\_\_\_\_\_\_ ,

Manga Bébé
\end{textebox}
\end{exercice}

\begin{exercice}[Appariement : formules et fonctions]
Relie chaque formule espagnole (1-6) à sa fonction (a-f).

\begin{center}
\begin{tabularx}{\linewidth}{|X|X|}
\hline
\rowcolor{popblueL} \textbf{Fórmula} & \textbf{Función} \\
\hline
1. \esp{Me dirijo a usted para...} & a. terminar la carta \\
\hline
2. \esp{Quisiera información sobre...} & b. empezar y presentar el tema \\
\hline
3. \esp{Le agradecería que me enviara...} & c. pedir algo con cortesía \\
\hline
4. \esp{En espera de su respuesta...} & d. saludar al final \\
\hline
5. \esp{Le saluda atentamente} & e. pedir información \\
\hline
6. \esp{Estimado señor director :} & f. saludar al principio \\
\hline
\end{tabularx}
\end{center}
\end{exercice}

\courssub{Je me prépare au Bac}

\begin{exercice}[Compréhension et langue : una carta de reclamación]
Lis le texte suivant, puis réponds aux questions.

\begin{textebox}[Una carta de reclamación]
Douala, 10 de junio

Estimados señores :

Me dirijo a ustedes para reclamar por un problema con un pedido. El día 25 de mayo compré en su tienda en línea un ordenador portátil. Lo recibí el 2 de junio, pero no funciona : la pantalla está rota y el ordenador no enciende.

Estoy muy decepcionado porque pagué 250 000 francos CFA. Quisiera que ustedes me cambiaran el ordenador o me devolvieran el dinero. Les ruego que me respondan antes del 30 de junio.

En espera de su respuesta, les saluda atentamente,

Etoundi Mbarga
\end{textebox}

\textbf{Questions de compréhension} (réponds en espagnol, avec des phrases complètes) :
\begin{enumerate}
\item ¿Qué compró Etoundi y cuándo ?
\item ¿Cuál es el problema con el ordenador ?
\item ¿Qué pide Etoundi a la empresa ?
\end{enumerate}

\textbf{Questions de langue} :
\begin{enumerate}
\item Relève deux formules de politesse du texte et traduis-les en français.
\item Conjugue au conditionnel : \esp{querer} (yo) et \esp{poder} (usted).
\item Pourquoi utilise-t-on \esp{cambiaran} et \esp{devolvieran} après \esp{quisiera que} ?
\end{enumerate}
\end{exercice}

\begin{exercice}[Thème et version]
\textbf{Thème.} Traduis en espagnol :
\begin{enumerate}
\item Je m'adresse à vous pour demander un stage dans votre entreprise pendant les vacances.
\item Je vous serais reconnaissant de bien vouloir m'envoyer le dossier d'inscription avant le 30 juin.
\item Dans l'attente de votre réponse, veuillez agréer, Monsieur, mes salutations distinguées.
\end{enumerate}

\textbf{Version.} Traduis en français l'extrait suivant :

\begin{textebox}[Extrait d'une lettre administrative]
Me dirijo a usted para informarle de que su solicitud ha sido aceptada. Les ruego que envíen los documentos antes del 15 de julio. En espera de su respuesta, reciba un cordial saludo.
\end{textebox}
\end{exercice}

\begin{exercice}[Expression écrite : lettre de réclamation ou de demande]
Choisis \textbf{un} des deux sujets suivants et rédige ta lettre en espagnol.

\textbf{Sujet A.} Vous avez acheté un ordinateur en ligne qui ne fonctionne pas. Rédigez une lettre de réclamation à l'entreprise : présentez le problème, exprimez votre mécontentement avec politesse et formulez une demande précise (échange ou remboursement). (80-100 mots)

\textbf{Sujet B.} Vous souhaitez vous inscrire à un cours d'été à l'université de Salamanque. Rédigez une lettre de demande d'inscription : présentez-vous brièvement, demandez des informations (dates, prix, logement) et utilisez les formules de politesse appropriées. (environ 100 mots)

\textbf{Consignes :} respecte la structure d'une lettre formelle (lieu et date, salutation, corps, formule de clôture, signature) ; utilise au moins deux conditionnels de courtoisie (\esp{quisiera}, \esp{le agradecería que} + subjonctif).
\end{exercice}

\newpage

\coursec{Corrigés des exercices}

\begin{corrige}[Exercice 1 — QCM : la correspondance formelle]
\begin{enumerate}
\item \textbf{b.} \esp{Estimado señor director :} — c'est la formule d'appel formelle, avec deux points (et non une virgule). \esp{Hola} et \esp{Querido amigo} relèvent du registre familier.
\item \textbf{b.} \esp{Quisiera} est une forme de courtoisie : il s'agit de l'imparfait du subjonctif de \esp{querer}, employé comme conditionnel de politesse (« je voudrais »).
\item \textbf{a.} Après \esp{Le agradecería que...}, on utilise le subjonctif (imparfait : \esp{me enviara}, ou présent : \esp{me envíe}), jamais l'indicatif.
\item \textbf{a.} \esp{En espera de su respuesta, le saluda atentamente} est la clôture formelle correcte. La réponse b est en français, et la réponse c est familière.
\end{enumerate}
\end{corrige}

\begin{corrige}[Exercice 2 — Vrai ou faux ? Justifie ta réponse]
\begin{enumerate}
\item \textbf{Faux.} Dans une lettre formelle, on vouvoie le destinataire : on utilise \esp{usted} (et \esp{ustedes} au pluriel), avec les verbes à la troisième personne. Le tutoiement (\esp{tú}) est réservé aux lettres amicales.
\item \textbf{Vrai.} \esp{Me dirijo a usted} signifie « Je m'adresse à vous » ; c'est la formule d'ouverture classique pour exposer le motif de la lettre : \esp{Me dirijo a usted para solicitar...}
\item \textbf{Faux.} Après \esp{Le agradecería que}, le verbe doit être au subjonctif : \esp{Le agradecería que me respondiera pronto} (imparfait du subjonctif) ou \esp{que me responda pronto} (présent du subjonctif). \esp{*me responde} est un indicatif, donc une faute.
\item \textbf{Vrai.} Une lettre de réclamation exprime un mécontentement mais doit rester courtoise : on utilise \esp{quisiera}, \esp{les ruego que} + subjonctif, et la formule de clôture polie. L'agressivité dessert la demande.
\end{enumerate}
\end{corrige}

\begin{corrige}[Exercice 3 — Vocabulaire de la correspondance]
\begin{enumerate}
\item El \textbf{remitente} es la persona que escribe y envía la carta. (L'expéditeur est la personne qui écrit et envoie la lettre.)
\item El \textbf{destinatario} es la persona que recibe la carta. (Le destinataire est la personne qui reçoit la lettre.)
\item Pongo la carta en un \textbf{sobre} y pego un \textbf{sello}. (Je mets la lettre dans une enveloppe et je colle un timbre.)
\item Escribo una carta de \textbf{solicitud} para pedir unas prácticas en una empresa. (J'écris une lettre de demande pour solliciter un stage dans une entreprise.)
\item Para estudiar en la universidad, hay que pagar la \textbf{matrícula}. (Pour étudier à l'université, il faut payer les frais d'inscription.)
\end{enumerate}
\textbf{Point de vigilance :} \esp{solicitud} est un faux ami partiel : il signifie « demande, requête », et non « sollicitude ».
\end{corrige}

\begin{corrige}[Exercice 4 — Conjugaison : formes de courtoisie]
\begin{enumerate}
\item \esp{Quisiera que usted me \textbf{enviara} (ou \textbf{enviase}) información sobre los cursos.} — Après \esp{quisiera que}, subjonctif imparfait : \esp{enviar} $\rightarrow$ \esp{enviara/enviase}. « Je voudrais que vous m'envoyiez des informations sur les cours. »
\item \esp{Le agradecería que \textbf{respondiera} (ou \textbf{respondiese}) pronto.} — Après \esp{le agradecería que}, subjonctif imparfait. « Je vous serais reconnaissant de répondre rapidement. »
\item \esp{¿\textbf{Podría} usted decirme el precio de la matrícula?} — Conditionnel de \esp{poder} (radical irrégulier \esp{podr-}) : forme de courtoisie pour demander. « Pourriez-vous me dire le prix de l'inscription ? » Attention aux points d'interrogation doubles : ¿ ... ?
\item \esp{\textbf{Me gustaría} hacer unas prácticas en su empresa.} — Conditionnel de \esp{gustar} avec le pronom \esp{me} ; suivi de l'infinitif. « J'aimerais faire un stage dans votre entreprise. »
\end{enumerate}
\textbf{Rappel :} \esp{quisiera} + infinitif (même sujet) ; \esp{quisiera que} + subjonctif (sujets différents).
\end{corrige}

\begin{corrige}[Exercice 5 — Traduction de formules de correspondance]
\begin{enumerate}
\item « Je vous serais reconnaissant de bien vouloir m'envoyer le dossier d'inscription avant le 30 juin. » \\
\esp{Le agradecería que me enviara el expediente de inscripción antes del 30 de junio.} \\
(Variante acceptée : \esp{Le agradecería que me envíe el expediente de inscripción antes del 30 de junio.})
\item « Dans l'attente de votre réponse, veuillez agréer, Monsieur, mes salutations distinguées. » \\
\esp{En espera de su respuesta, le ruego que acepte, señor, mis saludos más atentos.} \\
(Variante plus courante : \esp{En espera de su respuesta, reciba, señor, un cordial saludo.})
\item « J'ai l'honneur de vous demander des informations sur vos cours d'été. » \\
\esp{Tengo el honor de solicitar información sobre sus cursos de verano.}
\item « Je vous prie de bien vouloir accepter mes excuses pour ce retard. » \\
\esp{Le ruego que acepte mis disculpas por este retraso.}
\end{enumerate}
\textbf{Points de vigilance :} ne jamais traduire « veuillez agréer » mot à mot (\esp{*quiera usted agradar} est incorrect) ; utiliser les formules figées. Après \esp{le ruego que}, subjonctif présent (\esp{acepte}, \esp{envíe}). Le gallécisme \esp{*dossier} est à éviter : on dit \esp{expediente} ou \esp{documentación}.
\end{corrige}

\begin{corrige}[Exercice 6 — Transformation : du familier au formel]
\begin{enumerate}
\item \esp{Hola, quiero información sobre el curso.} $\rightarrow$ \esp{Estimado señor : me dirijo a usted para solicitar información sobre el curso.} \\
(On passe de \esp{quiero}, trop direct, à \esp{me dirijo a usted para solicitar} ou \esp{quisiera información}.)
\item \esp{Envíame el dossier rápido, por favor.} $\rightarrow$ \esp{Le agradecería que me enviara el expediente lo antes posible.} \\
(L'impératif \esp{envíame} et l'adverbe mal formé \esp{rápido} deviennent \esp{le agradecería que} + subjonctif imparfait + \esp{lo antes posible}.)
\item \esp{Gracias, escríbeme pronto.} $\rightarrow$ \esp{En espera de su respuesta, le saluda atentamente.} \\
(Formule de clôture figée à la place de l'impératif familier.)
\item \esp{¿Me puedes decir el precio?} $\rightarrow$ \esp{¿Podría usted decirme el precio?} \\
(\esp{puedes} (tú) devient \esp{podría usted}, conditionnel de courtoisie.)
\end{enumerate}
\end{corrige}

\begin{corrige}[Exercice 7 — Texte à trous : une lettre de demande]
Lettre complétée :

\esp{\textbf{Estimado} señor director :}

\esp{Me \textbf{dirijo} a usted para pedir información sobre los cursos de español de su instituto. \textbf{Quisiera} saber las fechas y los precios.}

\esp{Le \textbf{agradecería} que me enviara el expediente de inscripción.}

\esp{En \textbf{espera} de su respuesta, le saluda \textbf{atentamente},}

\esp{Manga Bébé}

\textbf{Justifications :} \esp{Estimado} s'accorde en genre avec \esp{señor} (masculin singulier) ; \esp{dirijo} est la première personne du présent de \esp{dirigirse} (verbe à changement orthographique g $\rightarrow$ j devant o) ; \esp{quisiera} + infinitif \esp{saber} (même sujet) ; \esp{agradecería} est suivi de \esp{que} + subjonctif ; \esp{en espera de} est une locution figée ; \esp{atentamente} est un adverbe en \esp{-mente} (sans accent).
\end{corrige}

\begin{corrige}[Exercice 8 — Appariement : formules et fonctions]
\begin{center}
\begin{tabularx}{\linewidth}{|X|X|}
\hline
\rowcolor{popblueL} \textbf{Fórmula} & \textbf{Función} \\
\hline
1. \esp{Me dirijo a usted para...} & \textbf{b.} empezar y presentar el tema \\
\hline
2. \esp{Quisiera información sobre...} & \textbf{e.} pedir información \\
\hline
3. \esp{Le agradecería que me enviara...} & \textbf{c.} pedir algo con cortesía \\
\hline
4. \esp{En espera de su respuesta...} & \textbf{a.} terminar la carta \\
\hline
5. \esp{Le saluda atentamente} & \textbf{d.} saludar al final \\
\hline
6. \esp{Estimado señor director :} & \textbf{f.} saludar al principio \\
\hline
\end{tabularx}
\end{center}
\textbf{Astuce :} repérer la place de chaque formule dans la lettre : appel (début), exposition du motif, demande(s), remerciement anticipé, clôture (fin).
\end{corrige}

\begin{corrige}[Exercice 9 — Compréhension et langue : una carta de reclamación]
\textbf{Questions de compréhension} (réponses en espagnol) :
\begin{enumerate}
\item \esp{Etoundi compró un ordenador portátil el día 25 de mayo en una tienda en línea.} (Etoundi a acheté un ordinateur portable le 25 mai dans une boutique en ligne.)
\item \esp{El problema es que el ordenador no funciona : la pantalla está rota y el ordenador no enciende.} (Le problème est que l'ordinateur ne fonctionne pas : l'écran est cassé et l'ordinateur ne s'allume pas.)
\item \esp{Etoundi pide que la empresa le cambie el ordenador o le devuelva el dinero, y que le responda antes del 30 de junio.} (Il demande que l'entreprise échange l'ordinateur ou lui rembourse l'argent, et qu'elle réponde avant le 30 juin.)
\end{enumerate}

\textbf{Questions de langue} :
\begin{enumerate}
\item Deux formules de politesse (au choix) :
\begin{itemize}
\item \esp{Quisiera que ustedes me cambiaran el ordenador} : « Je voudrais que vous m'échangiez l'ordinateur » (conditionnel de courtoisie + subjonctif imparfait).
\item \esp{Les ruego que me respondan antes del 30 de junio} : « Je vous prie de me répondre avant le 30 juin ».
\item \esp{En espera de su respuesta, les saluda atentamente} : « Dans l'attente de votre réponse, je vous salue respectueusement ».
\end{itemize}
\item Conditionnel : \esp{querer} (yo) $\rightarrow$ \esp{querría} (radical irrégulier \esp{querr-}) ; \esp{poder} (usted) $\rightarrow$ \esp{podría} (radical irrégulier \esp{podr-}).
\item On utilise \esp{cambiaran} et \esp{devolvieran} parce qu'après \esp{quisiera que}, le verbe de la subordonnée se met au subjonctif ; comme la demande porte sur une action future vue depuis un énoncé atténué, l'imparfait du subjonctif (terminaison \esp{-ra}) est la forme de courtoisie attendue.
\end{enumerate}
\textbf{Barème indicatif (10 points) :} compréhension 3 $\times$ 1,5 pt = 4,5 pts ; questions de langue : 2 pts (formules), 1,5 pt (conjugaison), 2 pts (justification du subjonctif).
\end{corrige}

\begin{corrige}[Exercice 10 — Thème et version]
\textbf{Thème.}
\begin{enumerate}
\item « Je m'adresse à vous pour demander un stage dans votre entreprise pendant les vacances. » \\
\esp{Me dirijo a usted para solicitar unas prácticas en su empresa durante las vacaciones.}
\item « Je vous serais reconnaissant de bien vouloir m'envoyer le dossier d'inscription avant le 30 juin. » \\
\esp{Le agradecería que me enviara el expediente de inscripción antes del 30 de junio.}
\item « Dans l'attente de votre réponse, veuillez agréer, Monsieur, mes salutations distinguées. » \\
\esp{En espera de su respuesta, reciba, señor, un cordial saludo.} (ou : \esp{...le ruego que acepte, señor, mis saludos más atentos.})
\end{enumerate}
\textbf{Points de vigilance :} \esp{prácticas} (et non \esp{*un stage}) ; subjonctif imparfait \esp{enviara} après \esp{le agradecería que} ; formule de clôture figée, sans traduction mot à mot ; \esp{expediente} au lieu du gallécisme \esp{*dossier}.

\textbf{Version.} \\
« Je m'adresse à vous pour vous informer que votre demande a été acceptée. Je vous prie d'envoyer les documents avant le 15 juillet. Dans l'attente de votre réponse, recevez un cordial salut (mes salutations cordiales). »

\textbf{Points de vigilance :} \esp{su solicitud ha sido aceptada} = « votre demande a été acceptée » (passif avec \esp{ser} + participe accordé) ; \esp{les ruego que envíen} = « je vous prie d'envoyer » (subjonctif présent rendu par un infinitif français) ; \esp{reciba un cordial saludo} est une formule de clôture.

\textbf{Barème indicatif (10 points) :} thème 3 $\times$ 2 pts = 6 pts ; version 4 pts (exactitude du sens 2 pts, qualité du français 1 pt, formules 1 pt).
\end{corrige}

\begin{corrige}[Exercice 11 — Expression écrite : lettre de réclamation ou de demande]
\textbf{Proposition de rédaction modèle — Sujet A (lettre de réclamation) :}

\esp{Douala, 12 de junio de 2025}

\esp{Estimados señores :}

\esp{Me dirijo a ustedes para reclamar por un problema con un pedido. El 5 de junio compré en su tienda en línea un ordenador portátil, pero no funciona : la pantalla está rota y no enciende.}

\esp{Estoy muy decepcionado porque pagué 250 000 francos CFA. Quisiera que me cambiaran el ordenador o me devolvieran el dinero. Les agradecería que me respondieran antes del 30 de junio.}

\esp{En espera de su respuesta, les saluda atentamente,}

\esp{Etoundi Mbarga}

(85 mots environ.)

\textbf{Proposition de rédaction modèle — Sujet B (lettre de demande d'inscription) :}

\esp{Yaoundé, 15 de mayo de 2025}

\esp{Estimado señor director :}

\esp{Me dirijo a usted para solicitar información sobre los cursos de verano de la Universidad de Salamanca. Soy estudiante de Primera A en el Liceo de Nkolbisson y quisiera mejorar mi español en España.}

\esp{¿Podría usted decirme las fechas de los cursos, el precio de la matrícula y las posibilidades de alojamiento ? Le agradecería que me enviara el expediente de inscripción antes del 30 de junio.}

\esp{En espera de su respuesta, le saluda atentamente,}

\esp{Manga Bébé}

(90 mots environ.)

\textbf{Barème indicatif (10 points) :} respect de la structure formelle (lieu/date, appel, corps, clôture, signature) : 3 pts ; emploi d'au moins deux conditionnels de courtoisie corrects (\esp{quisiera}, \esp{le agradecería que} + subjonctif, \esp{podría}) : 3 pts ; correction grammaticale et orthographe (accents, ¿ ¡) : 2 pts ; richesse du lexique de la correspondance : 2 pts.

\textbf{Points de vigilance :}
\begin{itemize}
\item Deux points après l'appel : \esp{Estimado señor director :} ; jamais de virgule.
\item \esp{usted} partout : pas de \esp{tú} ni d'impératif familier.
\item Après \esp{quisiera que} et \esp{le agradecería que}, subjonctif obligatoire (\esp{enviara}, \esp{respondieran}).
\item Ne pas écrire \esp{*un stage} : dire \esp{unas prácticas} ou \esp{un curso}.
\item Vérifier les accents : \esp{información}, \esp{matrícula}, \esp{atentamente} (sans accent).
\item Éviter le gallécisme \esp{*dossier} : utiliser \esp{expediente} ou \esp{documentación}.
\end{itemize}
\end{corrige}

\newpage

\coursec{Fiche bilan}

\begin{popfigure}
\begin{tikzpicture}[mindmap, grow cyclic, every node/.style={concept, text=white, font=\small},
  root concept/.append style={concept color=popblue, font=\small\bfseries},
  level 1/.append style={level distance=3.2cm, sibling angle=60},
  level 1 concept/.append style={font=\footnotesize}]
\node[root concept]{Cartas formales}
  child[concept color=poporange]{ node[align=center]{Solicitud\\ (stage, inscripción)} }
  child[concept color=popgreen]{ node {Reclamación} }
  child[concept color=poppurple]{ node[align=center]{Motivación\\ (100 palabras)} }
  child[concept color=poppink]{ node[align=center]{Condicional\\ de cortesía} }
  child[concept color=popgold]{ node[align=center]{Fórmulas\\ fijas} }
  child[concept color=popdark]{ node[align=center]{Estructura\\ de la carta} };
\end{tikzpicture}
\legende{Carte mentale de la leçon : la correspondance formelle}
\end{popfigure}

\begin{bilan}
\textbf{1. Lexique clé de la correspondance formelle}
\begin{itemize}
\item \esp{Estimado señor / Estimada señora} : « Monsieur / Madame » (formule d'appel)
\item \esp{Me dirijo a usted para...} : « Je m'adresse à vous pour... »
\item \esp{quisiera / me gustaría} : « je voudrais / j'aimerais »
\item \esp{Le agradecería que + subjuntivo} : « Je vous serais reconnaissant de... »
\item \esp{adjuntar el currículum} : « joindre le CV »
\item \esp{una solicitud de prácticas} : « une demande de stage »
\item \esp{presentar una reclamación} : « déposer une réclamation »
\item \esp{En espera de su respuesta} : « Dans l'attente de votre réponse »
\item \esp{Atentamente / Cordialmente} : « Cordialement / Veuillez agréer... »
\item \esp{el remitente / el destinatario} : « l'expéditeur / le destinataire »
\end{itemize}

\textbf{2. Grammaire à connaître par cœur : le conditionnel de courtoisie}
\begin{center}
\begin{tabularx}{\linewidth}{|l|X|X|}
\hline
\rowcolor{popblueL} Verbe & Forme & Exemple traduit \\
\hline
querer & \esp{quisiera} & \esp{Quisiera información.} « Je voudrais des renseignements. » \\
\hline
poder & \esp{podría} & \esp{¿Podría enviarme el folleto?} « Pourriez-vous m'envoyer la brochure ? » \\
\hline
agradecer & \esp{agradecería} & \esp{Le agradecería que me contestara.} « Je vous serais reconnaissant de me répondre. » \\
\hline
gustar & \esp{me gustaría} & \esp{Me gustaría hacer prácticas.} « J'aimerais faire un stage. » \\
\hline
deber & \esp{debería} & \esp{Debería devolverme el dinero.} « Vous devriez me rembourser. » \\
\hline
\end{tabularx}
\end{center}
Règle : infinitif + \textbf{-ía, -ías, -ía, -íamos, -íais, -ían} (mêmes radicaux que le futur : \esp{dir-ía}, \esp{har-ía}, \esp{podr-ía}). Après \esp{Le agradecería que}, toujours le \textbf{subjonctif imparfait} : \esp{...que me enviara}.

\textbf{3. Structure express d'une lettre formelle}
\begin{enumerate}
\item En-tête : lieu et date \esp{Yaoundé, a 15 de mayo de 2024}
\item Appel : \esp{Estimado señor Director:} (deux-points, jamais de virgule)
\item Objet : \esp{Asunto: solicitud de prácticas}
\item Corps : présentation, demande (conditionnel), remerciement
\item Formule finale : \esp{Le saluda atentamente,} + signature
\end{enumerate}

\textbf{4. Méthodes express}
\begin{itemize}
\item Lettre de demande : qui je suis, ce que je veux, pourquoi, remerciement.
\item Lettre de réclamation : faits précis (dates, produit), ton ferme mais poli, solution exigée (\esp{Le ruego que...}).
\item Lettre de motivation (100 mots) : 3 paragraphes = présentation, qualités et expérience, disponibilité et remerciement.
\item Traduction : repérer d'abord la formule figée, ne jamais traduire mot à mot (\esp{Le agradecería que} ne se traduit pas par « je vous remercierais que »).
\end{itemize}
\end{bilan}

\begin{aretenir}[Checklist avant le Bac]
\begin{itemize}
\item[$\square$] Je sais écrire l'appel \esp{Estimado señor:} avec les deux-points.
\item[$\square$] Je connais par cœur \esp{quisiera}, \esp{me gustaría}, \esp{Le agradecería que}.
\item[$\square$] Je sais conjuguer le conditionnel : infinitif + \textbf{-ía}.
\item[$\square$] Je mets le subjonctif imparfait après \esp{Le agradecería que}.
\item[$\square$] Je distingue lettre de demande, lettre de réclamation et lettre de motivation.
\item[$\square$] Je sais écrire la date en espagnol : \esp{a 15 de mayo de 2024}.
\item[$\square$] Je connais deux formules de politesse finales : \esp{Atentamente}, \esp{Le saluda cordialmente}.
\item[$\square$] Je sais joindre un document : \esp{Le adjunto mi currículum}.
\item[$\square$] Je peux rédiger une lettre de motivation de 100 mots en 3 paragraphes.
\item[$\square$] Je traduis les formules figées sans calquer le français.
\item[$\square$] Je vérifie les accents : \esp{información}, \esp{atención}, \esp{señor}.
\item[$\square$] Je relis ma lettre : appel, objet, corps, formule finale, signature.
\end{itemize}
\end{aretenir}

\findecours

\end{document}
